\documentclass[11pt]{article}
\usepackage{amsmath,amssymb}
\usepackage[margin=1in]{geometry}
\usepackage{hyperref}
\allowdisplaybreaks
\title{The VNB Library --- Results by Topic}
\author{Generated by the VNB proof checker}
\date{\today}
\begin{document}
\maketitle
\noindent Each result is shown three ways: an English paraphrase, the formal statement, and --- for machine-checked theorems --- a one-line idea of the proof together with the facts it is checked modulo.  Grouped by subject, like a textbook table of contents.

\tableofcontents
\bigskip

\section{Calculus basics}

\subsection{Vocabulary}

\subsubsection*{\texttt{is-diff-at}}
\emph{for every f, a, l, f is differentiable at a, with derivative l if and only if f is in fun(rr, rr) and a is in rr and l is in rr and there is some phi in fun(rr, rr) such that phi is continuous at a and phi(a) equals l and for every x in rr, f(x) - f(a) equals phi(x) * (x - a)}
{\small\[
\forall f, a, l.\; (\operatorname{is\text{-}diff\text{-}at}(f, a, l) \Leftrightarrow ((f \in (\mathbb{R} \to \mathbb{R})) \wedge ((a \in \mathbb{R}) \wedge ((l \in \mathbb{R}) \wedge \exists \varphi \in (\mathbb{R} \to \mathbb{R}).\; (\operatorname{is\text{-}continuous\text{-}at}(\operatorname{rr\text{-}ms}, \operatorname{rr\text{-}ms}, \varphi, a) \wedge ((\varphi(a) = l) \wedge \forall x \in \mathbb{R}.\; ((f(x) - f(a)) = (\varphi(x) \cdot (x - a)))))))))
\]}
\noindent\textbf{Definition.}

\subsubsection*{\texttt{deriv}}
\emph{Definition / vocabulary.}

\subsubsection*{\texttt{nth-deriv}}
\emph{Definition / vocabulary.}

\subsubsection*{\texttt{little-o-at}}
\emph{for every g, a, g is little-o at a if and only if g is in fun(rr, rr) and a is in rr and there is some eps in fun(rr, rr) such that eps is continuous at a and eps(a) equals 0 and for every x in rr, g(x) equals eps(x) * (x - a)}
{\small\[
\forall g, a.\; (\operatorname{little\text{-}o\text{-}at}(g, a) \Leftrightarrow ((g \in (\mathbb{R} \to \mathbb{R})) \wedge ((a \in \mathbb{R}) \wedge \exists \mathit{eps} \in (\mathbb{R} \to \mathbb{R}).\; (\operatorname{is\text{-}continuous\text{-}at}(\operatorname{rr\text{-}ms}, \operatorname{rr\text{-}ms}, \mathit{eps}, a) \wedge ((\operatorname{eps}(a) = 0) \wedge \forall x \in \mathbb{R}.\; (g(x) = (\operatorname{eps}(x) \cdot (x - a))))))))
\]}
\noindent\textbf{Definition.}

\subsubsection*{\texttt{taylor-poly}}
\emph{Definition / vocabulary.}

\subsubsection*{\texttt{taylor-differentiable}}
\emph{for every f, a, x, n, f is n times differentiable from a to x if and only if for every k, if k is in nn and k is at most n, then for every t in ccint(a, x), nth-deriv(f, k) is continuous at t and for every k, if k is in nn and k is at most n, then for every t, if a is less than t and t is less than x, then nth-deriv(f, k) is differentiable at t, with derivative (nth-deriv(f, succ(k)))(t)}
{\small\[
\forall f, a, x, n.\; (\operatorname{taylor\text{-}differentiable}(f, a, x, n) \Leftrightarrow (\forall k.\; (((k \in \mathbb{N}) \wedge (k \leq n)) \Rightarrow \forall t \in \operatorname{ccint}(a, x).\; \operatorname{is\text{-}continuous\text{-}at}(\operatorname{rr\text{-}ms}, \operatorname{rr\text{-}ms}, \operatorname{nth\text{-}deriv}(f, k), t)) \wedge \forall k.\; (((k \in \mathbb{N}) \wedge (k \leq n)) \Rightarrow \forall t.\; (((a < t) \wedge (t < x)) \Rightarrow \operatorname{is\text{-}diff\text{-}at}(\operatorname{nth\text{-}deriv}(f, k), t, \operatorname{nth\text{-}deriv}(f, (k + 1))(t))))))
\]}
\noindent\textbf{Definition.}

\subsection{Differentiation rules}

\subsubsection*{\texttt{derivative-unique}}
\emph{for every f, a, l, m, if f is differentiable at a, with derivative l and f is differentiable at a, with derivative m, then l equals m}
{\small\[
\forall f, a, l, m.\; ((\operatorname{is\text{-}diff\text{-}at}(f, a, l) \wedge \operatorname{is\text{-}diff\text{-}at}(f, a, m)) \Rightarrow (l = m))
\]}
\noindent\textbf{Machine-checked} (trust: proof). Skolemize the two Caratheodory factors phi, psi from the two differentiability hypotheses; off the point they agree by cancelling (x-a) =/= 0, and continuity at a forces phi(a)=psi(a), i.e. L=M.

\subsubsection*{\texttt{diff-implies-continuous}}
\emph{for every f, a, l, if f is differentiable at a, with derivative l, then f is continuous at a}
{\small\[
\forall f, a, l.\; (\operatorname{is\text{-}diff\text{-}at}(f, a, l) \Rightarrow \operatorname{is\text{-}continuous\text{-}at}(\operatorname{rr\text{-}ms}, \operatorname{rr\text{-}ms}, f, a))
\]}
\noindent\textbf{Machine-checked} (trust: proof). f coincides pointwise with the map x |-> f(a) + phi(x)*(x-a), which the continuity algebra makes continuous at a (a constant plus phi times x-a); continuity then transfers to f.

\subsubsection*{\texttt{deriv-const}}
\emph{for every c, a, if c is in rr and a is in rr, then vnb-lambda(x, rr, c) is differentiable at a, with derivative 0}
{\small\[
\forall c, a.\; (((c \in \mathbb{R}) \wedge (a \in \mathbb{R})) \Rightarrow \operatorname{is\text{-}diff\text{-}at}((\lambda x \in \mathbb{R}.\; c), a, 0))
\]}
\noindent\textbf{Machine-checked} (trust: proof). Witness phi := 0: the difference f(x)-f(a) = 0 = 0*(x-a), and the zero map is continuous with value 0.

\subsubsection*{\texttt{deriv-identity}}
\emph{for every a in rr, vnb-lambda(x, rr, x) is differentiable at a, with derivative 1}
{\small\[
\forall a \in \mathbb{R}.\; \operatorname{is\text{-}diff\text{-}at}((\lambda x \in \mathbb{R}.\; x), a, 1)
\]}
\noindent\textbf{Machine-checked} (trust: proof). Witness phi := 1: x-a = 1*(x-a), and the constant map 1 is continuous.

\subsubsection*{\texttt{deriv-sum}}
\emph{for every f, g, a, l, m, if f is differentiable at a, with derivative l and g is differentiable at a, with derivative m, then vnb-lambda(x, rr, f(x) + g(x)) is differentiable at a, with derivative l + m}
{\small\[
\forall f, g, a, l, m.\; ((\operatorname{is\text{-}diff\text{-}at}(f, a, l) \wedge \operatorname{is\text{-}diff\text{-}at}(g, a, m)) \Rightarrow \operatorname{is\text{-}diff\text{-}at}((\lambda x \in \mathbb{R}.\; (f(x) + g(x))), a, (l + m)))
\]}
\noindent\textbf{Machine-checked} (trust: proof).

\subsubsection*{\texttt{deriv-product}}
\emph{for every f, g, a, l, m, if f is differentiable at a, with derivative l and g is differentiable at a, with derivative m, then vnb-lambda(x, rr, f(x) * g(x)) is differentiable at a, with derivative l * g(a) + f(a) * m}
{\small\[
\forall f, g, a, l, m.\; ((\operatorname{is\text{-}diff\text{-}at}(f, a, l) \wedge \operatorname{is\text{-}diff\text{-}at}(g, a, m)) \Rightarrow \operatorname{is\text{-}diff\text{-}at}((\lambda x \in \mathbb{R}.\; (f(x) \cdot g(x))), a, ((l \cdot g(a)) + (f(a) \cdot m))))
\]}
\noindent\textbf{Machine-checked} (trust: proof).

\subsubsection*{\texttt{deriv-chain}}
\emph{for every f, g, a, l, m, if f is differentiable at a, with derivative l, then if g is differentiable at f(a), with derivative m, then compose(g, f) is differentiable at a, with derivative m * l}
{\small\[
\forall f, g, a, l, m.\; (\operatorname{is\text{-}diff\text{-}at}(f, a, l) \Rightarrow (\operatorname{is\text{-}diff\text{-}at}(g, f(a), m) \Rightarrow \operatorname{is\text{-}diff\text{-}at}(\operatorname{compose}(g, f), a, (m \cdot l))))
\]}
\noindent\textbf{Machine-checked} (trust: proof).

\subsubsection*{\texttt{deriv-chain-value}}
\emph{for every f, g, a, dl, dm, if f is differentiable at a, with derivative dl, then if g is differentiable at f(a), with derivative dm, then deriv(compose(g, f), a) equals deriv(g, f(a)) * deriv(f, a)}
{\small\[
\forall f, g, a, \mathit{dl}, \mathit{dm}.\; (\operatorname{is\text{-}diff\text{-}at}(f, a, \mathit{dl}) \Rightarrow (\operatorname{is\text{-}diff\text{-}at}(g, f(a), \mathit{dm}) \Rightarrow (\operatorname{deriv}(\operatorname{compose}(g, f), a) = (\operatorname{deriv}(g, f(a)) \cdot \operatorname{deriv}(f, a)))))
\]}
\noindent\textbf{Machine-checked} (trust: proof).

\subsubsection*{\texttt{deriv-of-is-diff-at}}
\emph{for every f, a, dv, if f is differentiable at a, with derivative dv, then deriv(f, a) equals dv}
{\small\[
\forall f, a, \mathit{dv}.\; (\operatorname{is\text{-}diff\text{-}at}(f, a, \mathit{dv}) \Rightarrow (\operatorname{deriv}(f, a) = \mathit{dv}))
\]}
\noindent\textbf{Machine-checked} (trust: proof).

\subsubsection*{\texttt{deriv-neg}}
\emph{for every f, a, l, if f is differentiable at a, with derivative l, then vnb-lambda(z, rr, -f(z)) is differentiable at a, with derivative -l}
{\small\[
\forall f, a, l.\; (\operatorname{is\text{-}diff\text{-}at}(f, a, l) \Rightarrow \operatorname{is\text{-}diff\text{-}at}((\lambda z \in \mathbb{R}.\; -f(z)), a, -l))
\]}
\noindent\textbf{Machine-checked} (trust: proof).

\subsubsection*{\texttt{nth-deriv-one}}
\emph{for every f, nth-deriv(f, 1) is identical to vnb-lambda(x, rr, deriv(f, x))}
{\small\[
\forall f.\; (\operatorname{nth\text{-}deriv}(f, 1) \simeq (\lambda x \in \mathbb{R}.\; \operatorname{deriv}(f, x)))
\]}
\noindent\textbf{Machine-checked} (trust: proof).

\subsection{Little-o calculus}

\subsubsection*{\texttt{diff-iff-little-o}}
\emph{for every f, a, l, f is differentiable at a, with derivative l if and only if f is in fun(rr, rr) and a is in rr and l is in rr and vnb-lambda(x, rr, f(x) - f(a) - l * (x - a)) is little-o at a}
{\small\[
\forall f, a, l.\; (\operatorname{is\text{-}diff\text{-}at}(f, a, l) \Leftrightarrow ((f \in (\mathbb{R} \to \mathbb{R})) \wedge ((a \in \mathbb{R}) \wedge ((l \in \mathbb{R}) \wedge \operatorname{little\text{-}o\text{-}at}((\lambda x \in \mathbb{R}.\; ((f(x) - f(a)) - (l \cdot (x - a)))), a)))))
\]}
\noindent\textbf{Machine-checked} (trust: proof).

\subsubsection*{\texttt{little-o-sum}}
\emph{for every g, h, a, if g is little-o at a, then if h is little-o at a, then vnb-lambda(x, rr, g(x) + h(x)) is little-o at a}
{\small\[
\forall g, h, a.\; (\operatorname{little\text{-}o\text{-}at}(g, a) \Rightarrow (\operatorname{little\text{-}o\text{-}at}(h, a) \Rightarrow \operatorname{little\text{-}o\text{-}at}((\lambda x \in \mathbb{R}.\; (g(x) + h(x))), a)))
\]}
\noindent\textbf{Machine-checked} (trust: proof).

\subsubsection*{\texttt{little-o-scalar}}
\emph{for every c, g, a, if c is in rr, then if g is little-o at a, then vnb-lambda(x, rr, c * g(x)) is little-o at a}
{\small\[
\forall c \in \mathbb{R}, g, a.\; (\operatorname{little\text{-}o\text{-}at}(g, a) \Rightarrow \operatorname{little\text{-}o\text{-}at}((\lambda x \in \mathbb{R}.\; (c \cdot g(x))), a))
\]}
\noindent\textbf{Machine-checked} (trust: proof).

\subsection{Extreme \& interior values (Fermat)}

\subsubsection*{\texttt{extreme-value-max}}
\emph{for every f, a, b, if f is in fun(rr, rr) and a is in rr and b is in rr and a is at most b, then if for every x in ccint(a, b), f is continuous at x, then there is some c in ccint(a, b) such that for every x in ccint(a, b), f(x) is at most f(c)}
{\small\[
\forall f, a, b.\; (((f \in (\mathbb{R} \to \mathbb{R})) \wedge ((a \in \mathbb{R}) \wedge ((b \in \mathbb{R}) \wedge (a \leq b)))) \Rightarrow (\forall x \in \operatorname{ccint}(a, b).\; \operatorname{is\text{-}continuous\text{-}at}(\operatorname{rr\text{-}ms}, \operatorname{rr\text{-}ms}, f, x) \Rightarrow \exists c \in \operatorname{ccint}(a, b).\; \forall x \in \operatorname{ccint}(a, b).\; (f(x) \leq f(c))))
\]}
\noindent\textbf{Machine-checked} (trust: proof).

\subsubsection*{\texttt{extreme-value-min}}
\emph{for every f, a, b, if f is in fun(rr, rr) and a is in rr and b is in rr and a is at most b, then if for every x in ccint(a, b), f is continuous at x, then there is some c in ccint(a, b) such that for every x in ccint(a, b), f(c) is at most f(x)}
{\small\[
\forall f, a, b.\; (((f \in (\mathbb{R} \to \mathbb{R})) \wedge ((a \in \mathbb{R}) \wedge ((b \in \mathbb{R}) \wedge (a \leq b)))) \Rightarrow (\forall x \in \operatorname{ccint}(a, b).\; \operatorname{is\text{-}continuous\text{-}at}(\operatorname{rr\text{-}ms}, \operatorname{rr\text{-}ms}, f, x) \Rightarrow \exists c \in \operatorname{ccint}(a, b).\; \forall x \in \operatorname{ccint}(a, b).\; (f(c) \leq f(x))))
\]}
\noindent\textbf{Machine-checked} (trust: proof).

\subsubsection*{\texttt{interior-max-deriv-zero}}
\emph{for every f, a, b, theta, l, if f is in fun(rr, rr) and a is in rr and b is in rr and theta is in rr and a is less than theta and theta is less than b, then if for every x in ccint(a, b), f(x) is at most f(theta), then if f is differentiable at theta, with derivative l, then l equals 0}
{\small\[
\forall f, a, b, \theta, l.\; (((f \in (\mathbb{R} \to \mathbb{R})) \wedge ((a \in \mathbb{R}) \wedge ((b \in \mathbb{R}) \wedge ((\theta \in \mathbb{R}) \wedge ((a < \theta) \wedge (\theta < b)))))) \Rightarrow (\forall x \in \operatorname{ccint}(a, b).\; (f(x) \leq f(\theta)) \Rightarrow (\operatorname{is\text{-}diff\text{-}at}(f, \theta, l) \Rightarrow (l = 0))))
\]}
\noindent\textbf{Machine-checked} (trust: proof).

\subsubsection*{\texttt{interior-min-deriv-zero}}
\emph{for every f, a, b, theta, l, if f is in fun(rr, rr) and a is in rr and b is in rr and theta is in rr and a is less than theta and theta is less than b, then if for every x in ccint(a, b), f(theta) is at most f(x), then if f is differentiable at theta, with derivative l, then l equals 0}
{\small\[
\forall f, a, b, \theta, l.\; (((f \in (\mathbb{R} \to \mathbb{R})) \wedge ((a \in \mathbb{R}) \wedge ((b \in \mathbb{R}) \wedge ((\theta \in \mathbb{R}) \wedge ((a < \theta) \wedge (\theta < b)))))) \Rightarrow (\forall x \in \operatorname{ccint}(a, b).\; (f(\theta) \leq f(x)) \Rightarrow (\operatorname{is\text{-}diff\text{-}at}(f, \theta, l) \Rightarrow (l = 0))))
\]}
\noindent\textbf{Machine-checked} (trust: proof).

\subsection{Mean value theorems}

\subsubsection*{\texttt{rolle}}
\emph{for every h, a, b, if h is in fun(rr, rr) and a is in rr and b is in rr and a is less than b, then if for every x in ccint(a, b), h is continuous at x, then if for every x, if x is in rr and a is less than x and x is less than b, then there is some l such that h is differentiable at x, with derivative l, then if h(a) equals h(b), then there is some theta in rr such that a is less than theta and theta is less than b and h is differentiable at theta, with derivative 0}
{\small\[
\forall h, a, b.\; (((h \in (\mathbb{R} \to \mathbb{R})) \wedge ((a \in \mathbb{R}) \wedge ((b \in \mathbb{R}) \wedge (a < b)))) \Rightarrow (\forall x \in \operatorname{ccint}(a, b).\; \operatorname{is\text{-}continuous\text{-}at}(\operatorname{rr\text{-}ms}, \operatorname{rr\text{-}ms}, h, x) \Rightarrow (\forall x.\; (((x \in \mathbb{R}) \wedge ((a < x) \wedge (x < b))) \Rightarrow \exists l.\; \operatorname{is\text{-}diff\text{-}at}(h, x, l)) \Rightarrow ((h(a) = h(b)) \Rightarrow \exists \theta \in \mathbb{R}.\; ((a < \theta) \wedge ((\theta < b) \wedge \operatorname{is\text{-}diff\text{-}at}(h, \theta, 0)))))))
\]}
\noindent\textbf{Machine-checked} (trust: proof). The extreme value theorem gives an interior maximum or minimum; Fermat's lemma kills the derivative there.  A three-way case split handles the degenerate case where both extrema sit at the (equal) endpoints.

\subsubsection*{\texttt{mvt}}
\emph{for every f, a, b, if f is in fun(rr, rr) and a is in rr and b is in rr and a is less than b, then if for every x in ccint(a, b), f is continuous at x, then if for every x, if x is in rr and a is less than x and x is less than b, then there is some l such that f is differentiable at x, with derivative l, then there is some theta in rr such that a is less than theta and theta is less than b and there is some l such that f is differentiable at theta, with derivative l and l * (b - a) equals f(b) - f(a)}
{\small\[
\forall f, a, b.\; (((f \in (\mathbb{R} \to \mathbb{R})) \wedge ((a \in \mathbb{R}) \wedge ((b \in \mathbb{R}) \wedge (a < b)))) \Rightarrow (\forall x \in \operatorname{ccint}(a, b).\; \operatorname{is\text{-}continuous\text{-}at}(\operatorname{rr\text{-}ms}, \operatorname{rr\text{-}ms}, f, x) \Rightarrow (\forall x.\; (((x \in \mathbb{R}) \wedge ((a < x) \wedge (x < b))) \Rightarrow \exists l.\; \operatorname{is\text{-}diff\text{-}at}(f, x, l)) \Rightarrow \exists \theta \in \mathbb{R}.\; ((a < \theta) \wedge ((\theta < b) \wedge \exists l.\; (\operatorname{is\text{-}diff\text{-}at}(f, \theta, l) \wedge ((l \cdot (b - a)) = (f(b) - f(a)))))))))
\]}
\noindent\textbf{Machine-checked} (trust: proof). Apply Rolle to the tilted auxiliary g(x) = f(x) - L(x), where L is the secant line through the two endpoints.

\subsubsection*{\texttt{generalized-mvt}}
\emph{for every f, g, a, b, if f is in fun(rr, rr) and g is in fun(rr, rr) and a is in rr and b is in rr and a is less than b, then if for every x in ccint(a, b), f is continuous at x and g is continuous at x, then if for every x, if x is in rr and a is less than x and x is less than b, then there is some l such that f is differentiable at x, with derivative l and there is some m such that g is differentiable at x, with derivative m, then there is some theta in rr such that a is less than theta and theta is less than b and there is some l, m such that f is differentiable at theta, with derivative l and g is differentiable at theta, with derivative m and l * (g(b) - g(a)) equals m * (f(b) - f(a))}
{\small\[
\forall f, g, a, b.\; (((f \in (\mathbb{R} \to \mathbb{R})) \wedge ((g \in (\mathbb{R} \to \mathbb{R})) \wedge ((a \in \mathbb{R}) \wedge ((b \in \mathbb{R}) \wedge (a < b))))) \Rightarrow (\forall x \in \operatorname{ccint}(a, b).\; (\operatorname{is\text{-}continuous\text{-}at}(\operatorname{rr\text{-}ms}, \operatorname{rr\text{-}ms}, f, x) \wedge \operatorname{is\text{-}continuous\text{-}at}(\operatorname{rr\text{-}ms}, \operatorname{rr\text{-}ms}, g, x)) \Rightarrow (\forall x.\; (((x \in \mathbb{R}) \wedge ((a < x) \wedge (x < b))) \Rightarrow (\exists l.\; \operatorname{is\text{-}diff\text{-}at}(f, x, l) \wedge \exists m.\; \operatorname{is\text{-}diff\text{-}at}(g, x, m))) \Rightarrow \exists \theta \in \mathbb{R}.\; ((a < \theta) \wedge ((\theta < b) \wedge \exists l, m.\; (\operatorname{is\text{-}diff\text{-}at}(f, \theta, l) \wedge (\operatorname{is\text{-}diff\text{-}at}(g, \theta, m) \wedge ((l \cdot (g(b) - g(a))) = (m \cdot (f(b) - f(a)))))))))))
\]}
\noindent\textbf{Machine-checked} (trust: proof). Rolle applied to a two-function auxiliary -- the Cauchy mean value theorem.

\subsection{Consequences of the mean value theorem}

\subsubsection*{\texttt{deriv-zero-implies-constant}}
\emph{for every f, a, b, if f is in fun(rr, rr) and a is in rr and b is in rr and a is less than b, then if for every x in ccint(a, b), f is continuous at x, then if for every x, if x is in rr and a is less than x and x is less than b, then f is differentiable at x, with derivative 0, then for every u, v, if u is in ccint(a, b) and v is in ccint(a, b), then f(u) equals f(v)}
{\small\[
\forall f, a, b.\; (((f \in (\mathbb{R} \to \mathbb{R})) \wedge ((a \in \mathbb{R}) \wedge ((b \in \mathbb{R}) \wedge (a < b)))) \Rightarrow (\forall x \in \operatorname{ccint}(a, b).\; \operatorname{is\text{-}continuous\text{-}at}(\operatorname{rr\text{-}ms}, \operatorname{rr\text{-}ms}, f, x) \Rightarrow (\forall x.\; (((x \in \mathbb{R}) \wedge ((a < x) \wedge (x < b))) \Rightarrow \operatorname{is\text{-}diff\text{-}at}(f, x, 0)) \Rightarrow \forall u, v.\; (((u \in \operatorname{ccint}(a, b)) \wedge (v \in \operatorname{ccint}(a, b))) \Rightarrow (f(u) = f(v))))))
\]}
\noindent\textbf{Machine-checked} (trust: proof). For any two points the mean value theorem on the subinterval gives f(y)-f(x) = f'(c)*(y-x) = 0.

\subsubsection*{\texttt{mvt-upper-bound}}
\emph{for every f, a, b, m, if f is in fun(rr, rr) and a is in rr and b is in rr and m is in rr and a is less than b, then if for every x in ccint(a, b), f is continuous at x, then if for every x, if x is in rr and a is less than x and x is less than b, then there is some l such that f is differentiable at x, with derivative l and l is at most m, then f(b) - f(a) is at most m * (b - a)}
{\small\[
\forall f, a, b, m.\; (((f \in (\mathbb{R} \to \mathbb{R})) \wedge ((a \in \mathbb{R}) \wedge ((b \in \mathbb{R}) \wedge ((m \in \mathbb{R}) \wedge (a < b))))) \Rightarrow (\forall x \in \operatorname{ccint}(a, b).\; \operatorname{is\text{-}continuous\text{-}at}(\operatorname{rr\text{-}ms}, \operatorname{rr\text{-}ms}, f, x) \Rightarrow (\forall x.\; (((x \in \mathbb{R}) \wedge ((a < x) \wedge (x < b))) \Rightarrow \exists l.\; (\operatorname{is\text{-}diff\text{-}at}(f, x, l) \wedge (l \leq m))) \Rightarrow ((f(b) - f(a)) \leq (m \cdot (b - a))))))
\]}
\noindent\textbf{Machine-checked} (trust: proof).

\subsubsection*{\texttt{mvt-lower-bound}}
\emph{for every f, a, b, m, if f is in fun(rr, rr) and a is in rr and b is in rr and m is in rr and a is less than b, then if for every x in ccint(a, b), f is continuous at x, then if for every x, if x is in rr and a is less than x and x is less than b, then there is some l such that f is differentiable at x, with derivative l and m is at most l, then m * (b - a) is at most f(b) - f(a)}
{\small\[
\forall f, a, b, m.\; (((f \in (\mathbb{R} \to \mathbb{R})) \wedge ((a \in \mathbb{R}) \wedge ((b \in \mathbb{R}) \wedge ((m \in \mathbb{R}) \wedge (a < b))))) \Rightarrow (\forall x \in \operatorname{ccint}(a, b).\; \operatorname{is\text{-}continuous\text{-}at}(\operatorname{rr\text{-}ms}, \operatorname{rr\text{-}ms}, f, x) \Rightarrow (\forall x.\; (((x \in \mathbb{R}) \wedge ((a < x) \wedge (x < b))) \Rightarrow \exists l.\; (\operatorname{is\text{-}diff\text{-}at}(f, x, l) \wedge (m \leq l))) \Rightarrow ((m \cdot (b - a)) \leq (f(b) - f(a))))))
\]}
\noindent\textbf{Machine-checked} (trust: proof).

\subsubsection*{\texttt{deriv-pos-strictly-increasing}}
\emph{for every f, a, b, if f is in fun(rr, rr) and a is in rr and b is in rr and a is less than b, then if for every x in ccint(a, b), f is continuous at x, then if for every x, if x is in rr and a is less than x and x is less than b, then there is some l such that f is differentiable at x, with derivative l and 0 is less than l, then for every u, v, if u is in ccint(a, b) and v is in ccint(a, b) and u is less than v, then f(u) is less than f(v)}
{\small\[
\forall f, a, b.\; (((f \in (\mathbb{R} \to \mathbb{R})) \wedge ((a \in \mathbb{R}) \wedge ((b \in \mathbb{R}) \wedge (a < b)))) \Rightarrow (\forall x \in \operatorname{ccint}(a, b).\; \operatorname{is\text{-}continuous\text{-}at}(\operatorname{rr\text{-}ms}, \operatorname{rr\text{-}ms}, f, x) \Rightarrow (\forall x.\; (((x \in \mathbb{R}) \wedge ((a < x) \wedge (x < b))) \Rightarrow \exists l.\; (\operatorname{is\text{-}diff\text{-}at}(f, x, l) \wedge (0 < l))) \Rightarrow \forall u, v.\; (((u \in \operatorname{ccint}(a, b)) \wedge ((v \in \operatorname{ccint}(a, b)) \wedge (u < v))) \Rightarrow (f(u) < f(v))))))
\]}
\noindent\textbf{Machine-checked} (trust: proof). The mean value theorem gives f(y)-f(x) = f'(c)*(y-x) > 0 whenever x < y.

\subsection{Taylor's theorem}

\subsubsection*{\texttt{taylor-poly-at-center}}
\emph{for every f, x, n, if x is in rr, then if n is in nn, then if for every k, if k is in nn and k is at most n, then (nth-deriv(f, k))(x) is in rr, then taylor-poly(f, x, n, x) equals f(x)}
{\small\[
\forall f, x \in \mathbb{R}, n \in \mathbb{N}.\; (\forall k.\; (((k \in \mathbb{N}) \wedge (k \leq n)) \Rightarrow (\operatorname{nth\text{-}deriv}(f, k)(x) \in \mathbb{R})) \Rightarrow (\operatorname{taylor\text{-}poly}(f, x, n, x) = f(x)))
\]}
\noindent\textbf{Machine-checked} (trust: proof).

\subsubsection*{\texttt{taylor-lagrange}}
\emph{for every f, a, x, n, if f is in fun(rr, rr) and a is in rr and x is in rr and n is in nn and a is less than x, then if f is n times differentiable from a to x, then there is some theta in rr such that a is less than theta and theta is less than x and factorial(succ(n)) * (f(x) - taylor-poly(f, a, n, x)) equals (nth-deriv(f, succ(n)))(theta) * (x - a) \textasciicircum{} succ(n)}
{\small\[
\forall f, a, x, n.\; (((f \in (\mathbb{R} \to \mathbb{R})) \wedge ((a \in \mathbb{R}) \wedge ((x \in \mathbb{R}) \wedge ((n \in \mathbb{N}) \wedge (a < x))))) \Rightarrow (\operatorname{taylor\text{-}differentiable}(f, a, x, n) \Rightarrow \exists \theta \in \mathbb{R}.\; ((a < \theta) \wedge ((\theta < x) \wedge ((\operatorname{factorial}((n + 1)) \cdot (f(x) - \operatorname{taylor\text{-}poly}(f, a, n, x))) = (\operatorname{nth\text{-}deriv}(f, (n + 1))(\theta) \cdot \operatorname{power}((x - a), (n + 1))))))))
\]}
\noindent\textbf{Machine-checked} (trust: proof). The Cauchy mean value theorem (generalized-mvt) applied to the Taylor remainder against (x-a)\textasciicircum{}(n+1).

\section{Vector calculus}

\subsection{Vocabulary}

\subsubsection*{\texttt{is-vector-space}}
\emph{Definition / vocabulary.}

\subsubsection*{\texttt{is-submodule}}
\emph{for every m, s, s is a submodule of m if and only if s is a subset of vec(m) and vzero(m) is in s and for every x\_ in s, y\_ in s, (vadd(m))(x\_, y\_) is in s and for every x\_ in s, (vneg(m))(x\_) is in s and for every r\_ in carr(scal(m)), x\_ in s, (act(m))(r\_, x\_) is in s}
{\small\[
\forall m, s.\; (\operatorname{is\text{-}submodule}(m, s) \Leftrightarrow ((s \subseteq \operatorname{vec}(m)) \wedge ((\operatorname{vzero}(m) \in s) \wedge (\forall \operatorname{x\_}, \operatorname{y\_} \in s.\; (\operatorname{vadd}(m)(\operatorname{x\_}, \operatorname{y\_}) \in s) \wedge (\forall \operatorname{x\_} \in s.\; (\operatorname{vneg}(m)(\operatorname{x\_}) \in s) \wedge \forall \operatorname{r\_} \in \operatorname{carr}(\operatorname{scal}(m)), \operatorname{x\_} \in s.\; (\operatorname{act}(m)(\operatorname{r\_}, \operatorname{x\_}) \in s))))))
\]}
\noindent\textbf{Definition.}

\subsubsection*{\texttt{is-subspace}}
\emph{for every m, s, s is a subspace of m if and only if s is a submodule of m}
{\small\[
\forall m, s.\; (\operatorname{is\text{-}subspace}(m, s) \Leftrightarrow \operatorname{is\text{-}submodule}(m, s))
\]}
\noindent\textbf{Definition.}

\subsubsection*{\texttt{is-noetherian}}
\emph{for every m, m is Noetherian if and only if m is a module and for every f\_ in fun(nn, power(vec(m))), if for every n\_ in nn, f\_(n\_) is a submodule of m and for every n\_ in nn, f\_(n\_) is a subset of f\_(succ(n\_)), then there is some k\_ in nn such that for every n\_, if n\_ is in nn and k\_ is at most n\_, then f\_(n\_) equals f\_(k\_)}
{\small\[
\forall m.\; (\operatorname{is\text{-}noetherian}(m) \Leftrightarrow (\operatorname{is\text{-}module}(m) \wedge \forall \operatorname{f\_} \in (\mathbb{N} \to \operatorname{power}(\operatorname{vec}(m))).\; ((\forall \operatorname{n\_} \in \mathbb{N}.\; \operatorname{is\text{-}submodule}(m, \operatorname{f\_}(\operatorname{n\_})) \wedge \forall \operatorname{n\_} \in \mathbb{N}.\; (\operatorname{f\_}(\operatorname{n\_}) \subseteq \operatorname{f\_}((\operatorname{n\_} + 1)))) \Rightarrow \exists \operatorname{k\_} \in \mathbb{N}.\; \forall \operatorname{n\_}.\; (((\operatorname{n\_} \in \mathbb{N}) \wedge (\operatorname{k\_} \leq \operatorname{n\_})) \Rightarrow (\operatorname{f\_}(\operatorname{n\_}) = \operatorname{f\_}(\operatorname{k\_}))))))
\]}
\noindent\textbf{Definition.}

\subsubsection*{\texttt{is-finite-dimensional}}
\emph{for every m, m is finite dimensional if and only if m is a vector space and m is Noetherian}
{\small\[
\forall m.\; (\operatorname{is\text{-}finite\text{-}dimensional}(m) \Leftrightarrow (\operatorname{is\text{-}vector\text{-}space}(m) \wedge \operatorname{is\text{-}noetherian}(m)))
\]}
\noindent\textbf{Definition.}

\subsubsection*{\texttt{span-add-one}}
\emph{Definition / vocabulary.}

\subsubsection*{\texttt{is-normed-vector-space}}
\emph{for every s, s is a normed vector space if and only if length(s) equals 7 and scal(s) is a ring and vec(s) is in set and vadd(s) is in fun(cartesian(vec(s), vec(s)), vec(s)) and vzero(s) is in vec(s) and vneg(s) is in fun(vec(s), vec(s)) and act(s) is in fun(cartesian(carr(scal(s)), vec(s)), vec(s)) and vnrm(s) is in fun(vec(s), rr) and vadd(s) is associative on vec(s) and vadd(s) is commutative on vec(s) and vzero(s) is an identity for vadd(s) on vec(s) and every element of vec(s) has an inverse under vadd(s), given by vneg(s), with unit vzero(s) and scal(s) equals normed-field-as-commutative-ring(rr-normed-field) and for every r\_ in carr(scal(s)), x\_ in vec(s), y\_ in vec(s), (act(s))(r\_, (vadd(s))(x\_, y\_)) equals (vadd(s))((act(s))(r\_, x\_), (act(s))(r\_, y\_)) and for every r\_ in carr(scal(s)), s\_ in carr(scal(s)), x\_ in vec(s), (act(s))((add(scal(s)))(r\_, s\_), x\_) equals (vadd(s))((act(s))(r\_, x\_), (act(s))(s\_, x\_)) and for every r\_ in carr(scal(s)), s\_ in carr(scal(s)), x\_ in vec(s), (act(s))((mul(scal(s)))(r\_, s\_), x\_) equals (act(s))(r\_, (act(s))(s\_, x\_)) and for every x\_ in vec(s), (act(s))(one(scal(s)), x\_) equals x\_ and for every x\_ in vec(s), 0 is at most (vnrm(s))(x\_) and for every x\_ in vec(s), (vnrm(s))(x\_) equals 0 if and only if x\_ equals vzero(s) and for every r\_ in carr(scal(s)), x\_ in vec(s), (vnrm(s))((act(s))(r\_, x\_)) equals abs(r\_) * (vnrm(s))(x\_) and for every x\_ in vec(s), y\_ in vec(s), (vnrm(s))((vadd(s))(x\_, y\_)) is at most (vnrm(s))(x\_) + (vnrm(s))(y\_)}
{\small\[
\forall s.\; (\operatorname{is\text{-}normed\text{-}vector\text{-}space}(s) \Leftrightarrow ((\operatorname{length}(s) = 7) \wedge (\operatorname{is\text{-}ring}(\operatorname{scal}(s)) \wedge ((\operatorname{vec}(s) \in \mathbf{Set}) \wedge ((\operatorname{vadd}(s) \in ((\operatorname{vec}(s) \times \operatorname{vec}(s)) \to \operatorname{vec}(s))) \wedge ((\operatorname{vzero}(s) \in \operatorname{vec}(s)) \wedge ((\operatorname{vneg}(s) \in (\operatorname{vec}(s) \to \operatorname{vec}(s))) \wedge ((\operatorname{act}(s) \in ((\operatorname{carr}(\operatorname{scal}(s)) \times \operatorname{vec}(s)) \to \operatorname{vec}(s))) \wedge ((\operatorname{vnrm}(s) \in (\operatorname{vec}(s) \to \mathbb{R})) \wedge (\operatorname{is\text{-}associative}(\operatorname{vadd}(s), \operatorname{vec}(s)) \wedge (\operatorname{is\text{-}commutative}(\operatorname{vadd}(s), \operatorname{vec}(s)) \wedge (\operatorname{is\text{-}identity}(\operatorname{vadd}(s), \operatorname{vzero}(s), \operatorname{vec}(s)) \wedge (\operatorname{has\text{-}inverses}(\operatorname{vadd}(s), \operatorname{vzero}(s), \operatorname{vneg}(s), \operatorname{vec}(s)) \wedge ((\operatorname{scal}(s) = \operatorname{normed\text{-}field\text{-}as\text{-}commutative\text{-}ring}(\operatorname{rr\text{-}normed\text{-}field})) \wedge (\forall \operatorname{r\_} \in \operatorname{carr}(\operatorname{scal}(s)), \operatorname{x\_}, \operatorname{y\_} \in \operatorname{vec}(s).\; (\operatorname{act}(s)(\operatorname{r\_}, \operatorname{vadd}(s)(\operatorname{x\_}, \operatorname{y\_})) = \operatorname{vadd}(s)(\operatorname{act}(s)(\operatorname{r\_}, \operatorname{x\_}), \operatorname{act}(s)(\operatorname{r\_}, \operatorname{y\_}))) \wedge (\forall \operatorname{r\_}, \operatorname{s\_} \in \operatorname{carr}(\operatorname{scal}(s)), \operatorname{x\_} \in \operatorname{vec}(s).\; (\operatorname{act}(s)(\operatorname{add}(\operatorname{scal}(s))(\operatorname{r\_}, \operatorname{s\_}), \operatorname{x\_}) = \operatorname{vadd}(s)(\operatorname{act}(s)(\operatorname{r\_}, \operatorname{x\_}), \operatorname{act}(s)(\operatorname{s\_}, \operatorname{x\_}))) \wedge (\forall \operatorname{r\_}, \operatorname{s\_} \in \operatorname{carr}(\operatorname{scal}(s)), \operatorname{x\_} \in \operatorname{vec}(s).\; (\operatorname{act}(s)(\operatorname{mul}(\operatorname{scal}(s))(\operatorname{r\_}, \operatorname{s\_}), \operatorname{x\_}) = \operatorname{act}(s)(\operatorname{r\_}, \operatorname{act}(s)(\operatorname{s\_}, \operatorname{x\_}))) \wedge (\forall \operatorname{x\_} \in \operatorname{vec}(s).\; (\operatorname{act}(s)(\operatorname{one}(\operatorname{scal}(s)), \operatorname{x\_}) = \operatorname{x\_}) \wedge (\forall \operatorname{x\_} \in \operatorname{vec}(s).\; (0 \leq \operatorname{vnrm}(s)(\operatorname{x\_})) \wedge (\forall \operatorname{x\_} \in \operatorname{vec}(s).\; ((\operatorname{vnrm}(s)(\operatorname{x\_}) = 0) \Leftrightarrow (\operatorname{x\_} = \operatorname{vzero}(s))) \wedge (\forall \operatorname{r\_} \in \operatorname{carr}(\operatorname{scal}(s)), \operatorname{x\_} \in \operatorname{vec}(s).\; (\operatorname{vnrm}(s)(\operatorname{act}(s)(\operatorname{r\_}, \operatorname{x\_})) = (\lvert \operatorname{r\_} \rvert \cdot \operatorname{vnrm}(s)(\operatorname{x\_}))) \wedge \forall \operatorname{x\_}, \operatorname{y\_} \in \operatorname{vec}(s).\; (\operatorname{vnrm}(s)(\operatorname{vadd}(s)(\operatorname{x\_}, \operatorname{y\_})) \leq (\operatorname{vnrm}(s)(\operatorname{x\_}) + \operatorname{vnrm}(s)(\operatorname{y\_}))))))))))))))))))))))))
\]}
\noindent\textbf{Definition.}

\subsubsection*{\texttt{is-linear-functional}}
\emph{for every m, f, f is a linear functional on m if and only if f is in fun(vec(m), rr) and for every x\_ in vec(m), y\_ in vec(m), f((vadd(m))(x\_, y\_)) equals f(x\_) + f(y\_) and for every r\_ in rr, x\_ in vec(m), f((act(m))(r\_, x\_)) equals r\_ * f(x\_)}
{\small\[
\forall m, f.\; (\operatorname{is\text{-}linear\text{-}functional}(m, f) \Leftrightarrow ((f \in (\operatorname{vec}(m) \to \mathbb{R})) \wedge (\forall \operatorname{x\_}, \operatorname{y\_} \in \operatorname{vec}(m).\; (f(\operatorname{vadd}(m)(\operatorname{x\_}, \operatorname{y\_})) = (f(\operatorname{x\_}) + f(\operatorname{y\_}))) \wedge \forall \operatorname{r\_} \in \mathbb{R}, \operatorname{x\_} \in \operatorname{vec}(m).\; (f(\operatorname{act}(m)(\operatorname{r\_}, \operatorname{x\_})) = (\operatorname{r\_} \cdot f(\operatorname{x\_}))))))
\]}
\noindent\textbf{Definition.}

\subsubsection*{\texttt{is-linear-functional-on}}
\emph{for every m, s, f, f is a linear functional on the subspace s of m if and only if f is in fun(s, rr) and for every x\_ in s, y\_ in s, f((vadd(m))(x\_, y\_)) equals f(x\_) + f(y\_) and for every r\_ in rr, x\_ in s, f((act(m))(r\_, x\_)) equals r\_ * f(x\_)}
{\small\[
\forall m, s, f.\; (\operatorname{is\text{-}linear\text{-}functional\text{-}on}(m, s, f) \Leftrightarrow ((f \in (s \to \mathbb{R})) \wedge (\forall \operatorname{x\_}, \operatorname{y\_} \in s.\; (f(\operatorname{vadd}(m)(\operatorname{x\_}, \operatorname{y\_})) = (f(\operatorname{x\_}) + f(\operatorname{y\_}))) \wedge \forall \operatorname{r\_} \in \mathbb{R}, \operatorname{x\_} \in s.\; (f(\operatorname{act}(m)(\operatorname{r\_}, \operatorname{x\_})) = (\operatorname{r\_} \cdot f(\operatorname{x\_}))))))
\]}
\noindent\textbf{Definition.}

\subsubsection*{\texttt{is-bounded-linear-functional}}
\emph{for every m, f, f is a bounded linear functional on m if and only if f is a linear functional on m and there is some c\_ in rr such that 0 is at most c\_ and for every x\_ in vec(m), abs(f(x\_)) is at most c\_ * (vnrm(m))(x\_)}
{\small\[
\forall m, f.\; (\operatorname{is\text{-}bounded\text{-}linear\text{-}functional}(m, f) \Leftrightarrow (\operatorname{is\text{-}linear\text{-}functional}(m, f) \wedge \exists \operatorname{c\_} \in \mathbb{R}.\; ((0 \leq \operatorname{c\_}) \wedge \forall \operatorname{x\_} \in \operatorname{vec}(m).\; (\lvert f(\operatorname{x\_}) \rvert \leq (\operatorname{c\_} \cdot \operatorname{vnrm}(m)(\operatorname{x\_}))))))
\]}
\noindent\textbf{Definition.}

\subsubsection*{\texttt{is-bounded-linear-functional-on}}
\emph{for every m, s, f, f is a bounded linear functional on the subspace s of m if and only if f is a linear functional on the subspace s of m and there is some c\_ in rr such that 0 is at most c\_ and for every x\_ in s, abs(f(x\_)) is at most c\_ * (vnrm(m))(x\_)}
{\small\[
\forall m, s, f.\; (\operatorname{is\text{-}bounded\text{-}linear\text{-}functional\text{-}on}(m, s, f) \Leftrightarrow (\operatorname{is\text{-}linear\text{-}functional\text{-}on}(m, s, f) \wedge \exists \operatorname{c\_} \in \mathbb{R}.\; ((0 \leq \operatorname{c\_}) \wedge \forall \operatorname{x\_} \in s.\; (\lvert f(\operatorname{x\_}) \rvert \leq (\operatorname{c\_} \cdot \operatorname{vnrm}(m)(\operatorname{x\_}))))))
\]}
\noindent\textbf{Definition.}

\subsubsection*{\texttt{dual-norm}}
\emph{Definition / vocabulary.}

\subsubsection*{\texttt{dual-norm-on}}
\emph{Definition / vocabulary.}

\subsubsection*{\texttt{extends-on}}
\emph{for every s, g, f, g agrees with f on s if and only if for every x\_ in s, g(x\_) equals f(x\_)}
{\small\[
\forall s, g, f.\; (\operatorname{extends\text{-}on}(s, g, f) \Leftrightarrow \forall \operatorname{x\_} \in s.\; (g(\operatorname{x\_}) = f(\operatorname{x\_})))
\]}
\noindent\textbf{Definition.}

\subsubsection*{\texttt{npe}}
\emph{for every m, s, f, t, g, g is a norm-preserving extension of f from s to t, in m if and only if t is a submodule of m and s is a subset of t and g is a linear functional on the subspace t of m and g agrees with f on s and for every w\_ in t, abs(g(w\_)) is at most dual-norm-on(m, s, f) * (vnrm(m))(w\_)}
{\small\[
\forall m, s, f, t, g.\; (\operatorname{npe}(m, s, f, t, g) \Leftrightarrow (\operatorname{is\text{-}submodule}(m, t) \wedge ((s \subseteq t) \wedge (\operatorname{is\text{-}linear\text{-}functional\text{-}on}(m, t, g) \wedge (\operatorname{extends\text{-}on}(s, g, f) \wedge \forall \operatorname{w\_} \in t.\; (\lvert g(\operatorname{w\_}) \rvert \leq (\operatorname{dual\text{-}norm\text{-}on}(m, s, f) \cdot \operatorname{vnrm}(m)(\operatorname{w\_}))))))))
\]}
\noindent\textbf{Definition.}

\subsubsection*{\texttt{good-sub}}
\emph{for every m, s, f, t, f has a norm-preserving extension from s to t, in m if and only if there is some g\_ such that g\_ is a norm-preserving extension of f from s to t, in m}
{\small\[
\forall m, s, f, t.\; (\operatorname{good\text{-}sub}(m, s, f, t) \Leftrightarrow \exists \operatorname{g\_}.\; \operatorname{npe}(m, s, f, t, \operatorname{g\_}))
\]}
\noindent\textbf{Definition.}

\subsubsection*{\texttt{line}}
\emph{Definition / vocabulary.}

\subsubsection*{\texttt{nvs-metric-space}}
\emph{Definition / vocabulary.}

\subsubsection*{\texttt{is-diff-at-v}}
\emph{for every m, f, a, l, f is differentiable at a with derivative l, as a curve in m if and only if m is a normed vector space and f is in fun(rr, vec(m)) and a is in rr and l is in vec(m) and there is some phi in fun(rr, vec(m)) such that phi is continuous at a and phi(a) equals l and for every x\_ in rr, (vadd(m))(f(x\_), (vneg(m))(f(a))) equals (act(m))(x\_ - a, phi(x\_))}
{\small\[
\forall m, f, a, l.\; (\operatorname{is\text{-}diff\text{-}at\text{-}v}(m, f, a, l) \Leftrightarrow (\operatorname{is\text{-}normed\text{-}vector\text{-}space}(m) \wedge ((f \in (\mathbb{R} \to \operatorname{vec}(m))) \wedge ((a \in \mathbb{R}) \wedge ((l \in \operatorname{vec}(m)) \wedge \exists \varphi \in (\mathbb{R} \to \operatorname{vec}(m)).\; (\operatorname{is\text{-}continuous\text{-}at}(\operatorname{rr\text{-}ms}, \operatorname{nvs\text{-}metric\text{-}space}(m), \varphi, a) \wedge ((\varphi(a) = l) \wedge \forall \operatorname{x\_} \in \mathbb{R}.\; (\operatorname{vadd}(m)(f(\operatorname{x\_}), \operatorname{vneg}(m)(f(a))) = \operatorname{act}(m)((\operatorname{x\_} - a), \varphi(\operatorname{x\_}))))))))))
\]}
\noindent\textbf{Definition.}

\subsubsection*{\texttt{deriv-v}}
\emph{Definition / vocabulary.}

\subsubsection*{\texttt{nth-deriv-v}}
\emph{Definition / vocabulary.}

\subsubsection*{\texttt{taylor-poly-v}}
\emph{Definition / vocabulary.}

\subsubsection*{\texttt{taylor-differentiable-v}}
\emph{for every m, f, a, x, n, f is n times differentiable from a to x, as a curve in m if and only if for every k, if k is in nn and k is at most n, then for every t in ccint(a, x), nth-deriv-v(m, f, k) is continuous at t and for every k, if k is in nn and k is at most n, then for every t, if a is less than t and t is less than x, then nth-deriv-v(m, f, k) is differentiable at t with derivative (nth-deriv-v(m, f, succ(k)))(t), as a curve in m}
{\small\[
\forall m, f, a, x, n.\; (\operatorname{taylor\text{-}differentiable\text{-}v}(m, f, a, x, n) \Leftrightarrow (\forall k.\; (((k \in \mathbb{N}) \wedge (k \leq n)) \Rightarrow \forall t \in \operatorname{ccint}(a, x).\; \operatorname{is\text{-}continuous\text{-}at}(\operatorname{rr\text{-}ms}, \operatorname{nvs\text{-}metric\text{-}space}(m), \operatorname{nth\text{-}deriv\text{-}v}(m, f, k), t)) \wedge \forall k.\; (((k \in \mathbb{N}) \wedge (k \leq n)) \Rightarrow \forall t.\; (((a < t) \wedge (t < x)) \Rightarrow \operatorname{is\text{-}diff\text{-}at\text{-}v}(m, \operatorname{nth\text{-}deriv\text{-}v}(m, f, k), t, \operatorname{nth\text{-}deriv\text{-}v}(m, f, (k + 1))(t))))))
\]}
\noindent\textbf{Definition.}

\subsection{Finite-dimensional spaces (noetherian / ascending chain condition)}

\subsubsection*{\texttt{noetherian-set-has-maximal}}
\emph{for every m, if m is Noetherian, then for every sig in set, if sig is a subset of power(vec(m)), then if for every t in sig, t is a submodule of m, then if there is some t0 such that t0 is in sig, then there is some t in sig such that for every u, if u is in sig and t is a subset of u, then t equals u}
{\small\[
\forall m.\; (\operatorname{is\text{-}noetherian}(m) \Rightarrow \forall \mathit{sig} \in \mathbf{Set}.\; ((\mathit{sig} \subseteq \operatorname{power}(\operatorname{vec}(m))) \Rightarrow (\forall t \in \mathit{sig}.\; \operatorname{is\text{-}submodule}(m, t) \Rightarrow (\exists t_{0}.\; (t_{0} \in \mathit{sig}) \Rightarrow \exists t \in \mathit{sig}.\; \forall u.\; (((u \in \mathit{sig}) \wedge (t \subseteq u)) \Rightarrow (t = u))))))
\]}
\noindent\textbf{Machine-checked} (trust: proof). The ascending chain condition: a strictly increasing infinite chain of subspaces would, via dependent choice, exceed every finite dimension.

\subsubsection*{\texttt{hb-good-has-maximal}}
\emph{for every m, s, f, if m is a normed vector space, then if normed-vector-space-as-module(m) is finite dimensional, then if s is a submodule of m, then if f is a bounded linear functional on the subspace s of m, then there is some t such that f has a norm-preserving extension from s to t, in m and for every u, if f has a norm-preserving extension from s to u, in m and t is a subset of u, then t equals u}
{\small\[
\forall m, s, f.\; (\operatorname{is\text{-}normed\text{-}vector\text{-}space}(m) \Rightarrow (\operatorname{is\text{-}finite\text{-}dimensional}(\operatorname{normed\text{-}vector\text{-}space\text{-}as\text{-}module}(m)) \Rightarrow (\operatorname{is\text{-}submodule}(m, s) \Rightarrow (\operatorname{is\text{-}bounded\text{-}linear\text{-}functional\text{-}on}(m, s, f) \Rightarrow \exists t.\; (\operatorname{good\text{-}sub}(m, s, f, t) \wedge \forall u.\; ((\operatorname{good\text{-}sub}(m, s, f, u) \wedge (t \subseteq u)) \Rightarrow (t = u)))))))
\]}
\noindent\textbf{Machine-checked} (trust: proof).

\subsection{Hahn-Banach extension}

\subsubsection*{\texttt{hahn-banach-extend-one}}
\emph{for every m, s, f, v, if m is a normed vector space and s is a submodule of m and f is a bounded linear functional on the subspace s of m and v is in vec(m) and v is not in s, then there is some g\_ such that g\_ is a linear functional on the subspace span-add-one(m, s, v) of m and g\_ agrees with f on s and for every w\_ in span-add-one(m, s, v), abs(g\_(w\_)) is at most dual-norm-on(m, s, f) * (vnrm(m))(w\_)}
{\small\[
\forall m, s, f, v.\; ((\operatorname{is\text{-}normed\text{-}vector\text{-}space}(m) \wedge (\operatorname{is\text{-}submodule}(m, s) \wedge (\operatorname{is\text{-}bounded\text{-}linear\text{-}functional\text{-}on}(m, s, f) \wedge ((v \in \operatorname{vec}(m)) \wedge \neg (v \in s))))) \Rightarrow \exists \operatorname{g\_}.\; (\operatorname{is\text{-}linear\text{-}functional\text{-}on}(m, \operatorname{span\text{-}add\text{-}one}(m, s, v), \operatorname{g\_}) \wedge (\operatorname{extends\text{-}on}(s, \operatorname{g\_}, f) \wedge \forall \operatorname{w\_} \in \operatorname{span\text{-}add\text{-}one}(m, s, v).\; (\lvert \operatorname{g\_}(\operatorname{w\_}) \rvert \leq (\operatorname{dual\text{-}norm\text{-}on}(m, s, f) \cdot \operatorname{vnrm}(m)(\operatorname{w\_}))))))
\]}
\noindent\textbf{Machine-checked} (trust: proof).

\subsubsection*{\texttt{good-step}}
\emph{for every m, s, f, t, x, if m is a normed vector space and f is a bounded linear functional on the subspace s of m and f has a norm-preserving extension from s to t, in m and x is in vec(m) and x is not in t, then f has a norm-preserving extension from s to span-add-one(m, t, x), in m}
{\small\[
\forall m, s, f, t, x.\; ((\operatorname{is\text{-}normed\text{-}vector\text{-}space}(m) \wedge (\operatorname{is\text{-}bounded\text{-}linear\text{-}functional\text{-}on}(m, s, f) \wedge (\operatorname{good\text{-}sub}(m, s, f, t) \wedge ((x \in \operatorname{vec}(m)) \wedge \neg (x \in t))))) \Rightarrow \operatorname{good\text{-}sub}(m, s, f, \operatorname{span\text{-}add\text{-}one}(m, t, x)))
\]}
\noindent\textbf{Machine-checked} (trust: proof).

\subsubsection*{\texttt{hahn-banach}}
\emph{for every m, s, f, if m is a normed vector space and normed-vector-space-as-module(m) is finite dimensional and s is a submodule of m and f is a bounded linear functional on the subspace s of m, then there is some g\_ such that g\_ is a linear functional on the subspace vec(m) of m and g\_ agrees with f on s and for every w\_ in vec(m), abs(g\_(w\_)) is at most dual-norm-on(m, s, f) * (vnrm(m))(w\_)}
{\small\[
\forall m, s, f.\; ((\operatorname{is\text{-}normed\text{-}vector\text{-}space}(m) \wedge (\operatorname{is\text{-}finite\text{-}dimensional}(\operatorname{normed\text{-}vector\text{-}space\text{-}as\text{-}module}(m)) \wedge (\operatorname{is\text{-}submodule}(m, s) \wedge \operatorname{is\text{-}bounded\text{-}linear\text{-}functional\text{-}on}(m, s, f)))) \Rightarrow \exists \operatorname{g\_}.\; (\operatorname{is\text{-}linear\text{-}functional\text{-}on}(m, \operatorname{vec}(m), \operatorname{g\_}) \wedge (\operatorname{extends\text{-}on}(s, \operatorname{g\_}, f) \wedge \forall \operatorname{w\_} \in \operatorname{vec}(m).\; (\lvert \operatorname{g\_}(\operatorname{w\_}) \rvert \leq (\operatorname{dual\text{-}norm\text{-}on}(m, s, f) \cdot \operatorname{vnrm}(m)(\operatorname{w\_}))))))
\]}
\noindent\textbf{Machine-checked} (trust: proof). A maximal norm-preserving extension exists by the ascending chain condition on the finite-dimensional space (via dependent choice on NN); were it proper, the one-step extension lemma would extend it further -- a contradiction.

\subsection{The norm as a supremum of functionals}

\subsubsection*{\texttt{norm-bounded-by-functionals}}
\emph{for every m, f, x, if m is a normed vector space and f is a bounded linear functional on m and x is in vec(m) and dual-norm(m, f) is at most 1, then abs(f(x)) is at most (vnrm(m))(x)}
{\small\[
\forall m, f, x.\; ((\operatorname{is\text{-}normed\text{-}vector\text{-}space}(m) \wedge (\operatorname{is\text{-}bounded\text{-}linear\text{-}functional}(m, f) \wedge ((x \in \operatorname{vec}(m)) \wedge (\operatorname{dual\text{-}norm}(m, f) \leq 1)))) \Rightarrow (\lvert f(x) \rvert \leq \operatorname{vnrm}(m)(x)))
\]}
\noindent\textbf{Machine-checked} (trust: proof).

\subsubsection*{\texttt{norm-attained-by-functional}}
\emph{for every m, x, if m is a normed vector space and normed-vector-space-as-module(m) is finite dimensional and x is in vec(m), then there is some g such that g is a bounded linear functional on m and dual-norm(m, g) is at most 1 and g(x) equals (vnrm(m))(x)}
{\small\[
\forall m, x.\; ((\operatorname{is\text{-}normed\text{-}vector\text{-}space}(m) \wedge (\operatorname{is\text{-}finite\text{-}dimensional}(\operatorname{normed\text{-}vector\text{-}space\text{-}as\text{-}module}(m)) \wedge (x \in \operatorname{vec}(m)))) \Rightarrow \exists g.\; (\operatorname{is\text{-}bounded\text{-}linear\text{-}functional}(m, g) \wedge ((\operatorname{dual\text{-}norm}(m, g) \leq 1) \wedge (g(x) = \operatorname{vnrm}(m)(x)))))
\]}
\noindent\textbf{Machine-checked} (trust: proof).

\subsubsection*{\texttt{norm-as-sup}}
\emph{for every m, x, if m is a normed vector space and normed-vector-space-as-module(m) is finite dimensional and x is in vec(m), then for every f, if f is a bounded linear functional on m, then if dual-norm(m, f) is at most 1, then abs(f(x)) is at most (vnrm(m))(x) and for every d, if d is in rr and for every f, if f is a bounded linear functional on m, then if dual-norm(m, f) is at most 1, then abs(f(x)) is at most d, then (vnrm(m))(x) is at most d}
{\small\[
\forall m, x.\; ((\operatorname{is\text{-}normed\text{-}vector\text{-}space}(m) \wedge (\operatorname{is\text{-}finite\text{-}dimensional}(\operatorname{normed\text{-}vector\text{-}space\text{-}as\text{-}module}(m)) \wedge (x \in \operatorname{vec}(m)))) \Rightarrow (\forall f.\; (\operatorname{is\text{-}bounded\text{-}linear\text{-}functional}(m, f) \Rightarrow ((\operatorname{dual\text{-}norm}(m, f) \leq 1) \Rightarrow (\lvert f(x) \rvert \leq \operatorname{vnrm}(m)(x)))) \wedge \forall d.\; (((d \in \mathbb{R}) \wedge \forall f.\; (\operatorname{is\text{-}bounded\text{-}linear\text{-}functional}(m, f) \Rightarrow ((\operatorname{dual\text{-}norm}(m, f) \leq 1) \Rightarrow (\lvert f(x) \rvert \leq d)))) \Rightarrow (\operatorname{vnrm}(m)(x) \leq d))))
\]}
\noindent\textbf{Machine-checked} (trust: proof). The Hahn-Banach payoff: on the line RR*x a norming functional attains ||x||, while every bounded functional is dominated by the norm, so ||x|| is exactly the supremum of the functional values.

\subsection{Vector-valued Taylor (remainder-norm bound, reduced to scalar)}

\subsubsection*{\texttt{vector-taylor-remainder-bound}}
\emph{for every m, f, a, x, n, if m is a normed vector space and normed-vector-space-as-module(m) is finite dimensional and f is in fun(rr, vec(m)) and a is in rr and x is in rr and n is in nn and a is less than x, then if f is n times differentiable from a to x, as a curve in m, then there is some theta in rr such that a is less than theta and theta is less than x and factorial(succ(n)) * (vnrm(m))((vadd(m))(f(x), (vneg(m))(taylor-poly-v(m, f, a, x, n)))) is at most (vnrm(m))((nth-deriv-v(m, f, succ(n)))(theta)) * (x - a) \textasciicircum{} succ(n)}
{\small\[
\forall m, f, a, x, n.\; ((\operatorname{is\text{-}normed\text{-}vector\text{-}space}(m) \wedge (\operatorname{is\text{-}finite\text{-}dimensional}(\operatorname{normed\text{-}vector\text{-}space\text{-}as\text{-}module}(m)) \wedge ((f \in (\mathbb{R} \to \operatorname{vec}(m))) \wedge ((a \in \mathbb{R}) \wedge ((x \in \mathbb{R}) \wedge ((n \in \mathbb{N}) \wedge (a < x))))))) \Rightarrow (\operatorname{taylor\text{-}differentiable\text{-}v}(m, f, a, x, n) \Rightarrow \exists \theta \in \mathbb{R}.\; ((a < \theta) \wedge ((\theta < x) \wedge ((\operatorname{factorial}((n + 1)) \cdot \operatorname{vnrm}(m)(\operatorname{vadd}(m)(f(x), \operatorname{vneg}(m)(\operatorname{taylor\text{-}poly\text{-}v}(m, f, a, x, n))))) \leq (\operatorname{vnrm}(m)(\operatorname{nth\text{-}deriv\text{-}v}(m, f, (n + 1))(\theta)) \cdot \operatorname{power}((x - a), (n + 1))))))))
\]}
\noindent\textbf{Machine-checked} (trust: proof). Reduce the vector remainder to a scalar: a norming functional g (Hahn-Banach) turns ||remainder|| into g(remainder); scalar Taylor on g o f, then bound g(f\textasciicircum{}(n+1)) by ||f\textasciicircum{}(n+1)||.

\section{Linear algebra}

\subsection{Vocabulary}

\subsubsection*{\texttt{mat}}
\emph{Definition / vocabulary.}

\subsubsection*{\texttt{entry}}
\emph{Definition / vocabulary.}

\subsubsection*{\texttt{matof}}
\emph{Definition / vocabulary.}

\subsubsection*{\texttt{matmul}}
\emph{Definition / vocabulary.}

\subsubsection*{\texttt{matadd}}
\emph{Definition / vocabulary.}

\subsubsection*{\texttt{matneg}}
\emph{Definition / vocabulary.}

\subsubsection*{\texttt{matscale}}
\emph{Definition / vocabulary.}

\subsubsection*{\texttt{identmat}}
\emph{Definition / vocabulary.}

\subsubsection*{\texttt{zeromat}}
\emph{Definition / vocabulary.}

\subsubsection*{\texttt{mat-ring}}
\emph{Definition / vocabulary.}

\subsubsection*{\texttt{submat}}
\emph{Definition / vocabulary.}

\subsubsection*{\texttt{block}}
\emph{Definition / vocabulary.}

\subsubsection*{\texttt{snoc-col}}
\emph{Definition / vocabulary.}

\subsubsection*{\texttt{snoc-row}}
\emph{Definition / vocabulary.}

\subsubsection*{\texttt{matunit}}
\emph{Definition / vocabulary.}

\subsubsection*{\texttt{elem-f}}
\emph{Definition / vocabulary.}

\subsubsection*{\texttt{elem-g}}
\emph{Definition / vocabulary.}

\subsubsection*{\texttt{elem-h}}
\emph{Definition / vocabulary.}

\subsubsection*{\texttt{mat-equiv}}
\emph{for every a, m, n, c, d, c and d are equivalent m-by-n matrices over a if and only if there is some u such that u is an invertible m-by-m matrix over a and there is some v such that v is an invertible n-by-n matrix over a and d equals matmul(a, matmul(a, u, c), v)}
{\small\[
\forall a, m, n, c, d.\; (\operatorname{mat\text{-}equiv}(a, m, n, c, d) \Leftrightarrow \exists u.\; (\operatorname{is\text{-}invertible\text{-}mat}(a, m, u) \wedge \exists v.\; (\operatorname{is\text{-}invertible\text{-}mat}(a, n, v) \wedge (d = \operatorname{matmul}(a, \operatorname{matmul}(a, u, c), v)))))
\]}
\noindent\textbf{Definition.}

\subsubsection*{\texttt{matact}}
\emph{Definition / vocabulary.}

\subsubsection*{\texttt{minor}}
\emph{Definition / vocabulary.}

\subsubsection*{\texttt{det}}
\emph{Definition / vocabulary.}

\subsubsection*{\texttt{interval}}
\emph{Definition / vocabulary.}

\subsection{Matrices as a set, and their entries}

\subsubsection*{\texttt{mat-is-set}}
\emph{for every x in set, m, n, mat(m, n, x) is in set}
{\small\[
\forall x \in \mathbf{Set}, m, n.\; (\operatorname{mat}(m, n, x) \in \mathbf{Set})
\]}
\noindent\textbf{Machine-checked} (trust: proof).

\subsubsection*{\texttt{matrix-sethood}}
\emph{for every s in set, matrix(s) is in set}
{\small\[
\forall s \in \mathbf{Set}.\; (\operatorname{matrix}(s) \in \mathbf{Set})
\]}
\noindent\textbf{Machine-checked} (trust: proof).

\subsubsection*{\texttt{matrix-membership}}
\emph{for every s, m, m is in matrix(s) if and only if m is in tuples(tuples(s)) and for every i, if i is in nn and 1 is at most i and i is at most length(m), then for every j, if j is in nn and 1 is at most j and j is at most length(m), then length(nth(i, m)) equals length(nth(j, m))}
{\small\[
\forall s, m.\; ((m \in \operatorname{matrix}(s)) \Leftrightarrow ((m \in \operatorname{tuples}(\operatorname{tuples}(s))) \wedge \forall i.\; (((i \in \mathbb{N}) \wedge ((1 \leq i) \wedge (i \leq \operatorname{length}(m)))) \Rightarrow \forall j.\; (((j \in \mathbb{N}) \wedge ((1 \leq j) \wedge (j \leq \operatorname{length}(m)))) \Rightarrow (\operatorname{length}(\operatorname{nth}(i, m)) = \operatorname{length}(\operatorname{nth}(j, m)))))))
\]}
\noindent\textbf{Definition.}

\subsubsection*{\texttt{matrix-entry-extensionality}}
\emph{for every m, n, x, p, q, if p is in mat(m, n, x), then if q is in mat(m, n, x), then if for every i in interval(1, m), j in interval(1, n), entry(p, i, j) equals entry(q, i, j), then p equals q}
{\small\[
\forall m, n, x, p, q \in \operatorname{mat}(m, n, x).\; (\forall i \in \operatorname{interval}(1, m), j \in \operatorname{interval}(1, n).\; (\operatorname{entry}(p, i, j) = \operatorname{entry}(q, i, j)) \Rightarrow (p = q))
\]}
\noindent\textbf{Machine-checked} (trust: proof).

\subsubsection*{\texttt{entry-in-carrier}}
\emph{for every m, n, x, p, i, j, if p is in mat(m, n, x), then if i is in interval(1, m), then if j is in interval(1, n), then entry(p, i, j) is in x}
{\small\[
\forall m, n, x, p \in \operatorname{mat}(m, n, x), i \in \operatorname{interval}(1, m), j \in \operatorname{interval}(1, n).\; (\operatorname{entry}(p, i, j) \in x)
\]}
\noindent\textbf{Machine-checked} (trust: proof).

\subsubsection*{\texttt{matof-exists}}
\emph{for every m, n, g, if m is in nn, then if n is in nn, then if for every i\_ in interval(1, m), j\_ in interval(1, n), g(i\_, j\_) is in set, then there is some p in mat(m, n, image(g, cartesian(interval(1, m), interval(1, n)))) such that for every i in interval(1, m), j in interval(1, n), entry(p, i, j) equals g(i, j)}
{\small\[
\forall m, n \in \mathbb{N}, g.\; (\forall \operatorname{i\_} \in \operatorname{interval}(1, m), \operatorname{j\_} \in \operatorname{interval}(1, n).\; (g(\operatorname{i\_}, \operatorname{j\_}) \in \mathbf{Set}) \Rightarrow \exists p \in \operatorname{mat}(m, n, \operatorname{image}(g, (\operatorname{interval}(1, m) \times \operatorname{interval}(1, n)))).\; \forall i \in \operatorname{interval}(1, m), j \in \operatorname{interval}(1, n).\; (\operatorname{entry}(p, i, j) = g(i, j)))
\]}
\noindent\textbf{Machine-checked} (trust: proof).

\subsubsection*{\texttt{matof-in-mat}}
\emph{for every m, n, x, g, if m is in nn, then if n is in nn, then if for every i in interval(1, m), j in interval(1, n), g(i, j) is in x, then matof(m, n, g) is in mat(m, n, x)}
{\small\[
\forall m, n \in \mathbb{N}, x, g.\; (\forall i \in \operatorname{interval}(1, m), j \in \operatorname{interval}(1, n).\; (g(i, j) \in x) \Rightarrow (\operatorname{matof}(m, n, g) \in \operatorname{mat}(m, n, x)))
\]}
\noindent\textbf{Machine-checked} (trust: proof).

\subsubsection*{\texttt{entry-of-matof}}
\emph{for every m, n, g, i, j, if m is in nn, then if n is in nn, then if for every i\_ in interval(1, m), j\_ in interval(1, n), g(i\_, j\_) is in set, then if i is in interval(1, m), then if j is in interval(1, n), then entry(matof(m, n, g), i, j) equals g(i, j)}
{\small\[
\forall m, n \in \mathbb{N}, g, i, j.\; (\forall \operatorname{i\_} \in \operatorname{interval}(1, m), \operatorname{j\_} \in \operatorname{interval}(1, n).\; (g(\operatorname{i\_}, \operatorname{j\_}) \in \mathbf{Set}) \Rightarrow ((i \in \operatorname{interval}(1, m)) \Rightarrow ((j \in \operatorname{interval}(1, n)) \Rightarrow (\operatorname{entry}(\operatorname{matof}(m, n, g), i, j) = g(i, j)))))
\]}
\noindent\textbf{Machine-checked} (trust: proof).

\subsubsection*{\texttt{mat-0-1-nonempty}}
\emph{for every x, there is some p such that p is in mat(0, 1, x)}
{\small\[
\forall x.\; \exists p.\; (p \in \operatorname{mat}(0, 1, x))
\]}
\noindent\textbf{Machine-checked} (trust: proof).

\subsubsection*{\texttt{mat-1-0-nonempty}}
\emph{for every x, there is some p such that p is in mat(1, 0, x)}
{\small\[
\forall x.\; \exists p.\; (p \in \operatorname{mat}(1, 0, x))
\]}
\noindent\textbf{Machine-checked} (trust: proof).

\subsection{Addition, negation, scaling}

\subsubsection*{\texttt{matadd-type}}
\emph{for every ring a, for every m, n, p, q, if p is in mat(m, n, carr(a)), then if q is in mat(m, n, carr(a)), then matadd(a, p, q) is in mat(m, n, carr(a))}
{\small\[
\forall a.\; (\operatorname{is\text{-}ring}(a) \Rightarrow \forall m, n, p, q \in \operatorname{mat}(m, n, \operatorname{carr}(a)).\; (\operatorname{matadd}(a, p, q) \in \operatorname{mat}(m, n, \operatorname{carr}(a))))
\]}
\noindent\textbf{Machine-checked} (trust: proof).

\subsubsection*{\texttt{matadd-entry}}
\emph{for every a, m, n, p, q, if a is a ring, then if p is in mat(m, n, carr(a)), then if q is in mat(m, n, carr(a)), then for every i in interval(1, m), j in interval(1, n), entry(matadd(a, p, q), i, j) equals (add(a))(entry(p, i, j), entry(q, i, j))}
{\small\[
\forall a, m, n, p, q.\; (\operatorname{is\text{-}ring}(a) \Rightarrow ((p \in \operatorname{mat}(m, n, \operatorname{carr}(a))) \Rightarrow ((q \in \operatorname{mat}(m, n, \operatorname{carr}(a))) \Rightarrow \forall i \in \operatorname{interval}(1, m), j \in \operatorname{interval}(1, n).\; (\operatorname{entry}(\operatorname{matadd}(a, p, q), i, j) = \operatorname{add}(a)(\operatorname{entry}(p, i, j), \operatorname{entry}(q, i, j))))))
\]}
\noindent\textbf{Machine-checked} (trust: proof).

\subsubsection*{\texttt{matadd-assoc}}
\emph{for every ring a, for every m, n, p, q, r, if p is in mat(m, n, carr(a)), then if q is in mat(m, n, carr(a)), then if r is in mat(m, n, carr(a)), then matadd(a, matadd(a, p, q), r) equals matadd(a, p, matadd(a, q, r))}
{\small\[
\forall a.\; (\operatorname{is\text{-}ring}(a) \Rightarrow \forall m, n, p, q, r \in \operatorname{mat}(m, n, \operatorname{carr}(a)).\; (\operatorname{matadd}(a, \operatorname{matadd}(a, p, q), r) = \operatorname{matadd}(a, p, \operatorname{matadd}(a, q, r))))
\]}
\noindent\textbf{Machine-checked} (trust: proof).

\subsubsection*{\texttt{matadd-comm}}
\emph{for every ring a, for every m, n, p, q, if p is in mat(m, n, carr(a)), then if q is in mat(m, n, carr(a)), then matadd(a, p, q) equals matadd(a, q, p)}
{\small\[
\forall a.\; (\operatorname{is\text{-}ring}(a) \Rightarrow \forall m, n, p, q \in \operatorname{mat}(m, n, \operatorname{carr}(a)).\; (\operatorname{matadd}(a, p, q) = \operatorname{matadd}(a, q, p)))
\]}
\noindent\textbf{Machine-checked} (trust: proof).

\subsubsection*{\texttt{matadd-zero-left}}
\emph{for every ring a, for every m, n, p in mat(m, n, carr(a)), matadd(a, zeromat(a, m, n), p) equals p}
{\small\[
\forall a.\; (\operatorname{is\text{-}ring}(a) \Rightarrow \forall m, n, p \in \operatorname{mat}(m, n, \operatorname{carr}(a)).\; (\operatorname{matadd}(a, \operatorname{zeromat}(a, m, n), p) = p))
\]}
\noindent\textbf{Machine-checked} (trust: proof).

\subsubsection*{\texttt{matadd-zero-right}}
\emph{for every ring a, for every m, n, p in mat(m, n, carr(a)), matadd(a, p, zeromat(a, m, n)) equals p}
{\small\[
\forall a.\; (\operatorname{is\text{-}ring}(a) \Rightarrow \forall m, n, p \in \operatorname{mat}(m, n, \operatorname{carr}(a)).\; (\operatorname{matadd}(a, p, \operatorname{zeromat}(a, m, n)) = p))
\]}
\noindent\textbf{Machine-checked} (trust: proof).

\subsubsection*{\texttt{matadd-neg-left}}
\emph{for every ring a, for every m, n, p in mat(m, n, carr(a)), matadd(a, matneg(a, p), p) equals zeromat(a, m, n)}
{\small\[
\forall a.\; (\operatorname{is\text{-}ring}(a) \Rightarrow \forall m, n, p \in \operatorname{mat}(m, n, \operatorname{carr}(a)).\; (\operatorname{matadd}(a, \operatorname{matneg}(a, p), p) = \operatorname{zeromat}(a, m, n)))
\]}
\noindent\textbf{Machine-checked} (trust: proof).

\subsubsection*{\texttt{matadd-neg-right}}
\emph{for every ring a, for every m, n, p in mat(m, n, carr(a)), matadd(a, p, matneg(a, p)) equals zeromat(a, m, n)}
{\small\[
\forall a.\; (\operatorname{is\text{-}ring}(a) \Rightarrow \forall m, n, p \in \operatorname{mat}(m, n, \operatorname{carr}(a)).\; (\operatorname{matadd}(a, p, \operatorname{matneg}(a, p)) = \operatorname{zeromat}(a, m, n)))
\]}
\noindent\textbf{Machine-checked} (trust: proof).

\subsubsection*{\texttt{matneg-type}}
\emph{for every ring a, for every m, n, p in mat(m, n, carr(a)), matneg(a, p) is in mat(m, n, carr(a))}
{\small\[
\forall a.\; (\operatorname{is\text{-}ring}(a) \Rightarrow \forall m, n, p \in \operatorname{mat}(m, n, \operatorname{carr}(a)).\; (\operatorname{matneg}(a, p) \in \operatorname{mat}(m, n, \operatorname{carr}(a))))
\]}
\noindent\textbf{Machine-checked} (trust: proof).

\subsubsection*{\texttt{matscale-type}}
\emph{for every ring a, for every m, n, r, p, if r is in carr(a), then if p is in mat(m, n, carr(a)), then matscale(a, r, p) is in mat(m, n, carr(a))}
{\small\[
\forall a.\; (\operatorname{is\text{-}ring}(a) \Rightarrow \forall m, n, r \in \operatorname{carr}(a), p \in \operatorname{mat}(m, n, \operatorname{carr}(a)).\; (\operatorname{matscale}(a, r, p) \in \operatorname{mat}(m, n, \operatorname{carr}(a))))
\]}
\noindent\textbf{Machine-checked} (trust: proof).

\subsubsection*{\texttt{matscale-entry}}
\emph{for every a, m, n, r, p, if a is a ring, then if r is in carr(a), then if p is in mat(m, n, carr(a)), then for every i in interval(1, m), j in interval(1, n), entry(matscale(a, r, p), i, j) equals (mul(a))(r, entry(p, i, j))}
{\small\[
\forall a, m, n, r, p.\; (\operatorname{is\text{-}ring}(a) \Rightarrow ((r \in \operatorname{carr}(a)) \Rightarrow ((p \in \operatorname{mat}(m, n, \operatorname{carr}(a))) \Rightarrow \forall i \in \operatorname{interval}(1, m), j \in \operatorname{interval}(1, n).\; (\operatorname{entry}(\operatorname{matscale}(a, r, p), i, j) = \operatorname{mul}(a)(r, \operatorname{entry}(p, i, j))))))
\]}
\noindent\textbf{Machine-checked} (trust: proof).

\subsection{Multiplication}

\subsubsection*{\texttt{matmul-type}}
\emph{for every ring a, for every m, n, k, p, q, if p is in mat(m, n, carr(a)), then if q is in mat(n, k, carr(a)), then if if n equals 0, then m equals 0 or k equals 0, then matmul(a, p, q) is in mat(m, k, carr(a))}
{\small\[
\forall a.\; (\operatorname{is\text{-}ring}(a) \Rightarrow \forall m, n, k, p \in \operatorname{mat}(m, n, \operatorname{carr}(a)), q \in \operatorname{mat}(n, k, \operatorname{carr}(a)).\; (((n = 0) \Rightarrow ((m = 0) \vee (k = 0))) \Rightarrow (\operatorname{matmul}(a, p, q) \in \operatorname{mat}(m, k, \operatorname{carr}(a)))))
\]}
\noindent\textbf{Machine-checked} (trust: proof).

\subsubsection*{\texttt{matmul-entry}}
\emph{for every ring a, for every m, n, k, p, q, if p is in mat(m, n, carr(a)), then if q is in mat(n, k, carr(a)), then if 1 is at most n, then for every i in interval(1, m), c in interval(1, k), entry(matmul(a, p, q), i, c) equals finsum(ring-additive-ag(a), vnb-lambda(j, interval(1, n), (mul(a))(entry(p, i, j), entry(q, j, c))), interval(1, n))}
{\small\[
\forall a.\; (\operatorname{is\text{-}ring}(a) \Rightarrow \forall m, n, k, p \in \operatorname{mat}(m, n, \operatorname{carr}(a)), q \in \operatorname{mat}(n, k, \operatorname{carr}(a)).\; ((1 \leq n) \Rightarrow \forall i \in \operatorname{interval}(1, m), c \in \operatorname{interval}(1, k).\; (\operatorname{entry}(\operatorname{matmul}(a, p, q), i, c) = \operatorname{finsum}(\operatorname{ring\text{-}additive\text{-}ag}(a), (\lambda j \in \operatorname{interval}(1, n).\; \operatorname{mul}(a)(\operatorname{entry}(p, i, j), \operatorname{entry}(q, j, c))), \operatorname{interval}(1, n)))))
\]}
\noindent\textbf{Machine-checked} (trust: proof).

\subsubsection*{\texttt{matmul-assoc}}
\emph{for every ring a, for every m, n, k, l, p, q, r, if p is in mat(m, n, carr(a)), then if q is in mat(n, k, carr(a)), then if r is in mat(k, l, carr(a)), then if if n equals 0, then m equals 0 or k equals 0 and l equals 0, then if if k equals 0, then l equals 0 or n equals 0 and m equals 0, then matmul(a, matmul(a, p, q), r) equals matmul(a, p, matmul(a, q, r))}
{\small\[
\forall a.\; (\operatorname{is\text{-}ring}(a) \Rightarrow \forall m, n, k, l, p \in \operatorname{mat}(m, n, \operatorname{carr}(a)), q \in \operatorname{mat}(n, k, \operatorname{carr}(a)), r \in \operatorname{mat}(k, l, \operatorname{carr}(a)).\; (((n = 0) \Rightarrow ((m = 0) \vee ((k = 0) \wedge (l = 0)))) \Rightarrow (((k = 0) \Rightarrow ((l = 0) \vee ((n = 0) \wedge (m = 0)))) \Rightarrow (\operatorname{matmul}(a, \operatorname{matmul}(a, p, q), r) = \operatorname{matmul}(a, p, \operatorname{matmul}(a, q, r))))))
\]}
\noindent\textbf{Machine-checked} (trust: proof). Both (PQ)R and P(QR) expand, via matmul-entry twice, to the same canonical double sum; finsum-fubini interchanges the two summation orders.

\subsubsection*{\texttt{matmul-left-dist-guarded}}
\emph{for every ring a, for every m, n, k, p, q, r, if p is in mat(m, n, carr(a)), then if q is in mat(n, k, carr(a)), then if r is in mat(n, k, carr(a)), then if if n equals 0, then m equals 0 or k equals 0, then matmul(a, p, matadd(a, q, r)) equals matadd(a, matmul(a, p, q), matmul(a, p, r))}
{\small\[
\forall a.\; (\operatorname{is\text{-}ring}(a) \Rightarrow \forall m, n, k, p \in \operatorname{mat}(m, n, \operatorname{carr}(a)), q, r \in \operatorname{mat}(n, k, \operatorname{carr}(a)).\; (((n = 0) \Rightarrow ((m = 0) \vee (k = 0))) \Rightarrow (\operatorname{matmul}(a, p, \operatorname{matadd}(a, q, r)) = \operatorname{matadd}(a, \operatorname{matmul}(a, p, q), \operatorname{matmul}(a, p, r)))))
\]}
\noindent\textbf{Machine-checked} (trust: proof).

\subsubsection*{\texttt{matmul-right-dist-guarded}}
\emph{for every ring a, for every m, n, k, p, q, r, if p is in mat(m, n, carr(a)), then if q is in mat(m, n, carr(a)), then if r is in mat(n, k, carr(a)), then if if n equals 0, then m equals 0 or k equals 0, then matmul(a, matadd(a, p, q), r) equals matadd(a, matmul(a, p, r), matmul(a, q, r))}
{\small\[
\forall a.\; (\operatorname{is\text{-}ring}(a) \Rightarrow \forall m, n, k, p, q \in \operatorname{mat}(m, n, \operatorname{carr}(a)), r \in \operatorname{mat}(n, k, \operatorname{carr}(a)).\; (((n = 0) \Rightarrow ((m = 0) \vee (k = 0))) \Rightarrow (\operatorname{matmul}(a, \operatorname{matadd}(a, p, q), r) = \operatorname{matadd}(a, \operatorname{matmul}(a, p, r), \operatorname{matmul}(a, q, r)))))
\]}
\noindent\textbf{Machine-checked} (trust: proof).

\subsubsection*{\texttt{matprod-summand-type}}
\emph{for every ring a, for every m, n, kc, p, q, i, c, if p is in mat(m, n, carr(a)), then if q is in mat(n, kc, carr(a)), then if i is in interval(1, m), then if c is in interval(1, kc), then vnb-lambda(j, interval(1, n), (mul(a))(entry(p, i, j), entry(q, j, c))) is in fun(interval(1, n), carr(ring-additive-ag(a)))}
{\small\[
\forall a.\; (\operatorname{is\text{-}ring}(a) \Rightarrow \forall m, n, \mathit{kc}, p \in \operatorname{mat}(m, n, \operatorname{carr}(a)), q \in \operatorname{mat}(n, \mathit{kc}, \operatorname{carr}(a)), i \in \operatorname{interval}(1, m), c \in \operatorname{interval}(1, \mathit{kc}).\; ((\lambda j \in \operatorname{interval}(1, n).\; \operatorname{mul}(a)(\operatorname{entry}(p, i, j), \operatorname{entry}(q, j, c))) \in (\operatorname{interval}(1, n) \to \operatorname{carr}(\operatorname{ring\text{-}additive\text{-}ag}(a)))))
\]}
\noindent\textbf{Machine-checked} (trust: proof).

\subsubsection*{\texttt{matmul-assoc-summand-type}}
\emph{for every ring a, for every m, n, k, l, p, q, r, row, col, if p is in mat(m, n, carr(a)), then if q is in mat(n, k, carr(a)), then if r is in mat(k, l, carr(a)), then if row is in interval(1, m), then if col is in interval(1, l), then vnb-lambda(z, cartesian(interval(1, k), interval(1, n)), (mul(a))((mul(a))(entry(p, row, nth(2, z)), entry(q, nth(2, z), nth(1, z))), entry(r, nth(1, z), col))) is in fun(cartesian(interval(1, k), interval(1, n)), carr(ring-additive-ag(a)))}
{\small\[
\forall a.\; (\operatorname{is\text{-}ring}(a) \Rightarrow \forall m, n, k, l, p \in \operatorname{mat}(m, n, \operatorname{carr}(a)), q \in \operatorname{mat}(n, k, \operatorname{carr}(a)), r \in \operatorname{mat}(k, l, \operatorname{carr}(a)), \mathit{row} \in \operatorname{interval}(1, m), \mathit{col} \in \operatorname{interval}(1, l).\; ((\lambda z \in (\operatorname{interval}(1, k) \times \operatorname{interval}(1, n)).\; \operatorname{mul}(a)(\operatorname{mul}(a)(\operatorname{entry}(p, \mathit{row}, \operatorname{nth}(2, z)), \operatorname{entry}(q, \operatorname{nth}(2, z), \operatorname{nth}(1, z))), \operatorname{entry}(r, \operatorname{nth}(1, z), \mathit{col}))) \in ((\operatorname{interval}(1, k) \times \operatorname{interval}(1, n)) \to \operatorname{carr}(\operatorname{ring\text{-}additive\text{-}ag}(a)))))
\]}
\noindent\textbf{Machine-checked} (trust: proof).

\subsection{The identity and zero matrices}

\subsubsection*{\texttt{identmat-type}}
\emph{for every ring a, for every n in nn, identmat(a, n) is in mat(n, n, carr(a))}
{\small\[
\forall a.\; (\operatorname{is\text{-}ring}(a) \Rightarrow \forall n \in \mathbb{N}.\; (\operatorname{identmat}(a, n) \in \operatorname{mat}(n, n, \operatorname{carr}(a))))
\]}
\noindent\textbf{Machine-checked} (trust: proof).

\subsubsection*{\texttt{identmat-entry-diag}}
\emph{for every a, n, i, if a is a ring, then if n is in nn, then if i is in interval(1, n), then entry(identmat(a, n), i, i) equals one(a)}
{\small\[
\forall a, n, i.\; (\operatorname{is\text{-}ring}(a) \Rightarrow ((n \in \mathbb{N}) \Rightarrow ((i \in \operatorname{interval}(1, n)) \Rightarrow (\operatorname{entry}(\operatorname{identmat}(a, n), i, i) = \operatorname{one}(a)))))
\]}
\noindent\textbf{Machine-checked} (trust: proof).

\subsubsection*{\texttt{identmat-entry-off}}
\emph{for every a, n, i, j, if a is a ring, then if n is in nn, then if i is in interval(1, n), then if j is in interval(1, n), then if i is not equal to j, then entry(identmat(a, n), i, j) equals zero(a)}
{\small\[
\forall a, n, i, j.\; (\operatorname{is\text{-}ring}(a) \Rightarrow ((n \in \mathbb{N}) \Rightarrow ((i \in \operatorname{interval}(1, n)) \Rightarrow ((j \in \operatorname{interval}(1, n)) \Rightarrow (\neg (i = j) \Rightarrow (\operatorname{entry}(\operatorname{identmat}(a, n), i, j) = \operatorname{zero}(a)))))))
\]}
\noindent\textbf{Machine-checked} (trust: proof).

\subsubsection*{\texttt{entry-of-identmat}}
\emph{for every a, n, i, j, if a is a ring, then if n is in nn, then if i is in interval(1, n), then if j is in interval(1, n), then entry(identmat(a, n), i, j) equals if(i = j, one(a), zero(a))}
{\small\[
\forall a, n, i, j.\; (\operatorname{is\text{-}ring}(a) \Rightarrow ((n \in \mathbb{N}) \Rightarrow ((i \in \operatorname{interval}(1, n)) \Rightarrow ((j \in \operatorname{interval}(1, n)) \Rightarrow (\operatorname{entry}(\operatorname{identmat}(a, n), i, j) = \left(\text{if }(i = j)\text{ then }\operatorname{one}(a)\text{ else }\operatorname{zero}(a)\right))))))
\]}
\noindent\textbf{Machine-checked} (trust: proof).

\subsubsection*{\texttt{entry-of-zeromat}}
\emph{for every a, m, n, i, j, if a is a ring, then if m is in nn, then if n is in nn, then if i is in interval(1, m), then if j is in interval(1, n), then entry(zeromat(a, m, n), i, j) equals zero(a)}
{\small\[
\forall a, m, n, i, j.\; (\operatorname{is\text{-}ring}(a) \Rightarrow ((m \in \mathbb{N}) \Rightarrow ((n \in \mathbb{N}) \Rightarrow ((i \in \operatorname{interval}(1, m)) \Rightarrow ((j \in \operatorname{interval}(1, n)) \Rightarrow (\operatorname{entry}(\operatorname{zeromat}(a, m, n), i, j) = \operatorname{zero}(a)))))))
\]}
\noindent\textbf{Machine-checked} (trust: proof).

\subsubsection*{\texttt{identmat-left-identity}}
\emph{for every ring a, for every m, n, p in mat(m, n, carr(a)), matmul(a, identmat(a, m), p) equals p}
{\small\[
\forall a.\; (\operatorname{is\text{-}ring}(a) \Rightarrow \forall m, n, p \in \operatorname{mat}(m, n, \operatorname{carr}(a)).\; (\operatorname{matmul}(a, \operatorname{identmat}(a, m), p) = p))
\]}
\noindent\textbf{Machine-checked} (trust: proof).

\subsubsection*{\texttt{identmat-right-identity}}
\emph{for every ring a, for every m, n, p in mat(m, n, carr(a)), matmul(a, p, identmat(a, n)) equals p}
{\small\[
\forall a.\; (\operatorname{is\text{-}ring}(a) \Rightarrow \forall m, n, p \in \operatorname{mat}(m, n, \operatorname{carr}(a)).\; (\operatorname{matmul}(a, p, \operatorname{identmat}(a, n)) = p))
\]}
\noindent\textbf{Machine-checked} (trust: proof).

\subsubsection*{\texttt{identmat-invertible}}
\emph{for every ring a, for every n in nn, identmat(a, n) is an invertible n-by-n matrix over a}
{\small\[
\forall a.\; (\operatorname{is\text{-}ring}(a) \Rightarrow \forall n \in \mathbb{N}.\; \operatorname{is\text{-}invertible\text{-}mat}(a, n, \operatorname{identmat}(a, n)))
\]}
\noindent\textbf{Machine-checked} (trust: proof).

\subsection{The matrix ring}

\subsubsection*{\texttt{mat-ring-is-ring}}
\emph{for every ring a, for every n in nn, mat-ring(a, n) is a ring}
{\small\[
\forall a.\; (\operatorname{is\text{-}ring}(a) \Rightarrow \forall n \in \mathbb{N}.\; \operatorname{is\text{-}ring}(\operatorname{mat\text{-}ring}(a, n)))
\]}
\noindent\textbf{Machine-checked} (trust: proof). Each of the 14 ring axioms is reduced by matrix-entry-extensionality to an entry identity, then discharged against the entrywise ring axioms and the FINSUM distribution / Fubini lemmas.

\subsubsection*{\texttt{mat-ring-carr}}
\emph{for every a, n, if a is a ring, then if n is in nn, then carr(mat-ring(a, n)) equals mat(n, n, carr(a))}
{\small\[
\forall a, n.\; (\operatorname{is\text{-}ring}(a) \Rightarrow ((n \in \mathbb{N}) \Rightarrow (\operatorname{carr}(\operatorname{mat\text{-}ring}(a, n)) = \operatorname{mat}(n, n, \operatorname{carr}(a)))))
\]}
\noindent\textbf{Machine-checked} (trust: proof).

\subsubsection*{\texttt{mat-ring-add}}
\emph{for every a, n, if a is a ring, then if n is in nn, then add(mat-ring(a, n)) equals vnb-lambda([p, q], cartesian(mat(n, n, carr(a)), mat(n, n, carr(a))), matadd(a, p, q))}
{\small\[
\forall a, n.\; (\operatorname{is\text{-}ring}(a) \Rightarrow ((n \in \mathbb{N}) \Rightarrow (\operatorname{add}(\operatorname{mat\text{-}ring}(a, n)) = (\lambda \langle p, q \rangle \in (\operatorname{mat}(n, n, \operatorname{carr}(a)) \times \operatorname{mat}(n, n, \operatorname{carr}(a))).\; \operatorname{matadd}(a, p, q)))))
\]}
\noindent\textbf{Machine-checked} (trust: proof).

\subsubsection*{\texttt{mat-ring-mul}}
\emph{for every a, n, if a is a ring, then if n is in nn, then mul(mat-ring(a, n)) equals vnb-lambda([p, q], cartesian(mat(n, n, carr(a)), mat(n, n, carr(a))), matmul(a, p, q))}
{\small\[
\forall a, n.\; (\operatorname{is\text{-}ring}(a) \Rightarrow ((n \in \mathbb{N}) \Rightarrow (\operatorname{mul}(\operatorname{mat\text{-}ring}(a, n)) = (\lambda \langle p, q \rangle \in (\operatorname{mat}(n, n, \operatorname{carr}(a)) \times \operatorname{mat}(n, n, \operatorname{carr}(a))).\; \operatorname{matmul}(a, p, q)))))
\]}
\noindent\textbf{Machine-checked} (trust: proof).

\subsubsection*{\texttt{mat-ring-zero}}
\emph{for every a, n, if a is a ring, then if n is in nn, then zero(mat-ring(a, n)) equals zeromat(a, n, n)}
{\small\[
\forall a, n.\; (\operatorname{is\text{-}ring}(a) \Rightarrow ((n \in \mathbb{N}) \Rightarrow (\operatorname{zero}(\operatorname{mat\text{-}ring}(a, n)) = \operatorname{zeromat}(a, n, n))))
\]}
\noindent\textbf{Machine-checked} (trust: proof).

\subsubsection*{\texttt{mat-ring-one}}
\emph{for every a, n, if a is a ring, then if n is in nn, then one(mat-ring(a, n)) equals identmat(a, n)}
{\small\[
\forall a, n.\; (\operatorname{is\text{-}ring}(a) \Rightarrow ((n \in \mathbb{N}) \Rightarrow (\operatorname{one}(\operatorname{mat\text{-}ring}(a, n)) = \operatorname{identmat}(a, n))))
\]}
\noindent\textbf{Machine-checked} (trust: proof).

\subsubsection*{\texttt{mat-ring-neg}}
\emph{for every a, n, if a is a ring, then if n is in nn, then neg(mat-ring(a, n)) equals vnb-lambda(p, mat(n, n, carr(a)), matneg(a, p))}
{\small\[
\forall a, n.\; (\operatorname{is\text{-}ring}(a) \Rightarrow ((n \in \mathbb{N}) \Rightarrow (\operatorname{neg}(\operatorname{mat\text{-}ring}(a, n)) = (\lambda p \in \operatorname{mat}(n, n, \operatorname{carr}(a)).\; \operatorname{matneg}(a, p)))))
\]}
\noindent\textbf{Machine-checked} (trust: proof).

\subsubsection*{\texttt{mat-ring-add-fun}}
\emph{for every ring a, for every n in nn, vnb-lambda([p, q], cartesian(mat(n, n, carr(a)), mat(n, n, carr(a))), matadd(a, p, q)) is in fun(cartesian(mat(n, n, carr(a)), mat(n, n, carr(a))), mat(n, n, carr(a)))}
{\small\[
\forall a.\; (\operatorname{is\text{-}ring}(a) \Rightarrow \forall n \in \mathbb{N}.\; ((\lambda \langle p, q \rangle \in (\operatorname{mat}(n, n, \operatorname{carr}(a)) \times \operatorname{mat}(n, n, \operatorname{carr}(a))).\; \operatorname{matadd}(a, p, q)) \in ((\operatorname{mat}(n, n, \operatorname{carr}(a)) \times \operatorname{mat}(n, n, \operatorname{carr}(a))) \to \operatorname{mat}(n, n, \operatorname{carr}(a)))))
\]}
\noindent\textbf{Machine-checked} (trust: proof).

\subsubsection*{\texttt{mat-ring-mul-fun}}
\emph{for every ring a, for every n in nn, vnb-lambda([p, q], cartesian(mat(n, n, carr(a)), mat(n, n, carr(a))), matmul(a, p, q)) is in fun(cartesian(mat(n, n, carr(a)), mat(n, n, carr(a))), mat(n, n, carr(a)))}
{\small\[
\forall a.\; (\operatorname{is\text{-}ring}(a) \Rightarrow \forall n \in \mathbb{N}.\; ((\lambda \langle p, q \rangle \in (\operatorname{mat}(n, n, \operatorname{carr}(a)) \times \operatorname{mat}(n, n, \operatorname{carr}(a))).\; \operatorname{matmul}(a, p, q)) \in ((\operatorname{mat}(n, n, \operatorname{carr}(a)) \times \operatorname{mat}(n, n, \operatorname{carr}(a))) \to \operatorname{mat}(n, n, \operatorname{carr}(a)))))
\]}
\noindent\textbf{Machine-checked} (trust: proof).

\subsection{Cutting and extending: submatrices, blocks, bordered matrices}

\subsubsection*{\texttt{submat-type}}
\emph{for every a, p, q, s, if p is in nn, then if q is in nn, then if s is in mat(succ(p), succ(q), carr(a)), then submat(s, p, q) is in mat(p, q, carr(a))}
{\small\[
\forall a, p, q \in \mathbb{N}, s \in \operatorname{mat}((p + 1), (q + 1), \operatorname{carr}(a)).\; (\operatorname{submat}(s, p, q) \in \operatorname{mat}(p, q, \operatorname{carr}(a)))
\]}
\noindent\textbf{Machine-checked} (trust: proof).

\subsubsection*{\texttt{entry-of-submat}}
\emph{for every s, p, q, i, j, x, if p is in nn, then if q is in nn, then if s is in mat(succ(p), succ(q), x), then if i is in interval(1, p), then if j is in interval(1, q), then entry(submat(s, p, q), i, j) equals entry(s, succ(i), succ(j))}
{\small\[
\forall s, p, q \in \mathbb{N}, i, j, x.\; ((s \in \operatorname{mat}((p + 1), (q + 1), x)) \Rightarrow ((i \in \operatorname{interval}(1, p)) \Rightarrow ((j \in \operatorname{interval}(1, q)) \Rightarrow (\operatorname{entry}(\operatorname{submat}(s, p, q), i, j) = \operatorname{entry}(s, (i + 1), (j + 1))))))
\]}
\noindent\textbf{Machine-checked} (trust: proof).

\subsubsection*{\texttt{entry-of-block}}
\emph{for every p, k, l, i, j, m, n, x, if k is in nn, then if l is in nn, then if p is in mat(m, n, x), then if k is at most m, then if l is at most n, then if i is in interval(1, k), then if j is in interval(1, l), then entry(block(p, k, l), i, j) equals entry(p, i, j)}
{\small\[
\forall p, k, l \in \mathbb{N}, i, j, m, n, x.\; ((p \in \operatorname{mat}(m, n, x)) \Rightarrow ((k \leq m) \Rightarrow ((l \leq n) \Rightarrow ((i \in \operatorname{interval}(1, k)) \Rightarrow ((j \in \operatorname{interval}(1, l)) \Rightarrow (\operatorname{entry}(\operatorname{block}(p, k, l), i, j) = \operatorname{entry}(p, i, j)))))))
\]}
\noindent\textbf{Machine-checked} (trust: proof).

\subsubsection*{\texttt{snoc-col-type}}
\emph{for every x, n, w, v, if n is in nn, then if w is in mat(n, 1, x), then if v is in x, then snoc-col(w, n, v) is in mat(succ(n), 1, x)}
{\small\[
\forall x, n \in \mathbb{N}, w \in \operatorname{mat}(n, 1, x), v \in x.\; (\operatorname{snoc\text{-}col}(w, n, v) \in \operatorname{mat}((n + 1), 1, x))
\]}
\noindent\textbf{Machine-checked} (trust: proof).

\subsubsection*{\texttt{snoc-col-last}}
\emph{for every w, n, v, x, if n is in nn, then if w is in mat(n, 1, x), then if v is in x, then entry(snoc-col(w, n, v), succ(n), 1) equals v}
{\small\[
\forall w, n \in \mathbb{N}, v, x.\; ((w \in \operatorname{mat}(n, 1, x)) \Rightarrow ((v \in x) \Rightarrow (\operatorname{entry}(\operatorname{snoc\text{-}col}(w, n, v), (n + 1), 1) = v)))
\]}
\noindent\textbf{Machine-checked} (trust: proof).

\subsubsection*{\texttt{entry-of-snoc-col}}
\emph{for every w, n, v, i, x, if n is in nn, then if w is in mat(n, 1, x), then if v is in x, then if i is in interval(1, n), then entry(snoc-col(w, n, v), i, 1) equals entry(w, i, 1)}
{\small\[
\forall w, n \in \mathbb{N}, v, i, x.\; ((w \in \operatorname{mat}(n, 1, x)) \Rightarrow ((v \in x) \Rightarrow ((i \in \operatorname{interval}(1, n)) \Rightarrow (\operatorname{entry}(\operatorname{snoc\text{-}col}(w, n, v), i, 1) = \operatorname{entry}(w, i, 1)))))
\]}
\noindent\textbf{Machine-checked} (trust: proof).

\subsubsection*{\texttt{snoc-row-type}}
\emph{for every x, n, c, r, if n is in nn, then if c is in mat(1, n, x), then if r is in x, then snoc-row(c, n, r) is in mat(1, succ(n), x)}
{\small\[
\forall x, n \in \mathbb{N}, c \in \operatorname{mat}(1, n, x), r \in x.\; (\operatorname{snoc\text{-}row}(c, n, r) \in \operatorname{mat}(1, (n + 1), x))
\]}
\noindent\textbf{Machine-checked} (trust: proof).

\subsubsection*{\texttt{snoc-row-last}}
\emph{for every c, n, r, x, if n is in nn, then if c is in mat(1, n, x), then if r is in x, then entry(snoc-row(c, n, r), 1, succ(n)) equals r}
{\small\[
\forall c, n \in \mathbb{N}, r, x.\; ((c \in \operatorname{mat}(1, n, x)) \Rightarrow ((r \in x) \Rightarrow (\operatorname{entry}(\operatorname{snoc\text{-}row}(c, n, r), 1, (n + 1)) = r)))
\]}
\noindent\textbf{Machine-checked} (trust: proof).

\subsubsection*{\texttt{entry-of-snoc-row}}
\emph{for every c, n, r, j, x, if n is in nn, then if c is in mat(1, n, x), then if r is in x, then if j is in interval(1, n), then entry(snoc-row(c, n, r), 1, j) equals entry(c, 1, j)}
{\small\[
\forall c, n \in \mathbb{N}, r, j, x.\; ((c \in \operatorname{mat}(1, n, x)) \Rightarrow ((r \in x) \Rightarrow ((j \in \operatorname{interval}(1, n)) \Rightarrow (\operatorname{entry}(\operatorname{snoc\text{-}row}(c, n, r), 1, j) = \operatorname{entry}(c, 1, j)))))
\]}
\noindent\textbf{Machine-checked} (trust: proof).

\subsection{The matrix unit and its column shift (Lemma 3.3)}

\subsubsection*{\texttt{matunit-type}}
\emph{for every ring a, for every n, k, l, if n is in nn, then matunit(a, n, k, l) is in mat(n, n, carr(a))}
{\small\[
\forall a.\; (\operatorname{is\text{-}ring}(a) \Rightarrow \forall n \in \mathbb{N}, k, l.\; (\operatorname{matunit}(a, n, k, l) \in \operatorname{mat}(n, n, \operatorname{carr}(a))))
\]}
\noindent\textbf{Machine-checked} (trust: proof).

\subsubsection*{\texttt{matunit-entry-k-row}}
\emph{for every a, n, k, l, c, if a is a ring, then if n is in nn, then if k is in interval(1, n), then if c is in interval(1, n), then entry(matunit(a, n, k, l), k, c) equals if(c = l, one(a), zero(a))}
{\small\[
\forall a, n, k, l, c.\; (\operatorname{is\text{-}ring}(a) \Rightarrow ((n \in \mathbb{N}) \Rightarrow ((k \in \operatorname{interval}(1, n)) \Rightarrow ((c \in \operatorname{interval}(1, n)) \Rightarrow (\operatorname{entry}(\operatorname{matunit}(a, n, k, l), k, c) = \left(\text{if }(c = l)\text{ then }\operatorname{one}(a)\text{ else }\operatorname{zero}(a)\right))))))
\]}
\noindent\textbf{Machine-checked} (trust: proof).

\subsubsection*{\texttt{matunit-entry-off-row}}
\emph{for every a, n, k, l, j, c, if a is a ring, then if n is in nn, then if j is in interval(1, n), then if c is in interval(1, n), then if j is not equal to k, then entry(matunit(a, n, k, l), j, c) equals zero(a)}
{\small\[
\forall a, n, k, l, j, c.\; (\operatorname{is\text{-}ring}(a) \Rightarrow ((n \in \mathbb{N}) \Rightarrow ((j \in \operatorname{interval}(1, n)) \Rightarrow ((c \in \operatorname{interval}(1, n)) \Rightarrow (\neg (j = k) \Rightarrow (\operatorname{entry}(\operatorname{matunit}(a, n, k, l), j, c) = \operatorname{zero}(a)))))))
\]}
\noindent\textbf{Machine-checked} (trust: proof).

\subsubsection*{\texttt{entry-of-matunit}}
\emph{for every a, n, k, l, i, j, if a is a ring, then if n is in nn, then if i is in interval(1, n), then if j is in interval(1, n), then entry(matunit(a, n, k, l), i, j) equals if(i = k and j = l, one(a), zero(a))}
{\small\[
\forall a, n, k, l, i, j.\; (\operatorname{is\text{-}ring}(a) \Rightarrow ((n \in \mathbb{N}) \Rightarrow ((i \in \operatorname{interval}(1, n)) \Rightarrow ((j \in \operatorname{interval}(1, n)) \Rightarrow (\operatorname{entry}(\operatorname{matunit}(a, n, k, l), i, j) = \left(\text{if }((i = k) \wedge (j = l))\text{ then }\operatorname{one}(a)\text{ else }\operatorname{zero}(a)\right))))))
\]}
\noindent\textbf{Machine-checked} (trust: proof).

\subsubsection*{\texttt{matunit-summand-type}}
\emph{for every ring a, for every m, n, p, k, l, i, c, if p is in mat(m, n, carr(a)), then if i is in interval(1, m), then if c is in interval(1, n), then vnb-lambda(j, interval(1, n), (mul(a))(entry(p, i, j), entry(matunit(a, n, k, l), j, c))) is in fun(interval(1, n), carr(ring-additive-ag(a)))}
{\small\[
\forall a.\; (\operatorname{is\text{-}ring}(a) \Rightarrow \forall m, n, p \in \operatorname{mat}(m, n, \operatorname{carr}(a)), k, l, i \in \operatorname{interval}(1, m), c \in \operatorname{interval}(1, n).\; ((\lambda j \in \operatorname{interval}(1, n).\; \operatorname{mul}(a)(\operatorname{entry}(p, i, j), \operatorname{entry}(\operatorname{matunit}(a, n, k, l), j, c))) \in (\operatorname{interval}(1, n) \to \operatorname{carr}(\operatorname{ring\text{-}additive\text{-}ag}(a)))))
\]}
\noindent\textbf{Machine-checked} (trust: proof).

\subsubsection*{\texttt{matunit-col-shift}}
\emph{for every ring a, for every m, n, p, k, l, if p is in mat(m, n, carr(a)), then if k is in interval(1, n), then if l is in interval(1, n), then for every i in interval(1, m), c in interval(1, n), entry(matmul(a, p, matunit(a, n, k, l)), i, c) equals if(c = l, entry(p, i, k), zero(a))}
{\small\[
\forall a.\; (\operatorname{is\text{-}ring}(a) \Rightarrow \forall m, n, p \in \operatorname{mat}(m, n, \operatorname{carr}(a)), k, l \in \operatorname{interval}(1, n).\; \forall i \in \operatorname{interval}(1, m), c \in \operatorname{interval}(1, n).\; (\operatorname{entry}(\operatorname{matmul}(a, p, \operatorname{matunit}(a, n, k, l)), i, c) = \left(\text{if }(c = l)\text{ then }\operatorname{entry}(p, i, k)\text{ else }\operatorname{zero}(a)\right)))
\]}
\noindent\textbf{Machine-checked} (trust: proof). matmul-entry expands (P.E[k,l])\_\{ic\} to FINSUM\_j P\_\{ij\}.E[k,l]\_\{jc\}; off the index k the summand vanishes, so finsum-single-support collapses the sum, and an excluded-middle split on c=l finishes.

\subsection{Elementary column operations (Prop 3.5)}

\subsubsection*{\texttt{elem-f-type}}
\emph{for every ring a, for every n, k, l, if n is in nn, then elem-f(a, n, k, l) is in mat(n, n, carr(a))}
{\small\[
\forall a.\; (\operatorname{is\text{-}ring}(a) \Rightarrow \forall n \in \mathbb{N}, k, l.\; (\operatorname{elem\text{-}f}(a, n, k, l) \in \operatorname{mat}(n, n, \operatorname{carr}(a))))
\]}
\noindent\textbf{Machine-checked} (trust: proof).

\subsubsection*{\texttt{elem-g-type}}
\emph{for every ring a, for every n, r, k, l, if n is in nn, then if r is in carr(a), then elem-g(a, n, r, k, l) is in mat(n, n, carr(a))}
{\small\[
\forall a.\; (\operatorname{is\text{-}ring}(a) \Rightarrow \forall n \in \mathbb{N}, r \in \operatorname{carr}(a), k, l.\; (\operatorname{elem\text{-}g}(a, n, r, k, l) \in \operatorname{mat}(n, n, \operatorname{carr}(a))))
\]}
\noindent\textbf{Machine-checked} (trust: proof).

\subsubsection*{\texttt{elem-h-type}}
\emph{for every ring a, for every n, r, k, if n is in nn, then if r is in carr(a), then elem-h(a, n, r, k) is in mat(n, n, carr(a))}
{\small\[
\forall a.\; (\operatorname{is\text{-}ring}(a) \Rightarrow \forall n \in \mathbb{N}, r \in \operatorname{carr}(a), k.\; (\operatorname{elem\text{-}h}(a, n, r, k) \in \operatorname{mat}(n, n, \operatorname{carr}(a))))
\]}
\noindent\textbf{Machine-checked} (trust: proof).

\subsubsection*{\texttt{elem-f-action}}
\emph{for every ring a, for every m, n, p, k, l, if p is in mat(m, n, carr(a)), then if k is in interval(1, n), then if l is in interval(1, n), then if k is not equal to l, then for every i in interval(1, m), c in interval(1, n), entry(matmul(a, p, elem-f(a, n, k, l)), i, c) equals if(c = k, entry(p, i, l), if(c = l, entry(p, i, k), entry(p, i, c)))}
{\small\[
\forall a.\; (\operatorname{is\text{-}ring}(a) \Rightarrow \forall m, n, p \in \operatorname{mat}(m, n, \operatorname{carr}(a)), k, l \in \operatorname{interval}(1, n).\; (\neg (k = l) \Rightarrow \forall i \in \operatorname{interval}(1, m), c \in \operatorname{interval}(1, n).\; (\operatorname{entry}(\operatorname{matmul}(a, p, \operatorname{elem\text{-}f}(a, n, k, l)), i, c) = \left(\text{if }(c = k)\text{ then }\operatorname{entry}(p, i, l)\text{ else }\left(\text{if }(c = l)\text{ then }\operatorname{entry}(p, i, k)\text{ else }\operatorname{entry}(p, i, c)\right)\right))))
\]}
\noindent\textbf{Machine-checked} (trust: proof). Right-multiplication by F[k,l] permutes columns k and l: a three-way column split, each branch a single-support finsum collapse of the transposition matrix.

\subsubsection*{\texttt{elem-g-action}}
\emph{for every ring a, for every m, n, p, r, k, l, if p is in mat(m, n, carr(a)), then if r is in carr(a), then if k is in interval(1, n), then if l is in interval(1, n), then if k is not equal to l, then for every i in interval(1, m), c in interval(1, n), entry(matmul(a, p, elem-g(a, n, r, k, l)), i, c) equals if(c = l, (add(a))(entry(p, i, l), (mul(a))(entry(p, i, k), r)), entry(p, i, c))}
{\small\[
\forall a.\; (\operatorname{is\text{-}ring}(a) \Rightarrow \forall m, n, p \in \operatorname{mat}(m, n, \operatorname{carr}(a)), r \in \operatorname{carr}(a), k, l \in \operatorname{interval}(1, n).\; (\neg (k = l) \Rightarrow \forall i \in \operatorname{interval}(1, m), c \in \operatorname{interval}(1, n).\; (\operatorname{entry}(\operatorname{matmul}(a, p, \operatorname{elem\text{-}g}(a, n, r, k, l)), i, c) = \left(\text{if }(c = l)\text{ then }\operatorname{add}(a)(\operatorname{entry}(p, i, l), \operatorname{mul}(a)(\operatorname{entry}(p, i, k), r))\text{ else }\operatorname{entry}(p, i, c)\right))))
\]}
\noindent\textbf{Machine-checked} (trust: proof). The l-column of G = I + r.E[k,l] has TWO supports (the diagonal l and the r-slot k), so a two-point finsum split gives P\_\{il\} + P\_\{ik\}.r; every other column is a single-support identity column.

\subsubsection*{\texttt{elem-h-action}}
\emph{for every ring a, for every m, n, p, r, k, if p is in mat(m, n, carr(a)), then if r is in carr(a), then if k is in interval(1, n), then for every i in interval(1, m), c in interval(1, n), entry(matmul(a, p, elem-h(a, n, r, k)), i, c) equals if(c = k, (mul(a))(entry(p, i, k), r), entry(p, i, c))}
{\small\[
\forall a.\; (\operatorname{is\text{-}ring}(a) \Rightarrow \forall m, n, p \in \operatorname{mat}(m, n, \operatorname{carr}(a)), r \in \operatorname{carr}(a), k \in \operatorname{interval}(1, n).\; \forall i \in \operatorname{interval}(1, m), c \in \operatorname{interval}(1, n).\; (\operatorname{entry}(\operatorname{matmul}(a, p, \operatorname{elem\text{-}h}(a, n, r, k)), i, c) = \left(\text{if }(c = k)\text{ then }\operatorname{mul}(a)(\operatorname{entry}(p, i, k), r)\text{ else }\operatorname{entry}(p, i, c)\right)))
\]}
\noindent\textbf{Machine-checked} (trust: proof). H[r,k] scales column k by r: a single-support finsum collapse at the diagonal, then a case split on c=k.

\subsubsection*{\texttt{entry-of-elem-f}}
\emph{for every a, n, k, l, i, j, if a is a ring, then if n is in nn, then if i is in interval(1, n), then if j is in interval(1, n), then entry(elem-f(a, n, k, l), i, j) equals if(i = k and j = l or i = l and j = k or i = j and not(i = k) and not(i = l), one(a), zero(a))}
{\small\[
\forall a, n, k, l, i, j.\; (\operatorname{is\text{-}ring}(a) \Rightarrow ((n \in \mathbb{N}) \Rightarrow ((i \in \operatorname{interval}(1, n)) \Rightarrow ((j \in \operatorname{interval}(1, n)) \Rightarrow (\operatorname{entry}(\operatorname{elem\text{-}f}(a, n, k, l), i, j) = \left(\text{if }(((i = k) \wedge (j = l)) \vee (((i = l) \wedge (j = k)) \vee ((i = j) \wedge (\neg (i = k) \wedge \neg (i = l)))))\text{ then }\operatorname{one}(a)\text{ else }\operatorname{zero}(a)\right))))))
\]}
\noindent\textbf{Machine-checked} (trust: proof).

\subsubsection*{\texttt{entry-of-elem-g}}
\emph{for every a, n, r, k, l, i, j, if a is a ring, then if n is in nn, then if r is in carr(a), then if i is in interval(1, n), then if j is in interval(1, n), then entry(elem-g(a, n, r, k, l), i, j) equals if(i = j, one(a), if(i = k and j = l, r, zero(a)))}
{\small\[
\forall a, n, r, k, l, i, j.\; (\operatorname{is\text{-}ring}(a) \Rightarrow ((n \in \mathbb{N}) \Rightarrow ((r \in \operatorname{carr}(a)) \Rightarrow ((i \in \operatorname{interval}(1, n)) \Rightarrow ((j \in \operatorname{interval}(1, n)) \Rightarrow (\operatorname{entry}(\operatorname{elem\text{-}g}(a, n, r, k, l), i, j) = \left(\text{if }(i = j)\text{ then }\operatorname{one}(a)\text{ else }\left(\text{if }((i = k) \wedge (j = l))\text{ then }r\text{ else }\operatorname{zero}(a)\right)\right)))))))
\]}
\noindent\textbf{Machine-checked} (trust: proof).

\subsubsection*{\texttt{entry-of-elem-h}}
\emph{for every a, n, r, k, i, j, if a is a ring, then if n is in nn, then if r is in carr(a), then if i is in interval(1, n), then if j is in interval(1, n), then entry(elem-h(a, n, r, k), i, j) equals if(i = j, if(i = k, r, one(a)), zero(a))}
{\small\[
\forall a, n, r, k, i, j.\; (\operatorname{is\text{-}ring}(a) \Rightarrow ((n \in \mathbb{N}) \Rightarrow ((r \in \operatorname{carr}(a)) \Rightarrow ((i \in \operatorname{interval}(1, n)) \Rightarrow ((j \in \operatorname{interval}(1, n)) \Rightarrow (\operatorname{entry}(\operatorname{elem\text{-}h}(a, n, r, k), i, j) = \left(\text{if }(i = j)\text{ then }\left(\text{if }(i = k)\text{ then }r\text{ else }\operatorname{one}(a)\right)\text{ else }\operatorname{zero}(a)\right)))))))
\]}
\noindent\textbf{Machine-checked} (trust: proof).

\subsection{Elementary row operations}

\subsubsection*{\texttt{elem-f-row-action}}
\emph{for every ring a, for every m, n, p, k, l, if p is in mat(m, n, carr(a)), then if k is in interval(1, m), then if l is in interval(1, m), then if k is not equal to l, then for every i in interval(1, m), c in interval(1, n), entry(matmul(a, elem-f(a, m, k, l), p), i, c) equals if(i = k, entry(p, l, c), if(i = l, entry(p, k, c), entry(p, i, c)))}
{\small\[
\forall a.\; (\operatorname{is\text{-}ring}(a) \Rightarrow \forall m, n, p \in \operatorname{mat}(m, n, \operatorname{carr}(a)), k, l \in \operatorname{interval}(1, m).\; (\neg (k = l) \Rightarrow \forall i \in \operatorname{interval}(1, m), c \in \operatorname{interval}(1, n).\; (\operatorname{entry}(\operatorname{matmul}(a, \operatorname{elem\text{-}f}(a, m, k, l), p), i, c) = \left(\text{if }(i = k)\text{ then }\operatorname{entry}(p, l, c)\text{ else }\left(\text{if }(i = l)\text{ then }\operatorname{entry}(p, k, c)\text{ else }\operatorname{entry}(p, i, c)\right)\right))))
\]}
\noindent\textbf{Machine-checked} (trust: proof).

\subsubsection*{\texttt{elem-g-row-action}}
\emph{for every ring a, for every m, n, p, r, k, l, if p is in mat(m, n, carr(a)), then if r is in carr(a), then if k is in interval(1, m), then if l is in interval(1, m), then if k is not equal to l, then for every i in interval(1, m), c in interval(1, n), entry(matmul(a, elem-g(a, m, r, k, l), p), i, c) equals if(i = k, (add(a))(entry(p, k, c), (mul(a))(r, entry(p, l, c))), entry(p, i, c))}
{\small\[
\forall a.\; (\operatorname{is\text{-}ring}(a) \Rightarrow \forall m, n, p \in \operatorname{mat}(m, n, \operatorname{carr}(a)), r \in \operatorname{carr}(a), k, l \in \operatorname{interval}(1, m).\; (\neg (k = l) \Rightarrow \forall i \in \operatorname{interval}(1, m), c \in \operatorname{interval}(1, n).\; (\operatorname{entry}(\operatorname{matmul}(a, \operatorname{elem\text{-}g}(a, m, r, k, l), p), i, c) = \left(\text{if }(i = k)\text{ then }\operatorname{add}(a)(\operatorname{entry}(p, k, c), \operatorname{mul}(a)(r, \operatorname{entry}(p, l, c)))\text{ else }\operatorname{entry}(p, i, c)\right))))
\]}
\noindent\textbf{Machine-checked} (trust: proof).

\subsubsection*{\texttt{elem-h-row-action}}
\emph{for every ring a, for every m, n, p, r, k, if p is in mat(m, n, carr(a)), then if r is in carr(a), then if k is in interval(1, m), then for every i in interval(1, m), c in interval(1, n), entry(matmul(a, elem-h(a, m, r, k), p), i, c) equals if(i = k, (mul(a))(r, entry(p, k, c)), entry(p, i, c))}
{\small\[
\forall a.\; (\operatorname{is\text{-}ring}(a) \Rightarrow \forall m, n, p \in \operatorname{mat}(m, n, \operatorname{carr}(a)), r \in \operatorname{carr}(a), k \in \operatorname{interval}(1, m).\; \forall i \in \operatorname{interval}(1, m), c \in \operatorname{interval}(1, n).\; (\operatorname{entry}(\operatorname{matmul}(a, \operatorname{elem\text{-}h}(a, m, r, k), p), i, c) = \left(\text{if }(i = k)\text{ then }\operatorname{mul}(a)(r, \operatorname{entry}(p, k, c))\text{ else }\operatorname{entry}(p, i, c)\right)))
\]}
\noindent\textbf{Machine-checked} (trust: proof).

\subsection{Elementary inverses (Cor 3.6)}

\subsubsection*{\texttt{elem-f-inverse}}
\emph{for every ring a, for every n, k, l, if n is in nn, then if k is in interval(1, n), then if l is in interval(1, n), then if k is not equal to l, then matmul(a, elem-f(a, n, k, l), elem-f(a, n, l, k)) equals identmat(a, n)}
{\small\[
\forall a.\; (\operatorname{is\text{-}ring}(a) \Rightarrow \forall n \in \mathbb{N}, k, l \in \operatorname{interval}(1, n).\; (\neg (k = l) \Rightarrow (\operatorname{matmul}(a, \operatorname{elem\text{-}f}(a, n, k, l), \operatorname{elem\text{-}f}(a, n, l, k)) = \operatorname{identmat}(a, n))))
\]}
\noindent\textbf{Machine-checked} (trust: proof). F[k,l] is its own inverse: composing its column-swap action (Prop 3.5) with F[l,k] swaps the columns back, matching the identity entrywise.

\subsubsection*{\texttt{elem-g-inverse}}
\emph{for every ring a, for every n, r, k, l, if n is in nn, then if r is in carr(a), then if k is in interval(1, n), then if l is in interval(1, n), then if k is not equal to l, then matmul(a, elem-g(a, n, r, k, l), elem-g(a, n, (neg(a))(r), k, l)) equals identmat(a, n)}
{\small\[
\forall a.\; (\operatorname{is\text{-}ring}(a) \Rightarrow \forall n \in \mathbb{N}, r \in \operatorname{carr}(a), k, l \in \operatorname{interval}(1, n).\; (\neg (k = l) \Rightarrow (\operatorname{matmul}(a, \operatorname{elem\text{-}g}(a, n, r, k, l), \operatorname{elem\text{-}g}(a, n, \operatorname{neg}(a)(r), k, l)) = \operatorname{identmat}(a, n))))
\]}
\noindent\textbf{Machine-checked} (trust: proof). G[r,k,l]\textasciicircum{}\{-1\} = G[-r,k,l]: on column l the added term P\_\{ik\}.(-r) cancels the r via r + (-r) = 0; every other column is unchanged.

\subsubsection*{\texttt{elem-h-inverse}}
\emph{for every ring a, for every n, r, s, k, if n is in nn, then if r is in carr(a), then if s is in carr(a), then if (mul(a))(r, s) equals one(a), then if k is in interval(1, n), then matmul(a, elem-h(a, n, r, k), elem-h(a, n, s, k)) equals identmat(a, n)}
{\small\[
\forall a.\; (\operatorname{is\text{-}ring}(a) \Rightarrow \forall n \in \mathbb{N}, r, s \in \operatorname{carr}(a), k.\; ((\operatorname{mul}(a)(r, s) = \operatorname{one}(a)) \Rightarrow ((k \in \operatorname{interval}(1, n)) \Rightarrow (\operatorname{matmul}(a, \operatorname{elem\text{-}h}(a, n, r, k), \operatorname{elem\text{-}h}(a, n, s, k)) = \operatorname{identmat}(a, n)))))
\]}
\noindent\textbf{Machine-checked} (trust: proof). H[r,k]\textasciicircum{}\{-1\} = H[r\textasciicircum{}\{-1\},k]: the (k,k) diagonal entry becomes r.s = 1 by the unit hypothesis; the rest of the diagonal stays 1 and off-diagonal 0.

\subsubsection*{\texttt{elem-f-invertible}}
\emph{for every ring a, for every n, k, l, if n is in nn, then if k is in interval(1, n), then if l is in interval(1, n), then if k is not equal to l, then elem-f(a, n, k, l) is an invertible n-by-n matrix over a}
{\small\[
\forall a.\; (\operatorname{is\text{-}ring}(a) \Rightarrow \forall n \in \mathbb{N}, k, l \in \operatorname{interval}(1, n).\; (\neg (k = l) \Rightarrow \operatorname{is\text{-}invertible\text{-}mat}(a, n, \operatorname{elem\text{-}f}(a, n, k, l))))
\]}
\noindent\textbf{Machine-checked} (trust: proof).

\subsubsection*{\texttt{elem-g-invertible}}
\emph{for every ring a, for every n, r, k, l, if n is in nn, then if r is in carr(a), then if k is in interval(1, n), then if l is in interval(1, n), then if k is not equal to l, then elem-g(a, n, r, k, l) is an invertible n-by-n matrix over a}
{\small\[
\forall a.\; (\operatorname{is\text{-}ring}(a) \Rightarrow \forall n \in \mathbb{N}, r \in \operatorname{carr}(a), k, l \in \operatorname{interval}(1, n).\; (\neg (k = l) \Rightarrow \operatorname{is\text{-}invertible\text{-}mat}(a, n, \operatorname{elem\text{-}g}(a, n, r, k, l))))
\]}
\noindent\textbf{Machine-checked} (trust: proof).

\subsection{Matrix equivalence}

\subsubsection*{\texttt{mat-equiv}}
\emph{for every a, m, n, c, d, c and d are equivalent m-by-n matrices over a if and only if there is some u such that u is an invertible m-by-m matrix over a and there is some v such that v is an invertible n-by-n matrix over a and d equals matmul(a, matmul(a, u, c), v)}
{\small\[
\forall a, m, n, c, d.\; (\operatorname{mat\text{-}equiv}(a, m, n, c, d) \Leftrightarrow \exists u.\; (\operatorname{is\text{-}invertible\text{-}mat}(a, m, u) \wedge \exists v.\; (\operatorname{is\text{-}invertible\text{-}mat}(a, n, v) \wedge (d = \operatorname{matmul}(a, \operatorname{matmul}(a, u, c), v)))))
\]}
\noindent\textbf{Definition.}

\subsubsection*{\texttt{mat-equiv-refl}}
\emph{for every ring a, for every m, n, c in mat(m, n, carr(a)), c and c are equivalent m-by-n matrices over a}
{\small\[
\forall a.\; (\operatorname{is\text{-}ring}(a) \Rightarrow \forall m, n, c \in \operatorname{mat}(m, n, \operatorname{carr}(a)).\; \operatorname{mat\text{-}equiv}(a, m, n, c, c))
\]}
\noindent\textbf{Machine-checked} (trust: proof).

\subsubsection*{\texttt{mat-equiv-trans}}
\emph{for every ring a, for every m, n, c, d, e, if c is in mat(m, n, carr(a)), then if c and d are equivalent m-by-n matrices over a, then if d and e are equivalent m-by-n matrices over a, then c and e are equivalent m-by-n matrices over a}
{\small\[
\forall a.\; (\operatorname{is\text{-}ring}(a) \Rightarrow \forall m, n, c \in \operatorname{mat}(m, n, \operatorname{carr}(a)), d, e.\; (\operatorname{mat\text{-}equiv}(a, m, n, c, d) \Rightarrow (\operatorname{mat\text{-}equiv}(a, m, n, d, e) \Rightarrow \operatorname{mat\text{-}equiv}(a, m, n, c, e))))
\]}
\noindent\textbf{Machine-checked} (trust: proof).

\subsubsection*{\texttt{mat-equiv-left-mult}}
\emph{for every ring a, for every m, n, c, u, if c is in mat(m, n, carr(a)), then if u is an invertible m-by-m matrix over a, then c and matmul(a, u, c) are equivalent m-by-n matrices over a}
{\small\[
\forall a.\; (\operatorname{is\text{-}ring}(a) \Rightarrow \forall m, n, c \in \operatorname{mat}(m, n, \operatorname{carr}(a)), u.\; (\operatorname{is\text{-}invertible\text{-}mat}(a, m, u) \Rightarrow \operatorname{mat\text{-}equiv}(a, m, n, c, \operatorname{matmul}(a, u, c))))
\]}
\noindent\textbf{Machine-checked} (trust: proof).

\subsubsection*{\texttt{mat-equiv-right-mult}}
\emph{for every ring a, for every m, n, c, v, if c is in mat(m, n, carr(a)), then if v is an invertible n-by-n matrix over a, then c and matmul(a, c, v) are equivalent m-by-n matrices over a}
{\small\[
\forall a.\; (\operatorname{is\text{-}ring}(a) \Rightarrow \forall m, n, c \in \operatorname{mat}(m, n, \operatorname{carr}(a)), v.\; (\operatorname{is\text{-}invertible\text{-}mat}(a, n, v) \Rightarrow \operatorname{mat\text{-}equiv}(a, m, n, c, \operatorname{matmul}(a, c, v))))
\]}
\noindent\textbf{Machine-checked} (trust: proof).

\subsubsection*{\texttt{mat-equiv-target-is-mat}}
\emph{for every ring a, for every m, n, c, d, if c is in mat(m, n, carr(a)), then if c and d are equivalent m-by-n matrices over a, then d is in mat(m, n, carr(a))}
{\small\[
\forall a.\; (\operatorname{is\text{-}ring}(a) \Rightarrow \forall m, n, c \in \operatorname{mat}(m, n, \operatorname{carr}(a)), d.\; (\operatorname{mat\text{-}equiv}(a, m, n, c, d) \Rightarrow (d \in \operatorname{mat}(m, n, \operatorname{carr}(a)))))
\]}
\noindent\textbf{Machine-checked} (trust: proof).

\subsection{Matrices acting on sequences of module elements}

\subsubsection*{\texttt{matact-type}}
\emph{for every module md, for every m, n, q, p, u, if p is in mat(m, n, carr(scal(md))), then if u is in mat(n, q, vec(md)), then if if n equals 0, then m equals 0 or q equals 0, then matact(md, p, u) is in mat(m, q, vec(md))}
{\small\[
\forall \mathit{md}.\; (\operatorname{is\text{-}module}(\mathit{md}) \Rightarrow \forall m, n, q, p \in \operatorname{mat}(m, n, \operatorname{carr}(\operatorname{scal}(\mathit{md}))), u \in \operatorname{mat}(n, q, \operatorname{vec}(\mathit{md})).\; (((n = 0) \Rightarrow ((m = 0) \vee (q = 0))) \Rightarrow (\operatorname{matact}(\mathit{md}, p, u) \in \operatorname{mat}(m, q, \operatorname{vec}(\mathit{md})))))
\]}
\noindent\textbf{Machine-checked} (trust: proof).

\subsubsection*{\texttt{matact-entry}}
\emph{for every module md, for every m, n, q, p, u, if p is in mat(m, n, carr(scal(md))), then if u is in mat(n, q, vec(md)), then if 1 is at most n, then for every i in interval(1, m), c in interval(1, q), entry(matact(md, p, u), i, c) equals finsum(module-vector-ag(md), vnb-lambda(j, interval(1, n), (act(md))(entry(p, i, j), entry(u, j, c))), interval(1, n))}
{\small\[
\forall \mathit{md}.\; (\operatorname{is\text{-}module}(\mathit{md}) \Rightarrow \forall m, n, q, p \in \operatorname{mat}(m, n, \operatorname{carr}(\operatorname{scal}(\mathit{md}))), u \in \operatorname{mat}(n, q, \operatorname{vec}(\mathit{md})).\; ((1 \leq n) \Rightarrow \forall i \in \operatorname{interval}(1, m), c \in \operatorname{interval}(1, q).\; (\operatorname{entry}(\operatorname{matact}(\mathit{md}, p, u), i, c) = \operatorname{finsum}(\operatorname{module\text{-}vector\text{-}ag}(\mathit{md}), (\lambda j \in \operatorname{interval}(1, n).\; \operatorname{act}(\mathit{md})(\operatorname{entry}(p, i, j), \operatorname{entry}(u, j, c))), \operatorname{interval}(1, n)))))
\]}
\noindent\textbf{Machine-checked} (trust: proof).

\subsubsection*{\texttt{matact-identmat}}
\emph{for every module md, for every n, q, u in mat(n, q, vec(md)), matact(md, identmat(scal(md), n), u) equals u}
{\small\[
\forall \mathit{md}.\; (\operatorname{is\text{-}module}(\mathit{md}) \Rightarrow \forall n, q, u \in \operatorname{mat}(n, q, \operatorname{vec}(\mathit{md})).\; (\operatorname{matact}(\mathit{md}, \operatorname{identmat}(\operatorname{scal}(\mathit{md}), n), u) = u))
\]}
\noindent\textbf{Machine-checked} (trust: proof).

\subsubsection*{\texttt{matact-assoc}}
\emph{for every module md, for every m, n, k, l, p, q, u, if p is in mat(m, n, carr(scal(md))), then if q is in mat(n, k, carr(scal(md))), then if u is in mat(k, l, vec(md)), then if if n equals 0, then m equals 0 or k equals 0 and l equals 0, then if if k equals 0, then l equals 0 or n equals 0 and m equals 0, then matact(md, matmul(scal(md), p, q), u) equals matact(md, p, matact(md, q, u))}
{\small\[
\forall \mathit{md}.\; (\operatorname{is\text{-}module}(\mathit{md}) \Rightarrow \forall m, n, k, l, p \in \operatorname{mat}(m, n, \operatorname{carr}(\operatorname{scal}(\mathit{md}))), q \in \operatorname{mat}(n, k, \operatorname{carr}(\operatorname{scal}(\mathit{md}))), u \in \operatorname{mat}(k, l, \operatorname{vec}(\mathit{md})).\; (((n = 0) \Rightarrow ((m = 0) \vee ((k = 0) \wedge (l = 0)))) \Rightarrow (((k = 0) \Rightarrow ((l = 0) \vee ((n = 0) \wedge (m = 0)))) \Rightarrow (\operatorname{matact}(\mathit{md}, \operatorname{matmul}(\operatorname{scal}(\mathit{md}), p, q), u) = \operatorname{matact}(\mathit{md}, p, \operatorname{matact}(\mathit{md}, q, u))))))
\]}
\noindent\textbf{Machine-checked} (trust: proof).

\subsubsection*{\texttt{lincomb-row-add}}
\emph{for every module md, for every n, c1, c2, u, if c1 is in mat(1, n, carr(scal(md))), then if c2 is in mat(1, n, carr(scal(md))), then if u is in mat(n, 1, vec(md)), then lincomb(md, n, matadd(scal(md), c1, c2), u) equals (vadd(md))(lincomb(md, n, c1, u), lincomb(md, n, c2, u))}
{\small\[
\forall \mathit{md}.\; (\operatorname{is\text{-}module}(\mathit{md}) \Rightarrow \forall n, c_{1}, c_{2} \in \operatorname{mat}(1, n, \operatorname{carr}(\operatorname{scal}(\mathit{md}))), u \in \operatorname{mat}(n, 1, \operatorname{vec}(\mathit{md})).\; (\operatorname{lincomb}(\mathit{md}, n, \operatorname{matadd}(\operatorname{scal}(\mathit{md}), c_{1}, c_{2}), u) = \operatorname{vadd}(\mathit{md})(\operatorname{lincomb}(\mathit{md}, n, c_{1}, u), \operatorname{lincomb}(\mathit{md}, n, c_{2}, u))))
\]}
\noindent\textbf{Machine-checked} (trust: proof).

\subsubsection*{\texttt{lincomb-row-scale}}
\emph{for every module md, for every n, r, c, u, if r is in carr(scal(md)), then if c is in mat(1, n, carr(scal(md))), then if u is in mat(n, 1, vec(md)), then lincomb(md, n, matscale(scal(md), r, c), u) equals (act(md))(r, lincomb(md, n, c, u))}
{\small\[
\forall \mathit{md}.\; (\operatorname{is\text{-}module}(\mathit{md}) \Rightarrow \forall n, r \in \operatorname{carr}(\operatorname{scal}(\mathit{md})), c \in \operatorname{mat}(1, n, \operatorname{carr}(\operatorname{scal}(\mathit{md}))), u \in \operatorname{mat}(n, 1, \operatorname{vec}(\mathit{md})).\; (\operatorname{lincomb}(\mathit{md}, n, \operatorname{matscale}(\operatorname{scal}(\mathit{md}), r, c), u) = \operatorname{act}(\mathit{md})(r, \operatorname{lincomb}(\mathit{md}, n, c, u))))
\]}
\noindent\textbf{Machine-checked} (trust: proof).

\subsubsection*{\texttt{lincomb-row-peel}}
\emph{for every module md, for every n in nn, c in mat(1, succ(n), carr(scal(md))), u in mat(succ(n), 1, vec(md)), lincomb(md, succ(n), c, u) equals (vadd(md))(lincomb(md, n, block(c, 1, n), block(u, n, 1)), (act(md))(entry(c, 1, succ(n)), entry(u, succ(n), 1)))}
{\small\[
\forall \mathit{md}.\; (\operatorname{is\text{-}module}(\mathit{md}) \Rightarrow \forall n \in \mathbb{N}, c \in \operatorname{mat}(1, (n + 1), \operatorname{carr}(\operatorname{scal}(\mathit{md}))), u \in \operatorname{mat}((n + 1), 1, \operatorname{vec}(\mathit{md})).\; (\operatorname{lincomb}(\mathit{md}, (n + 1), c, u) = \operatorname{vadd}(\mathit{md})(\operatorname{lincomb}(\mathit{md}, n, \operatorname{block}(c, 1, n), \operatorname{block}(u, n, 1)), \operatorname{act}(\mathit{md})(\operatorname{entry}(c, 1, (n + 1)), \operatorname{entry}(u, (n + 1), 1)))))
\]}
\noindent\textbf{Machine-checked} (trust: proof).

\subsubsection*{\texttt{lincomb-snoc}}
\emph{for every module md, for every n in nn, c in mat(1, n, carr(scal(md))), u in mat(n, 1, vec(md)), r in carr(scal(md)), x in vec(md), lincomb(md, succ(n), snoc-row(c, n, r), snoc-col(u, n, x)) equals (vadd(md))(lincomb(md, n, c, u), (act(md))(r, x))}
{\small\[
\forall \mathit{md}.\; (\operatorname{is\text{-}module}(\mathit{md}) \Rightarrow \forall n \in \mathbb{N}, c \in \operatorname{mat}(1, n, \operatorname{carr}(\operatorname{scal}(\mathit{md}))), u \in \operatorname{mat}(n, 1, \operatorname{vec}(\mathit{md})), r \in \operatorname{carr}(\operatorname{scal}(\mathit{md})), x \in \operatorname{vec}(\mathit{md}).\; (\operatorname{lincomb}(\mathit{md}, (n + 1), \operatorname{snoc\text{-}row}(c, n, r), \operatorname{snoc\text{-}col}(u, n, x)) = \operatorname{vadd}(\mathit{md})(\operatorname{lincomb}(\mathit{md}, n, c, u), \operatorname{act}(\mathit{md})(r, x))))
\]}
\noindent\textbf{Machine-checked} (trust: proof).

\subsubsection*{\texttt{lincomb-unitrow}}
\emph{for every module md, for every n, i, u in mat(n, 1, vec(md)), if i is in interval(1, n), then lincomb(md, n, unitrow(scal(md), n, i), u) equals entry(u, i, 1)}
{\small\[
\forall \mathit{md}.\; (\operatorname{is\text{-}module}(\mathit{md}) \Rightarrow \forall n, i \in \operatorname{interval}(1, n), u \in \operatorname{mat}(n, 1, \operatorname{vec}(\mathit{md})).\; (\operatorname{lincomb}(\mathit{md}, n, \operatorname{unitrow}(\operatorname{scal}(\mathit{md}), n, i), u) = \operatorname{entry}(u, i, 1)))
\]}
\noindent\textbf{Machine-checked} (trust: proof).

\subsubsection*{\texttt{lincomb-zerorow}}
\emph{for every module md, for every n, u in mat(n, 1, vec(md)), lincomb(md, n, zeromat(scal(md), 1, n), u) equals vzero(md)}
{\small\[
\forall \mathit{md}.\; (\operatorname{is\text{-}module}(\mathit{md}) \Rightarrow \forall n, u \in \operatorname{mat}(n, 1, \operatorname{vec}(\mathit{md})).\; (\operatorname{lincomb}(\mathit{md}, n, \operatorname{zeromat}(\operatorname{scal}(\mathit{md}), 1, n), u) = \operatorname{vzero}(\mathit{md})))
\]}
\noindent\textbf{Machine-checked} (trust: proof).

\subsubsection*{\texttt{lincomb-empty}}
\emph{for every module md, for every c, u, lincomb(md, 0, c, u) equals vzero(md)}
{\small\[
\forall \mathit{md}.\; (\operatorname{is\text{-}module}(\mathit{md}) \Rightarrow \forall c, u.\; (\operatorname{lincomb}(\mathit{md}, 0, c, u) = \operatorname{vzero}(\mathit{md})))
\]}
\noindent\textbf{Machine-checked} (trust: proof).

\subsubsection*{\texttt{lincomb-unfold}}
\emph{for every md, n, c, u, lincomb(md, n, c, u) is identical to finsum(module-vector-ag(md), vnb-lambda(j, interval(1, n), (act(md))(entry(c, 1, j), entry(u, j, 1))), interval(1, n))}
{\small\[
\forall \mathit{md}, n, c, u.\; (\operatorname{lincomb}(\mathit{md}, n, c, u) \simeq \operatorname{finsum}(\operatorname{module\text{-}vector\text{-}ag}(\mathit{md}), (\lambda j \in \operatorname{interval}(1, n).\; \operatorname{act}(\mathit{md})(\operatorname{entry}(c, 1, j), \operatorname{entry}(u, j, 1))), \operatorname{interval}(1, n)))
\]}
\noindent\textbf{Machine-checked} (trust: proof).

\subsubsection*{\texttt{lincomb-type}}
\emph{for every module md, for every n, c, u, if c is in mat(1, n, carr(scal(md))), then if u is in mat(n, 1, vec(md)), then lincomb(md, n, c, u) is in vec(md)}
{\small\[
\forall \mathit{md}.\; (\operatorname{is\text{-}module}(\mathit{md}) \Rightarrow \forall n, c \in \operatorname{mat}(1, n, \operatorname{carr}(\operatorname{scal}(\mathit{md}))), u \in \operatorname{mat}(n, 1, \operatorname{vec}(\mathit{md})).\; (\operatorname{lincomb}(\mathit{md}, n, c, u) \in \operatorname{vec}(\mathit{md})))
\]}
\noindent\textbf{Machine-checked} (trust: proof).

\subsubsection*{\texttt{matact-lincomb}}
\emph{for every module md, for every n, c, u, if c is in mat(1, n, carr(scal(md))), then if u is in mat(n, 1, vec(md)), then if 1 is at most n, then entry(matact(md, c, u), 1, 1) equals lincomb(md, n, c, u)}
{\small\[
\forall \mathit{md}.\; (\operatorname{is\text{-}module}(\mathit{md}) \Rightarrow \forall n, c \in \operatorname{mat}(1, n, \operatorname{carr}(\operatorname{scal}(\mathit{md}))), u \in \operatorname{mat}(n, 1, \operatorname{vec}(\mathit{md})).\; ((1 \leq n) \Rightarrow (\operatorname{entry}(\operatorname{matact}(\mathit{md}, c, u), 1, 1) = \operatorname{lincomb}(\mathit{md}, n, c, u))))
\]}
\noindent\textbf{Machine-checked} (trust: proof).

\subsubsection*{\texttt{matact-triple-left}}
\emph{for every module md, for every m, n, k, l, p, q, u, row, col, if p is in mat(m, n, carr(scal(md))), then if q is in mat(n, k, carr(scal(md))), then if u is in mat(k, l, vec(md)), then if 1 is at most n, then if 1 is at most k, then if row is in interval(1, m), then if col is in interval(1, l), then entry(matact(md, matmul(scal(md), p, q), u), row, col) equals finsum(module-vector-ag(md), vnb-lambda(c, interval(1, k), finsum(module-vector-ag(md), vnb-lambda(j, interval(1, n), (vnb-lambda(z, cartesian(interval(1, k), interval(1, n)), (act(md))((mul(scal(md)))(entry(p, row, nth(2, z)), entry(q, nth(2, z), nth(1, z))), entry(u, nth(1, z), col))))([c, j])), interval(1, n))), interval(1, k))}
{\small\[
\forall \mathit{md}.\; (\operatorname{is\text{-}module}(\mathit{md}) \Rightarrow \forall m, n, k, l, p \in \operatorname{mat}(m, n, \operatorname{carr}(\operatorname{scal}(\mathit{md}))), q \in \operatorname{mat}(n, k, \operatorname{carr}(\operatorname{scal}(\mathit{md}))), u \in \operatorname{mat}(k, l, \operatorname{vec}(\mathit{md})), \mathit{row}, \mathit{col}.\; ((1 \leq n) \Rightarrow ((1 \leq k) \Rightarrow ((\mathit{row} \in \operatorname{interval}(1, m)) \Rightarrow ((\mathit{col} \in \operatorname{interval}(1, l)) \Rightarrow (\operatorname{entry}(\operatorname{matact}(\mathit{md}, \operatorname{matmul}(\operatorname{scal}(\mathit{md}), p, q), u), \mathit{row}, \mathit{col}) = \operatorname{finsum}(\operatorname{module\text{-}vector\text{-}ag}(\mathit{md}), (\lambda c \in \operatorname{interval}(1, k).\; \operatorname{finsum}(\operatorname{module\text{-}vector\text{-}ag}(\mathit{md}), (\lambda j \in \operatorname{interval}(1, n).\; (\lambda z \in (\operatorname{interval}(1, k) \times \operatorname{interval}(1, n)).\; \operatorname{act}(\mathit{md})(\operatorname{mul}(\operatorname{scal}(\mathit{md}))(\operatorname{entry}(p, \mathit{row}, \operatorname{nth}(2, z)), \operatorname{entry}(q, \operatorname{nth}(2, z), \operatorname{nth}(1, z))), \operatorname{entry}(u, \operatorname{nth}(1, z), \mathit{col})))(\langle c, j \rangle)), \operatorname{interval}(1, n))), \operatorname{interval}(1, k))))))))
\]}
\noindent\textbf{Machine-checked} (trust: proof).

\subsubsection*{\texttt{matact-triple-right}}
\emph{for every module md, for every m, n, k, l, p, q, u, row, col, if p is in mat(m, n, carr(scal(md))), then if q is in mat(n, k, carr(scal(md))), then if u is in mat(k, l, vec(md)), then if 1 is at most n, then if 1 is at most k, then if row is in interval(1, m), then if col is in interval(1, l), then entry(matact(md, p, matact(md, q, u)), row, col) equals finsum(module-vector-ag(md), vnb-lambda(j, interval(1, n), finsum(module-vector-ag(md), vnb-lambda(c, interval(1, k), (vnb-lambda(z, cartesian(interval(1, k), interval(1, n)), (act(md))((mul(scal(md)))(entry(p, row, nth(2, z)), entry(q, nth(2, z), nth(1, z))), entry(u, nth(1, z), col))))([c, j])), interval(1, k))), interval(1, n))}
{\small\[
\forall \mathit{md}.\; (\operatorname{is\text{-}module}(\mathit{md}) \Rightarrow \forall m, n, k, l, p \in \operatorname{mat}(m, n, \operatorname{carr}(\operatorname{scal}(\mathit{md}))), q \in \operatorname{mat}(n, k, \operatorname{carr}(\operatorname{scal}(\mathit{md}))), u \in \operatorname{mat}(k, l, \operatorname{vec}(\mathit{md})), \mathit{row}, \mathit{col}.\; ((1 \leq n) \Rightarrow ((1 \leq k) \Rightarrow ((\mathit{row} \in \operatorname{interval}(1, m)) \Rightarrow ((\mathit{col} \in \operatorname{interval}(1, l)) \Rightarrow (\operatorname{entry}(\operatorname{matact}(\mathit{md}, p, \operatorname{matact}(\mathit{md}, q, u)), \mathit{row}, \mathit{col}) = \operatorname{finsum}(\operatorname{module\text{-}vector\text{-}ag}(\mathit{md}), (\lambda j \in \operatorname{interval}(1, n).\; \operatorname{finsum}(\operatorname{module\text{-}vector\text{-}ag}(\mathit{md}), (\lambda c \in \operatorname{interval}(1, k).\; (\lambda z \in (\operatorname{interval}(1, k) \times \operatorname{interval}(1, n)).\; \operatorname{act}(\mathit{md})(\operatorname{mul}(\operatorname{scal}(\mathit{md}))(\operatorname{entry}(p, \mathit{row}, \operatorname{nth}(2, z)), \operatorname{entry}(q, \operatorname{nth}(2, z), \operatorname{nth}(1, z))), \operatorname{entry}(u, \operatorname{nth}(1, z), \mathit{col})))(\langle c, j \rangle)), \operatorname{interval}(1, k))), \operatorname{interval}(1, n))))))))
\]}
\noindent\textbf{Machine-checked} (trust: proof).

\subsection{Determinants: minors and cofactor expansion}

\subsubsection*{\texttt{minor}}
\emph{Definition / vocabulary.}

\subsubsection*{\texttt{det}}
\emph{Definition / vocabulary.}

\subsubsection*{\texttt{minor-type}}
\emph{for every r, s, p, q, n, if r is a ring, then if p is in nn, then if q is in nn, then if n is in nn, then if s is in mat(succ(n), succ(n), carr(r)), then minor(s, p, q, n) is in mat(n, n, carr(r))}
{\small\[
\forall r, s, p, q, n.\; (\operatorname{is\text{-}ring}(r) \Rightarrow ((p \in \mathbb{N}) \Rightarrow ((q \in \mathbb{N}) \Rightarrow ((n \in \mathbb{N}) \Rightarrow ((s \in \operatorname{mat}((n + 1), (n + 1), \operatorname{carr}(r))) \Rightarrow (\operatorname{minor}(s, p, q, n) \in \operatorname{mat}(n, n, \operatorname{carr}(r))))))))
\]}
\noindent\textbf{Machine-checked} (trust: proof).

\subsubsection*{\texttt{det-zero}}
\emph{for every r, a, det(r, 0, a) is identical to one(r)}
{\small\[
\forall r, a.\; (\operatorname{det}(r, 0, a) \simeq \operatorname{one}(r))
\]}
\noindent\textbf{Definition.}

\subsubsection*{\texttt{det-cofactor}}
\emph{for every r, n, a, if r is a ring, then if n is in nn, then if a is in mat(succ(n), succ(n), carr(r)), then det(r, succ(n), a) is identical to finsum(ring-additive-ag(r), vnb-lambda(j, interval(1, succ(n)), (mul(r))(mpow(ring-multiplicative-monoid(r), (neg(r))(one(r)), succ(j)), (mul(r))(entry(a, 1, j), det(r, n, minor(a, 1, j, n))))), interval(1, succ(n)))}
{\small\[
\forall r, n, a.\; (\operatorname{is\text{-}ring}(r) \Rightarrow ((n \in \mathbb{N}) \Rightarrow ((a \in \operatorname{mat}((n + 1), (n + 1), \operatorname{carr}(r))) \Rightarrow (\operatorname{det}(r, (n + 1), a) \simeq \operatorname{finsum}(\operatorname{ring\text{-}additive\text{-}ag}(r), (\lambda j \in \operatorname{interval}(1, (n + 1)).\; \operatorname{mul}(r)(\operatorname{mpow}(\operatorname{ring\text{-}multiplicative\text{-}monoid}(r), \operatorname{neg}(r)(\operatorname{one}(r)), (j + 1)), \operatorname{mul}(r)(\operatorname{entry}(a, 1, j), \operatorname{det}(r, n, \operatorname{minor}(a, 1, j, n))))), \operatorname{interval}(1, (n + 1)))))))
\]}
\noindent\textbf{Definition.}

\subsubsection*{\texttt{det-in-carrier}}
\emph{for every r, n, a, if r is a ring, then if n is in nn, then if a is in mat(n, n, carr(r)), then det(r, n, a) is in carr(r)}
{\small\[
\forall r, n, a.\; (\operatorname{is\text{-}ring}(r) \Rightarrow ((n \in \mathbb{N}) \Rightarrow ((a \in \operatorname{mat}(n, n, \operatorname{carr}(r))) \Rightarrow (\operatorname{det}(r, n, a) \in \operatorname{carr}(r)))))
\]}
\noindent\textbf{Machine-checked} (trust: proof).

\subsubsection*{\texttt{det-1x1}}
\emph{for every r, a, if r is a ring, then if a is in mat(1, 1, carr(r)), then det(r, 1, a) equals entry(a, 1, 1)}
{\small\[
\forall r, a.\; (\operatorname{is\text{-}ring}(r) \Rightarrow ((a \in \operatorname{mat}(1, 1, \operatorname{carr}(r))) \Rightarrow (\operatorname{det}(r, 1, a) = \operatorname{entry}(a, 1, 1))))
\]}
\noindent\textbf{Machine-checked} (trust: proof).

\subsubsection*{\texttt{det-2x2}}
\emph{for every r, a, if r is a ring, then if a is in mat(2, 2, carr(r)), then det(r, 2, a) equals (add(r))((mul(r))(entry(a, 1, 1), entry(a, 2, 2)), (neg(r))((mul(r))(entry(a, 1, 2), entry(a, 2, 1))))}
{\small\[
\forall r, a.\; (\operatorname{is\text{-}ring}(r) \Rightarrow ((a \in \operatorname{mat}(2, 2, \operatorname{carr}(r))) \Rightarrow (\operatorname{det}(r, 2, a) = \operatorname{add}(r)(\operatorname{mul}(r)(\operatorname{entry}(a, 1, 1), \operatorname{entry}(a, 2, 2)), \operatorname{neg}(r)(\operatorname{mul}(r)(\operatorname{entry}(a, 1, 2), \operatorname{entry}(a, 2, 1)))))))
\]}
\noindent\textbf{Machine-checked} (trust: proof).

\subsubsection*{\texttt{det-identity}}
\emph{for every r, n, if r is a ring, then if n is in nn, then det(r, n, one(mat-ring(r, n))) equals one(r)}
{\small\[
\forall r, n.\; (\operatorname{is\text{-}ring}(r) \Rightarrow ((n \in \mathbb{N}) \Rightarrow (\operatorname{det}(r, n, \operatorname{one}(\operatorname{mat\text{-}ring}(r, n))) = \operatorname{one}(r))))
\]}
\noindent\textbf{Machine-checked} (trust: proof).

\subsubsection*{\texttt{det-row-linear}}
\emph{for every r, n, a, b, c, p, sc, if r is a commutative ring, then if n is in nn, then if a is in mat(n, n, carr(r)), then if b is in mat(n, n, carr(r)), then if c is in mat(n, n, carr(r)), then if p is in interval(1, n), then if sc is in carr(r), then if for every dri\_ in interval(1, n), if dri\_ is not equal to p, then for every drk\_ in interval(1, n), entry(b, dri\_, drk\_) equals entry(a, dri\_, drk\_), then if for every dri\_ in interval(1, n), if dri\_ is not equal to p, then for every drk\_ in interval(1, n), entry(c, dri\_, drk\_) equals entry(a, dri\_, drk\_), then if for every drk\_ in interval(1, n), entry(c, p, drk\_) equals (add(r))((mul(r))(sc, entry(a, p, drk\_)), entry(b, p, drk\_)), then det(r, n, c) equals (add(r))((mul(r))(sc, det(r, n, a)), det(r, n, b))}
{\small\[
\forall r, n, a, b, c, p, \mathit{sc}.\; (\operatorname{is\text{-}commutative\text{-}ring}(r) \Rightarrow ((n \in \mathbb{N}) \Rightarrow ((a \in \operatorname{mat}(n, n, \operatorname{carr}(r))) \Rightarrow ((b \in \operatorname{mat}(n, n, \operatorname{carr}(r))) \Rightarrow ((c \in \operatorname{mat}(n, n, \operatorname{carr}(r))) \Rightarrow ((p \in \operatorname{interval}(1, n)) \Rightarrow ((\mathit{sc} \in \operatorname{carr}(r)) \Rightarrow (\forall \operatorname{dri\_} \in \operatorname{interval}(1, n).\; (\neg (\operatorname{dri\_} = p) \Rightarrow \forall \operatorname{drk\_} \in \operatorname{interval}(1, n).\; (\operatorname{entry}(b, \operatorname{dri\_}, \operatorname{drk\_}) = \operatorname{entry}(a, \operatorname{dri\_}, \operatorname{drk\_}))) \Rightarrow (\forall \operatorname{dri\_} \in \operatorname{interval}(1, n).\; (\neg (\operatorname{dri\_} = p) \Rightarrow \forall \operatorname{drk\_} \in \operatorname{interval}(1, n).\; (\operatorname{entry}(c, \operatorname{dri\_}, \operatorname{drk\_}) = \operatorname{entry}(a, \operatorname{dri\_}, \operatorname{drk\_}))) \Rightarrow (\forall \operatorname{drk\_} \in \operatorname{interval}(1, n).\; (\operatorname{entry}(c, p, \operatorname{drk\_}) = \operatorname{add}(r)(\operatorname{mul}(r)(\mathit{sc}, \operatorname{entry}(a, p, \operatorname{drk\_})), \operatorname{entry}(b, p, \operatorname{drk\_}))) \Rightarrow (\operatorname{det}(r, n, c) = \operatorname{add}(r)(\operatorname{mul}(r)(\mathit{sc}, \operatorname{det}(r, n, a)), \operatorname{det}(r, n, b)))))))))))))
\]}
\noindent\textbf{Machine-checked} (trust: proof).

\subsubsection*{\texttt{det-swap-rows}}
\emph{for every r, n, a, b, p, q, if r is a commutative ring, then if n is in nn, then if a is in mat(n, n, carr(r)), then if b is in mat(n, n, carr(r)), then if p is in interval(1, n), then if q is in interval(1, n), then if p is not equal to q, then if for every dri\_ in interval(1, n), if dri\_ is not equal to p, then if dri\_ is not equal to q, then for every drk\_ in interval(1, n), entry(b, dri\_, drk\_) equals entry(a, dri\_, drk\_), then if for every drk\_ in interval(1, n), entry(b, p, drk\_) equals entry(a, q, drk\_), then if for every drk\_ in interval(1, n), entry(b, q, drk\_) equals entry(a, p, drk\_), then det(r, n, b) equals (neg(r))(det(r, n, a))}
{\small\[
\forall r, n, a, b, p, q.\; (\operatorname{is\text{-}commutative\text{-}ring}(r) \Rightarrow ((n \in \mathbb{N}) \Rightarrow ((a \in \operatorname{mat}(n, n, \operatorname{carr}(r))) \Rightarrow ((b \in \operatorname{mat}(n, n, \operatorname{carr}(r))) \Rightarrow ((p \in \operatorname{interval}(1, n)) \Rightarrow ((q \in \operatorname{interval}(1, n)) \Rightarrow (\neg (p = q) \Rightarrow (\forall \operatorname{dri\_} \in \operatorname{interval}(1, n).\; (\neg (\operatorname{dri\_} = p) \Rightarrow (\neg (\operatorname{dri\_} = q) \Rightarrow \forall \operatorname{drk\_} \in \operatorname{interval}(1, n).\; (\operatorname{entry}(b, \operatorname{dri\_}, \operatorname{drk\_}) = \operatorname{entry}(a, \operatorname{dri\_}, \operatorname{drk\_})))) \Rightarrow (\forall \operatorname{drk\_} \in \operatorname{interval}(1, n).\; (\operatorname{entry}(b, p, \operatorname{drk\_}) = \operatorname{entry}(a, q, \operatorname{drk\_})) \Rightarrow (\forall \operatorname{drk\_} \in \operatorname{interval}(1, n).\; (\operatorname{entry}(b, q, \operatorname{drk\_}) = \operatorname{entry}(a, p, \operatorname{drk\_})) \Rightarrow (\operatorname{det}(r, n, b) = \operatorname{neg}(r)(\operatorname{det}(r, n, a)))))))))))))
\]}
\noindent\textbf{Machine-checked} (trust: proof).

\subsubsection*{\texttt{det-alternating-rows}}
\emph{for every r, n, a, p, q, if r is a commutative ring, then if n is in nn, then if a is in mat(n, n, carr(r)), then if p is in interval(1, n), then if q is in interval(1, n), then if p is not equal to q, then if for every k in interval(1, n), entry(a, p, k) equals entry(a, q, k), then det(r, n, a) equals zero(r)}
{\small\[
\forall r, n, a, p, q.\; (\operatorname{is\text{-}commutative\text{-}ring}(r) \Rightarrow ((n \in \mathbb{N}) \Rightarrow ((a \in \operatorname{mat}(n, n, \operatorname{carr}(r))) \Rightarrow ((p \in \operatorname{interval}(1, n)) \Rightarrow ((q \in \operatorname{interval}(1, n)) \Rightarrow (\neg (p = q) \Rightarrow (\forall k \in \operatorname{interval}(1, n).\; (\operatorname{entry}(a, p, k) = \operatorname{entry}(a, q, k)) \Rightarrow (\operatorname{det}(r, n, a) = \operatorname{zero}(r)))))))))
\]}
\noindent\textbf{Machine-checked} (trust: proof).

\subsubsection*{\texttt{det-expand-row}}
\emph{for every r, n, a, i, if r is a commutative ring, then if n is in nn, then if a is in mat(succ(n), succ(n), carr(r)), then if i is in interval(1, succ(n)), then det(r, succ(n), a) equals finsum(ring-additive-ag(r), vnb-lambda(dxj\_, interval(1, succ(n)), (mul(r))(mpow(ring-multiplicative-monoid(r), (neg(r))(one(r)), i + dxj\_), (mul(r))(entry(a, i, dxj\_), det(r, n, minor(a, i, dxj\_, n))))), interval(1, succ(n)))}
{\small\[
\forall r, n, a, i.\; (\operatorname{is\text{-}commutative\text{-}ring}(r) \Rightarrow ((n \in \mathbb{N}) \Rightarrow ((a \in \operatorname{mat}((n + 1), (n + 1), \operatorname{carr}(r))) \Rightarrow ((i \in \operatorname{interval}(1, (n + 1))) \Rightarrow (\operatorname{det}(r, (n + 1), a) = \operatorname{finsum}(\operatorname{ring\text{-}additive\text{-}ag}(r), (\lambda \operatorname{dxj\_} \in \operatorname{interval}(1, (n + 1)).\; \operatorname{mul}(r)(\operatorname{mpow}(\operatorname{ring\text{-}multiplicative\text{-}monoid}(r), \operatorname{neg}(r)(\operatorname{one}(r)), (i + \operatorname{dxj\_})), \operatorname{mul}(r)(\operatorname{entry}(a, i, \operatorname{dxj\_}), \operatorname{det}(r, n, \operatorname{minor}(a, i, \operatorname{dxj\_}, n))))), \operatorname{interval}(1, (n + 1))))))))
\]}
\noindent\textbf{Machine-checked} (trust: proof).

\subsubsection*{\texttt{det-alternating-form}}
\emph{for every r, n, d, a, if r is a commutative ring, then if n is in nn, then if d is in fun(mat(n, n, carr(r)), carr(r)), then if for every dfa\_, dfb\_, dfc\_, dfp\_, dfs\_, if dfa\_ is in mat(n, n, carr(r)), then if dfb\_ is in mat(n, n, carr(r)), then if dfc\_ is in mat(n, n, carr(r)), then if dfp\_ is in interval(1, n), then if dfs\_ is in carr(r), then if for every dri\_ in interval(1, n), if dri\_ is not equal to dfp\_, then for every drk\_ in interval(1, n), entry(dfb\_, dri\_, drk\_) equals entry(dfa\_, dri\_, drk\_), then if for every dri\_ in interval(1, n), if dri\_ is not equal to dfp\_, then for every drk\_ in interval(1, n), entry(dfc\_, dri\_, drk\_) equals entry(dfa\_, dri\_, drk\_), then if for every drk\_ in interval(1, n), entry(dfc\_, dfp\_, drk\_) equals (add(r))((mul(r))(dfs\_, entry(dfa\_, dfp\_, drk\_)), entry(dfb\_, dfp\_, drk\_)), then d(dfc\_) equals (add(r))((mul(r))(dfs\_, d(dfa\_)), d(dfb\_)), then if for every dfa\_, dfp\_, dfq\_, if dfa\_ is in mat(n, n, carr(r)), then if dfp\_ is in interval(1, n), then if dfq\_ is in interval(1, n), then if dfp\_ is not equal to dfq\_, then if for every drk\_ in interval(1, n), entry(dfa\_, dfp\_, drk\_) equals entry(dfa\_, dfq\_, drk\_), then d(dfa\_) equals zero(r), then if a is in mat(n, n, carr(r)), then d(a) equals (mul(r))(det(r, n, a), d(identmat(r, n)))}
{\small\[
\forall r, n, d, a.\; (\operatorname{is\text{-}commutative\text{-}ring}(r) \Rightarrow ((n \in \mathbb{N}) \Rightarrow ((d \in (\operatorname{mat}(n, n, \operatorname{carr}(r)) \to \operatorname{carr}(r))) \Rightarrow (\forall \operatorname{dfa\_}, \operatorname{dfb\_}, \operatorname{dfc\_} \in \operatorname{mat}(n, n, \operatorname{carr}(r)), \operatorname{dfp\_} \in \operatorname{interval}(1, n), \operatorname{dfs\_} \in \operatorname{carr}(r).\; (\forall \operatorname{dri\_} \in \operatorname{interval}(1, n).\; (\neg (\operatorname{dri\_} = \operatorname{dfp\_}) \Rightarrow \forall \operatorname{drk\_} \in \operatorname{interval}(1, n).\; (\operatorname{entry}(\operatorname{dfb\_}, \operatorname{dri\_}, \operatorname{drk\_}) = \operatorname{entry}(\operatorname{dfa\_}, \operatorname{dri\_}, \operatorname{drk\_}))) \Rightarrow (\forall \operatorname{dri\_} \in \operatorname{interval}(1, n).\; (\neg (\operatorname{dri\_} = \operatorname{dfp\_}) \Rightarrow \forall \operatorname{drk\_} \in \operatorname{interval}(1, n).\; (\operatorname{entry}(\operatorname{dfc\_}, \operatorname{dri\_}, \operatorname{drk\_}) = \operatorname{entry}(\operatorname{dfa\_}, \operatorname{dri\_}, \operatorname{drk\_}))) \Rightarrow (\forall \operatorname{drk\_} \in \operatorname{interval}(1, n).\; (\operatorname{entry}(\operatorname{dfc\_}, \operatorname{dfp\_}, \operatorname{drk\_}) = \operatorname{add}(r)(\operatorname{mul}(r)(\operatorname{dfs\_}, \operatorname{entry}(\operatorname{dfa\_}, \operatorname{dfp\_}, \operatorname{drk\_})), \operatorname{entry}(\operatorname{dfb\_}, \operatorname{dfp\_}, \operatorname{drk\_}))) \Rightarrow (d(\operatorname{dfc\_}) = \operatorname{add}(r)(\operatorname{mul}(r)(\operatorname{dfs\_}, d(\operatorname{dfa\_})), d(\operatorname{dfb\_})))))) \Rightarrow (\forall \operatorname{dfa\_} \in \operatorname{mat}(n, n, \operatorname{carr}(r)), \operatorname{dfp\_}, \operatorname{dfq\_} \in \operatorname{interval}(1, n).\; (\neg (\operatorname{dfp\_} = \operatorname{dfq\_}) \Rightarrow (\forall \operatorname{drk\_} \in \operatorname{interval}(1, n).\; (\operatorname{entry}(\operatorname{dfa\_}, \operatorname{dfp\_}, \operatorname{drk\_}) = \operatorname{entry}(\operatorname{dfa\_}, \operatorname{dfq\_}, \operatorname{drk\_})) \Rightarrow (d(\operatorname{dfa\_}) = \operatorname{zero}(r)))) \Rightarrow ((a \in \operatorname{mat}(n, n, \operatorname{carr}(r))) \Rightarrow (d(a) = \operatorname{mul}(r)(\operatorname{det}(r, n, a), d(\operatorname{identmat}(r, n))))))))))
\]}
\noindent\textbf{Machine-checked} (trust: proof).

\subsubsection*{\texttt{det-multiplicative}}
\emph{for every r, n, p, q, if r is a commutative ring, then if n is in nn, then if p is in mat(n, n, carr(r)), then if q is in mat(n, n, carr(r)), then det(r, n, matmul(r, p, q)) equals (mul(r))(det(r, n, p), det(r, n, q))}
{\small\[
\forall r, n, p, q.\; (\operatorname{is\text{-}commutative\text{-}ring}(r) \Rightarrow ((n \in \mathbb{N}) \Rightarrow ((p \in \operatorname{mat}(n, n, \operatorname{carr}(r))) \Rightarrow ((q \in \operatorname{mat}(n, n, \operatorname{carr}(r))) \Rightarrow (\operatorname{det}(r, n, \operatorname{matmul}(r, p, q)) = \operatorname{mul}(r)(\operatorname{det}(r, n, p), \operatorname{det}(r, n, q)))))))
\]}
\noindent\textbf{Machine-checked} (trust: proof).

\subsubsection*{\texttt{det-mul-invertible}}
\emph{for every r, n, a, v, if r is a field, then if n is in nn, then if a is in mat(n, n, carr(r)), then if v is an invertible n-by-n matrix over r, then det(r, n, matmul(r, a, v)) equals (mul(r))(det(r, n, a), det(r, n, v))}
{\small\[
\forall r, n, a, v.\; (\operatorname{is\text{-}field\text{-}ring}(r) \Rightarrow ((n \in \mathbb{N}) \Rightarrow ((a \in \operatorname{mat}(n, n, \operatorname{carr}(r))) \Rightarrow (\operatorname{is\text{-}invertible\text{-}mat}(r, n, v) \Rightarrow (\operatorname{det}(r, n, \operatorname{matmul}(r, a, v)) = \operatorname{mul}(r)(\operatorname{det}(r, n, a), \operatorname{det}(r, n, v)))))))
\]}
\noindent\textbf{Machine-checked} (trust: proof).

\subsection{The binomial theorem}

\subsubsection*{\texttt{sum-expansion}}
\emph{for every n in nn, r, if r is a commutative ring, then for every x in carr(r), y in carr(r), g in fun(zz, carr(r)), sum(r, vnb-lambda(k, zz, (add(r))((mul(r))(x, g(k - 1)), (mul(r))(y, g(k)))), succ(n)) equals (add(r))((add(r))((mul(r))((add(r))(x, y), sum(r, g, n)), (mul(r))(x, g(0 - 1))), (mul(r))(y, g(n)))}
{\small\[
\forall n \in \mathbb{N}, r.\; (\operatorname{is\text{-}commutative\text{-}ring}(r) \Rightarrow \forall x, y \in \operatorname{carr}(r), g \in (\mathbb{Z} \to \operatorname{carr}(r)).\; (\operatorname{sum}(r, (\lambda k \in \mathbb{Z}.\; \operatorname{add}(r)(\operatorname{mul}(r)(x, g((k - 1))), \operatorname{mul}(r)(y, g(k)))), (n + 1)) = \operatorname{add}(r)(\operatorname{add}(r)(\operatorname{mul}(r)(\operatorname{add}(r)(x, y), \operatorname{sum}(r, g, n)), \operatorname{mul}(r)(x, g((0 - 1)))), \operatorname{mul}(r)(y, g(n)))))
\]}
\noindent\textbf{Machine-checked} (trust: proof).

\subsubsection*{\texttt{binomial-theorem}}
\emph{for every n in nn, r, if r is a commutative ring, then for every x in carr(r), y in carr(r), ring-power(r, (add(r))(x, y), n) equals sum(r, comb-kk(r, x, y, n), succ(n))}
{\small\[
\forall n \in \mathbb{N}, r.\; (\operatorname{is\text{-}commutative\text{-}ring}(r) \Rightarrow \forall x, y \in \operatorname{carr}(r).\; (\operatorname{ring\text{-}power}(r, \operatorname{add}(r)(x, y), n) = \operatorname{sum}(r, \operatorname{comb\text{-}kk}(r, x, y, n), (n + 1))))
\]}
\noindent\textbf{Machine-checked} (trust: proof). Induction on n via a multiply-and-shift sum expansion and Pascal's recurrence on the binomial coefficients, in integer-range SUM form (no bijections).

\section{Combinatorial}

\subsection{Vocabulary}

\subsubsection*{\texttt{inf-subsets}}
\emph{Definition / vocabulary.}

\subsubsection*{\texttt{is-finite-cover}}
\emph{for every c, a, c is a finite cover of a if and only if c is in set and card(c) is in nn and a is a subset of big-union(u, c, u)}
{\small\[
\forall c, a.\; (\operatorname{is\text{-}finite\text{-}cover}(c, a) \Leftrightarrow ((c \in \mathbf{Set}) \wedge ((|c| \in \mathbb{N}) \wedge (a \subseteq \operatorname{big\text{-}union}(u, c, u)))))
\]}
\noindent\textbf{Definition.}

\subsection{Dependent choice / recursion on NN}

\subsubsection*{\texttt{dc-on-nn}}
\emph{for every x in set, a in x, r in set, if for every k in nn, u in x, there is some y in x such that [k, u, y] is in r, then there is some f in fun(nn, x) such that f(0) equals a and for every k in nn, [k, f(k), f(succ(k))] is in r}
{\small\[
\forall x \in \mathbf{Set}, a \in x, r \in \mathbf{Set}.\; (\forall k \in \mathbb{N}, u \in x.\; \exists y \in x.\; (\langle k, u, y \rangle \in r) \Rightarrow \exists f \in (\mathbb{N} \to x).\; ((f(0) = a) \wedge \forall k \in \mathbb{N}.\; (\langle k, f(k), f((k + 1)) \rangle \in r)))
\]}
\noindent\textbf{Machine-checked} (trust: proof).

\subsubsection*{\texttt{dc-on-nn-pred}}
\emph{for every x in set, a in x, nxt, if for every k in nn, u in x, there is some y in x such that y is in nxt(k, u), then there is some f in fun(nn, x) such that f(0) equals a and for every k in nn, f(succ(k)) is in nxt(k, f(k))}
{\small\[
\forall x \in \mathbf{Set}, a \in x, \mathit{nxt}.\; (\forall k \in \mathbb{N}, u \in x.\; \exists y \in x.\; (y \in \operatorname{nxt}(k, u)) \Rightarrow \exists f \in (\mathbb{N} \to x).\; ((f(0) = a) \wedge \forall k \in \mathbb{N}.\; (f((k + 1)) \in \operatorname{nxt}(k, f(k)))))
\]}
\noindent\textbf{Machine-checked} (trust: proof).

\subsection{Infinite pigeonhole and the block family}

\subsubsection*{\texttt{pigeonhole-infinite}}
\emph{for every s, if s is in set and card(s) is not in nn, then for every f, if f is in set and card(f) is in nn, then for every pmap\_ in fun(s, f), there is some c in f such that card(\{x in s: pmap\_(x) = c\}) is not in nn}
{\small\[
\forall s.\; (((s \in \mathbf{Set}) \wedge \neg (|s| \in \mathbb{N})) \Rightarrow \forall f.\; (((f \in \mathbf{Set}) \wedge (|f| \in \mathbb{N})) \Rightarrow \forall \operatorname{pmap\_} \in (s \to f).\; \exists c \in f.\; \neg (|\{x \in s : (\operatorname{pmap\_}(x) = c)\}| \in \mathbb{N})))
\]}
\noindent\textbf{Machine-checked} (trust: proof).

\subsubsection*{\texttt{cover-block-step}}
\emph{for every v in set, f in fun(nn, v), c, if c is a finite cover of v, then for every j in inf-subsets(nn), there is some j\_ in inf-subsets(nn) such that j\_ is a subset of j and there is some u in c such that for every i in j\_, f(i) is in u}
{\small\[
\forall v \in \mathbf{Set}, f \in (\mathbb{N} \to v), c.\; (\operatorname{is\text{-}finite\text{-}cover}(c, v) \Rightarrow \forall j \in \mathrm{Inf}(\mathbb{N}).\; \exists \operatorname{j\_} \in \mathrm{Inf}(\mathbb{N}).\; ((\operatorname{j\_} \subseteq j) \wedge \exists u \in c.\; \forall i \in \operatorname{j\_}.\; (f(i) \in u)))
\]}
\noindent\textbf{Machine-checked} (trust: proof).

\subsubsection*{\texttt{block-family-combinatorial}}
\emph{for every v in set, f in fun(nn, v), cov, if for every k in nn, cov(k) is a finite cover of v, then there is some blk in fun(nn, inf-subsets(nn)) such that for every k in nn, blk(succ(k)) is a subset of blk(k) and for every k in nn, there is some u in cov(k) such that for every i in blk(k), f(i) is in u}
{\small\[
\forall v \in \mathbf{Set}, f \in (\mathbb{N} \to v), \mathit{cov}.\; (\forall k \in \mathbb{N}.\; \operatorname{is\text{-}finite\text{-}cover}(\operatorname{cov}(k), v) \Rightarrow \exists \mathit{blk} \in (\mathbb{N} \to \mathrm{Inf}(\mathbb{N})).\; (\forall k \in \mathbb{N}.\; (\operatorname{blk}((k + 1)) \subseteq \operatorname{blk}(k)) \wedge \forall k \in \mathbb{N}.\; \exists u \in \operatorname{cov}(k).\; \forall i \in \operatorname{blk}(k).\; (f(i) \in u)))
\]}
\noindent\textbf{Machine-checked} (trust: proof). Dependent choice on NN whose step is one infinite-pigeonhole refinement (cover-block-step): each block is an infinite monochromatic sub-block of the previous one for the next cover.  The index shift blk(k)=aux(succ k) aligns capture with level.

\subsection{Diagonalization}

\subsubsection*{\texttt{diagonalization}}
\emph{for every s in fun(nn, inf-subsets(nn)), if for every k in nn, s(succ(k)) is a subset of s(k), then there is some f such that f is a strictly increasing sequence of naturals and for every k in nn, j in nn, if k is at most j, then f(j) is in s(k)}
{\small\[
\forall s \in (\mathbb{N} \to \mathrm{Inf}(\mathbb{N})).\; (\forall k \in \mathbb{N}.\; (s((k + 1)) \subseteq s(k)) \Rightarrow \exists f.\; (\operatorname{strictly\text{-}mono\text{-}nn}(f) \wedge \forall k, j \in \mathbb{N}.\; ((k \leq j) \Rightarrow (f(j) \in s(k)))))
\]}
\noindent\textbf{Machine-checked} (trust: proof). Dependent choice on NN (dc-on-nn-pred) builds a strictly increasing sequence, each term hopping into the next nested set; the nested chain then puts the whole tail inside every member.

\subsubsection*{\texttt{nn-step-strictly-mono}}
\emph{for every g in fun(nn, nn), if for every k in nn, g(k) is less than g(succ(k)), then g is a strictly increasing sequence of naturals}
{\small\[
\forall g \in (\mathbb{N} \to \mathbb{N}).\; (\forall k \in \mathbb{N}.\; (g(k) < g((k + 1))) \Rightarrow \operatorname{strictly\text{-}mono\text{-}nn}(g))
\]}
\noindent\textbf{Machine-checked} (trust: proof).

\subsubsection*{\texttt{nn-nested-subset-chain}}
\emph{for every t, if for every k in nn, t(succ(k)) is a subset of t(k), then for every k in nn, j in nn, if k is at most j, then t(j) is a subset of t(k)}
{\small\[
\forall t.\; (\forall k \in \mathbb{N}.\; (t((k + 1)) \subseteq t(k)) \Rightarrow \forall k, j \in \mathbb{N}.\; ((k \leq j) \Rightarrow (t(j) \subseteq t(k))))
\]}
\noindent\textbf{Machine-checked} (trust: proof).

\subsubsection*{\texttt{inf-subset-nn-unbounded}}
\emph{for every t in inf-subsets(nn), u in nn, there is some y in nn such that y is in t and u is less than y}
{\small\[
\forall t \in \mathrm{Inf}(\mathbb{N}), u \in \mathbb{N}.\; \exists y \in \mathbb{N}.\; ((y \in t) \wedge (u < y))
\]}
\noindent\textbf{Machine-checked} (trust: proof).

\subsection{Well-ordering of NN}

\subsubsection*{\texttt{well-ordering-principle}}
\emph{for every s in set, there is some phi such that phi is in bijection(ord-segment(card(s)), s)}
{\small\[
\forall s \in \mathbf{Set}.\; \exists \varphi.\; (\varphi \in \mathrm{Bij}(\mathrm{OS}(|s|), s))
\]}
\noindent\textbf{Machine-checked} (trust: proof).

\subsubsection*{\texttt{nn-least-element}}
\emph{for every t, if t is a subset of nn and there is some n such that n is in t, then there is some m in t such that for every k in t, m is at most k}
{\small\[
\forall t.\; (((t \subseteq \mathbb{N}) \wedge \exists n.\; (n \in t)) \Rightarrow \exists m \in t.\; \forall k \in t.\; (m \leq k))
\]}
\noindent\textbf{Machine-checked} (trust: proof).

\section{Metric spaces}

\subsection{Vocabulary}

\subsubsection*{\texttt{is-ms-sequence}}
\emph{for every ms, ms is a sequence of metric spaces if and only if for every n in nn, ms(n) is a metric space}
{\small\[
\forall \mathit{ms}.\; (\operatorname{is\text{-}ms\text{-}sequence}(\mathit{ms}) \Leftrightarrow \forall n \in \mathbb{N}.\; \operatorname{is\text{-}metric\text{-}space}(\operatorname{ms}(n)))
\]}
\noindent\textbf{Definition.}

\subsubsection*{\texttt{converges-to}}
\emph{for every s, f, l, f converges to l in s if and only if s is a metric space and f is in fun(nn, pts(s)) and l is in pts(s) and for every positive real eps, there is some n in nn such that for every n\_ in nn, if n is at most n\_, then (dist(s))(f(n\_), l) is at most eps}
{\small\[
\forall s, f, l.\; (\operatorname{converges\text{-}to}(s, f, l) \Leftrightarrow (\operatorname{is\text{-}metric\text{-}space}(s) \wedge ((f \in (\mathbb{N} \to \operatorname{pts}(s))) \wedge ((l \in \operatorname{pts}(s)) \wedge \forall \mathit{eps}.\; (\operatorname{pos\text{-}rr}(\mathit{eps}) \Rightarrow \exists n \in \mathbb{N}.\; \forall \operatorname{n\_} \in \mathbb{N}.\; ((n \leq \operatorname{n\_}) \Rightarrow (\operatorname{dist}(s)(f(\operatorname{n\_}), l) \leq \mathit{eps})))))))
\]}
\noindent\textbf{Definition.}

\subsubsection*{\texttt{converges}}
\emph{for every s, f, f converges in s if and only if there is some l such that f converges to l in s}
{\small\[
\forall s, f.\; (\operatorname{converges}(s, f) \Leftrightarrow \exists l.\; \operatorname{converges\text{-}to}(s, f, l))
\]}
\noindent\textbf{Definition.}

\subsubsection*{\texttt{is-cauchy-seq}}
\emph{for every s, f, is-cauchy-seq(s, f) if and only if s is a metric space and f is in fun(nn, pts(s)) and for every positive real eps, there is some n in nn such that for every m in nn, n\_ in nn, if n is at most m and n is at most n\_, then (dist(s))(f(m), f(n\_)) is at most eps}
{\small\[
\forall s, f.\; (\operatorname{is\text{-}cauchy\text{-}seq}(s, f) \Leftrightarrow (\operatorname{is\text{-}metric\text{-}space}(s) \wedge ((f \in (\mathbb{N} \to \operatorname{pts}(s))) \wedge \forall \mathit{eps}.\; (\operatorname{pos\text{-}rr}(\mathit{eps}) \Rightarrow \exists n \in \mathbb{N}.\; \forall m, \operatorname{n\_} \in \mathbb{N}.\; (((n \leq m) \wedge (n \leq \operatorname{n\_})) \Rightarrow (\operatorname{dist}(s)(f(m), f(\operatorname{n\_})) \leq \mathit{eps}))))))
\]}
\noindent\textbf{Definition.}

\subsubsection*{\texttt{is-complete}}
\emph{for every s, s is complete if and only if s is a metric space and for every f, if is-cauchy-seq(s, f), then f converges in s}
{\small\[
\forall s.\; (\operatorname{is\text{-}complete}(s) \Leftrightarrow (\operatorname{is\text{-}metric\text{-}space}(s) \wedge \forall f.\; (\operatorname{is\text{-}cauchy\text{-}seq}(s, f) \Rightarrow \operatorname{converges}(s, f))))
\]}
\noindent\textbf{Definition.}

\subsubsection*{\texttt{ball}}
\emph{Definition / vocabulary.}

\subsubsection*{\texttt{is-open}}
\emph{for every s, u, u is open in s if and only if s is a metric space and u is a subset of pts(s) and for every y in u, there is some positive real r such that ball(s, y, r) is a subset of u}
{\small\[
\forall s, u.\; (\operatorname{is\text{-}open}(s, u) \Leftrightarrow (\operatorname{is\text{-}metric\text{-}space}(s) \wedge ((u \subseteq \operatorname{pts}(s)) \wedge \forall y \in u.\; \exists r.\; (\operatorname{pos\text{-}rr}(r) \wedge (\operatorname{ball}(s, y, r) \subseteq u)))))
\]}
\noindent\textbf{Definition.}

\subsubsection*{\texttt{is-closed}}
\emph{for every s, a, a is closed in s if and only if s is a metric space and a is a subset of pts(s) and complement-in(pts(s), a) is open in s}
{\small\[
\forall s, a.\; (\operatorname{is\text{-}closed}(s, a) \Leftrightarrow (\operatorname{is\text{-}metric\text{-}space}(s) \wedge ((a \subseteq \operatorname{pts}(s)) \wedge \operatorname{is\text{-}open}(s, \operatorname{complement\text{-}in}(\operatorname{pts}(s), a)))))
\]}
\noindent\textbf{Definition.}

\subsubsection*{\texttt{is-continuous-at}}
\emph{for every s, t, f, a, f is continuous at a if and only if s is a metric space and t is a metric space and f is in fun(pts(s), pts(t)) and a is in pts(s) and for every positive real eps, there is some positive real delta such that for every b in pts(s), if (dist(s))(a, b) is at most delta, then (dist(t))(f(a), f(b)) is at most eps}
{\small\[
\forall s, t, f, a.\; (\operatorname{is\text{-}continuous\text{-}at}(s, t, f, a) \Leftrightarrow (\operatorname{is\text{-}metric\text{-}space}(s) \wedge (\operatorname{is\text{-}metric\text{-}space}(t) \wedge ((f \in (\operatorname{pts}(s) \to \operatorname{pts}(t))) \wedge ((a \in \operatorname{pts}(s)) \wedge \forall \mathit{eps}.\; (\operatorname{pos\text{-}rr}(\mathit{eps}) \Rightarrow \exists \delta.\; (\operatorname{pos\text{-}rr}(\delta) \wedge \forall b \in \operatorname{pts}(s).\; ((\operatorname{dist}(s)(a, b) \leq \delta) \Rightarrow (\operatorname{dist}(t)(f(a), f(b)) \leq \mathit{eps})))))))))
\]}
\noindent\textbf{Definition.}

\subsubsection*{\texttt{is-continuous}}
\emph{for every s, t, f, f is continuous from s to t if and only if s is a metric space and t is a metric space and f is in fun(pts(s), pts(t)) and for every a in pts(s), f is continuous at a}
{\small\[
\forall s, t, f.\; (\operatorname{is\text{-}continuous}(s, t, f) \Leftrightarrow (\operatorname{is\text{-}metric\text{-}space}(s) \wedge (\operatorname{is\text{-}metric\text{-}space}(t) \wedge ((f \in (\operatorname{pts}(s) \to \operatorname{pts}(t))) \wedge \forall a \in \operatorname{pts}(s).\; \operatorname{is\text{-}continuous\text{-}at}(s, t, f, a)))))
\]}
\noindent\textbf{Definition.}

\subsubsection*{\texttt{is-uniformly-continuous}}
\emph{for every s, t, f, f is uniformly continuous from s to t if and only if s is a metric space and t is a metric space and f is in fun(pts(s), pts(t)) and for every positive real eps, there is some positive real delta such that for every a in pts(s), b in pts(s), if (dist(s))(a, b) is at most delta, then (dist(t))(f(a), f(b)) is at most eps}
{\small\[
\forall s, t, f.\; (\operatorname{is\text{-}uniformly\text{-}continuous}(s, t, f) \Leftrightarrow (\operatorname{is\text{-}metric\text{-}space}(s) \wedge (\operatorname{is\text{-}metric\text{-}space}(t) \wedge ((f \in (\operatorname{pts}(s) \to \operatorname{pts}(t))) \wedge \forall \mathit{eps}.\; (\operatorname{pos\text{-}rr}(\mathit{eps}) \Rightarrow \exists \delta.\; (\operatorname{pos\text{-}rr}(\delta) \wedge \forall a, b \in \operatorname{pts}(s).\; ((\operatorname{dist}(s)(a, b) \leq \delta) \Rightarrow (\operatorname{dist}(t)(f(a), f(b)) \leq \mathit{eps}))))))))
\]}
\noindent\textbf{Definition.}

\subsubsection*{\texttt{totally-bounded}}
\emph{for every s, s is totally bounded if and only if s is a metric space and for every r, if r is in rr and 0 is at most r and 0 is not equal to r, then there is some f such that card(f) is in nn and f is an r-net for pts(s) in s}
{\small\[
\forall s.\; (\operatorname{totally\text{-}bounded}(s) \Leftrightarrow (\operatorname{is\text{-}metric\text{-}space}(s) \wedge \forall r.\; (((r \in \mathbb{R}) \wedge ((0 \leq r) \wedge \neg (0 = r))) \Rightarrow \exists f.\; ((|f| \in \mathbb{N}) \wedge \operatorname{is\text{-}r\text{-}net}(s, f, \operatorname{pts}(s), r)))))
\]}
\noindent\textbf{Definition.}

\subsubsection*{\texttt{is-r-net}}
\emph{for every s, f, a, r, f is an r-net for a in s if and only if f is a subset of a and for every p in a, there is some c in f such that (dist(s))(c, p) is at most r and (dist(s))(c, p) is not equal to r}
{\small\[
\forall s, f, a, r.\; (\operatorname{is\text{-}r\text{-}net}(s, f, a, r) \Leftrightarrow ((f \subseteq a) \wedge \forall p \in a.\; \exists c \in f.\; ((\operatorname{dist}(s)(c, p) \leq r) \wedge \neg (\operatorname{dist}(s)(c, p) = r))))
\]}
\noindent\textbf{Definition.}

\subsubsection*{\texttt{product-metric}}
\emph{Definition / vocabulary.}

\subsubsection*{\texttt{completion}}
\emph{Definition / vocabulary.}

\subsection{Sequences, limits, completeness}

\subsubsection*{\texttt{complete-cauchy-converges}}
\emph{for every complete s, for every f, if is-cauchy-seq(s, f), then f converges in s}
{\small\[
\forall s.\; (\operatorname{is\text{-}complete}(s) \Rightarrow \forall f.\; (\operatorname{is\text{-}cauchy\text{-}seq}(s, f) \Rightarrow \operatorname{converges}(s, f)))
\]}
\noindent\textbf{Definition.}

\subsection{Open / closed sets, balls}

\subsubsection*{\texttt{ball-is-open}}
\emph{for every metric space s, for every x in pts(s), r, if r is in rr and 0 is at most r and 0 is not equal to r, then ball(s, x, r) is open in s}
{\small\[
\forall s.\; (\operatorname{is\text{-}metric\text{-}space}(s) \Rightarrow \forall x \in \operatorname{pts}(s), r.\; (((r \in \mathbb{R}) \wedge ((0 \leq r) \wedge \neg (0 = r))) \Rightarrow \operatorname{is\text{-}open}(s, \operatorname{ball}(s, x, r))))
\]}
\noindent\textbf{Machine-checked} (trust: proof).

\subsubsection*{\texttt{ball-membership}}
\emph{for every s, x, r, y, y is in ball(s, x, r) if and only if y is in pts(s) and (dist(s))(x, y) is at most r and (dist(s))(x, y) is not equal to r}
{\small\[
\forall s, x, r, y.\; ((y \in \operatorname{ball}(s, x, r)) \Leftrightarrow ((y \in \operatorname{pts}(s)) \wedge ((\operatorname{dist}(s)(x, y) \leq r) \wedge \neg (\operatorname{dist}(s)(x, y) = r))))
\]}
\noindent\textbf{Machine-checked} (trust: proof).

\subsubsection*{\texttt{ball-center-in}}
\emph{for every metric space s, for every x in pts(s), r, if r is in rr and 0 is at most r and 0 is not equal to r, then x is in ball(s, x, r)}
{\small\[
\forall s.\; (\operatorname{is\text{-}metric\text{-}space}(s) \Rightarrow \forall x \in \operatorname{pts}(s), r.\; (((r \in \mathbb{R}) \wedge ((0 \leq r) \wedge \neg (0 = r))) \Rightarrow (x \in \operatorname{ball}(s, x, r))))
\]}
\noindent\textbf{Machine-checked} (trust: proof).

\subsubsection*{\texttt{empty-is-open}}
\emph{for every metric space s, empty-set is open in s}
{\small\[
\forall s.\; (\operatorname{is\text{-}metric\text{-}space}(s) \Rightarrow \operatorname{is\text{-}open}(s, \emptyset))
\]}
\noindent\textbf{Machine-checked} (trust: proof).

\subsubsection*{\texttt{carrier-is-open}}
\emph{for every metric space s, pts(s) is open in s}
{\small\[
\forall s.\; (\operatorname{is\text{-}metric\text{-}space}(s) \Rightarrow \operatorname{is\text{-}open}(s, \operatorname{pts}(s)))
\]}
\noindent\textbf{Machine-checked} (trust: proof).

\subsubsection*{\texttt{inter-of-opens-open}}
\emph{for every metric space s, for every u, w, if u is open in s and w is open in s, then intersection(u, w) is open in s}
{\small\[
\forall s.\; (\operatorname{is\text{-}metric\text{-}space}(s) \Rightarrow \forall u, w.\; ((\operatorname{is\text{-}open}(s, u) \wedge \operatorname{is\text{-}open}(s, w)) \Rightarrow \operatorname{is\text{-}open}(s, (u \cap w))))
\]}
\noindent\textbf{Machine-checked} (trust: proof).

\subsection{Continuity}

\subsubsection*{\texttt{continuous-is-continuous-at}}
\emph{for every s, t, f, if f is continuous from s to t, then for every a in pts(s), f is continuous at a}
{\small\[
\forall s, t, f.\; (\operatorname{is\text{-}continuous}(s, t, f) \Rightarrow \forall a \in \operatorname{pts}(s).\; \operatorname{is\text{-}continuous\text{-}at}(s, t, f, a))
\]}
\noindent\textbf{Machine-checked} (trust: proof).

\subsubsection*{\texttt{continuous-implies-open-preimage}}
\emph{for every s, t, f, if f is continuous from s to t, then for every v, if v is open in t, then preimage(s, f, v) is open in s}
{\small\[
\forall s, t, f.\; (\operatorname{is\text{-}continuous}(s, t, f) \Rightarrow \forall v.\; (\operatorname{is\text{-}open}(t, v) \Rightarrow \operatorname{is\text{-}open}(s, \operatorname{preimage}(s, f, v))))
\]}
\noindent\textbf{Machine-checked} (trust: proof).

\subsubsection*{\texttt{open-preimage-implies-continuous}}
\emph{for every s, t, f, if s is a metric space and t is a metric space and f is in fun(pts(s), pts(t)), then if for every v, if v is open in t, then preimage(s, f, v) is open in s, then f is continuous from s to t}
{\small\[
\forall s, t, f.\; ((\operatorname{is\text{-}metric\text{-}space}(s) \wedge (\operatorname{is\text{-}metric\text{-}space}(t) \wedge (f \in (\operatorname{pts}(s) \to \operatorname{pts}(t))))) \Rightarrow (\forall v.\; (\operatorname{is\text{-}open}(t, v) \Rightarrow \operatorname{is\text{-}open}(s, \operatorname{preimage}(s, f, v))) \Rightarrow \operatorname{is\text{-}continuous}(s, t, f)))
\]}
\noindent\textbf{Machine-checked} (trust: proof).

\subsubsection*{\texttt{continuous-implies-closed-preimage}}
\emph{for every s, t, f, if f is continuous from s to t, then for every a, if a is closed in t, then preimage(s, f, a) is closed in s}
{\small\[
\forall s, t, f.\; (\operatorname{is\text{-}continuous}(s, t, f) \Rightarrow \forall a.\; (\operatorname{is\text{-}closed}(t, a) \Rightarrow \operatorname{is\text{-}closed}(s, \operatorname{preimage}(s, f, a))))
\]}
\noindent\textbf{Machine-checked} (trust: proof).

\subsubsection*{\texttt{closed-preimage-implies-continuous}}
\emph{for every s, t, f, if s is a metric space and t is a metric space and f is in fun(pts(s), pts(t)), then if for every a, if a is closed in t, then preimage(s, f, a) is closed in s, then f is continuous from s to t}
{\small\[
\forall s, t, f.\; ((\operatorname{is\text{-}metric\text{-}space}(s) \wedge (\operatorname{is\text{-}metric\text{-}space}(t) \wedge (f \in (\operatorname{pts}(s) \to \operatorname{pts}(t))))) \Rightarrow (\forall a.\; (\operatorname{is\text{-}closed}(t, a) \Rightarrow \operatorname{is\text{-}closed}(s, \operatorname{preimage}(s, f, a))) \Rightarrow \operatorname{is\text{-}continuous}(s, t, f)))
\]}
\noindent\textbf{Machine-checked} (trust: proof).

\subsection{Completion}

\subsubsection*{\texttt{completion-is-metric-space}}
\emph{for every metric space m, completion(m) is a metric space}
{\small\[
\forall m.\; (\operatorname{is\text{-}metric\text{-}space}(m) \Rightarrow \operatorname{is\text{-}metric\text{-}space}(\operatorname{completion}(m)))
\]}
\noindent\textbf{Machine-checked} (trust: proof).

\subsubsection*{\texttt{completion-is-complete}}
\emph{for every metric space m, completion(m) is complete}
{\small\[
\forall m.\; (\operatorname{is\text{-}metric\text{-}space}(m) \Rightarrow \operatorname{is\text{-}complete}(\operatorname{completion}(m)))
\]}
\noindent\textbf{Machine-checked} (trust: proof).

\subsubsection*{\texttt{embed-isometry}}
\emph{for every metric space m, for every u in pts(m), v in pts(m), (dist(completion(m)))((embed(m))(u), (embed(m))(v)) equals (dist(m))(u, v)}
{\small\[
\forall m.\; (\operatorname{is\text{-}metric\text{-}space}(m) \Rightarrow \forall u, v \in \operatorname{pts}(m).\; (\operatorname{dist}(\operatorname{completion}(m))(\operatorname{embed}(m)(u), \operatorname{embed}(m)(v)) = \operatorname{dist}(m)(u, v)))
\]}
\noindent\textbf{Machine-checked} (trust: proof).

\subsection{Compactness / total boundedness}

\subsubsection*{\texttt{totally-bounded-has-cauchy-subsequence}}
\emph{for every s, if s is totally bounded, then for every f in fun(nn, pts(s)), there is some phi such that phi is a strictly increasing sequence of naturals and is-cauchy-seq(s, subseq(f, phi))}
{\small\[
\forall s.\; (\operatorname{totally\text{-}bounded}(s) \Rightarrow \forall f \in (\mathbb{N} \to \operatorname{pts}(s)).\; \exists \varphi.\; (\operatorname{strictly\text{-}mono\text{-}nn}(\varphi) \wedge \operatorname{is\text{-}cauchy\text{-}seq}(s, \operatorname{subseq}(f, \varphi))))
\]}
\noindent\textbf{Machine-checked} (trust: proof). Total boundedness supplies a finite eps-net cover at each level; block-family-combinatorial builds nested infinite blocks pinned into shrinking balls; diagonalization extracts one subsequence, and the 2r triangle estimate makes it Cauchy.

\subsubsection*{\texttt{cauchy-rapid-subsequence}}
\emph{for every s, f, rad, if is-cauchy-seq(s, f) and rad is in fun(nn, rr) and for every k in nn, rad(k) is a positive real, then there is some phi in fun(nn, nn) such that for every m in nn, n\_ in nn, if m is less than n\_, then phi(m) is less than phi(n\_) and for every k in nn, (dist(s))(f(phi(k)), f(phi(succ(k)))) is at most rad(k)}
{\small\[
\forall s, f, \mathit{rad}.\; ((\operatorname{is\text{-}cauchy\text{-}seq}(s, f) \wedge ((\mathit{rad} \in (\mathbb{N} \to \mathbb{R})) \wedge \forall k \in \mathbb{N}.\; \operatorname{pos\text{-}rr}(\operatorname{rad}(k)))) \Rightarrow \exists \varphi \in (\mathbb{N} \to \mathbb{N}).\; (\forall m, \operatorname{n\_} \in \mathbb{N}.\; ((m < \operatorname{n\_}) \Rightarrow (\varphi(m) < \varphi(\operatorname{n\_}))) \wedge \forall k \in \mathbb{N}.\; (\operatorname{dist}(s)(f(\varphi(k)), f(\varphi((k + 1)))) \leq \operatorname{rad}(k))))
\]}
\noindent\textbf{Machine-checked} (trust: proof).

\subsection{Product metric spaces}

\subsubsection*{\texttt{product-is-metric-space}}
\emph{for every ms, if ms is a sequence of metric spaces, then for every w, if w is a summable sequence of positive weights, then product-metric-w(ms, w) is a metric space}
{\small\[
\forall \mathit{ms}.\; (\operatorname{is\text{-}ms\text{-}sequence}(\mathit{ms}) \Rightarrow \forall w.\; (\operatorname{summable\text{-}weight}(w) \Rightarrow \operatorname{is\text{-}metric\text{-}space}(\operatorname{product\text{-}metric\text{-}w}(\mathit{ms}, w))))
\]}
\noindent\textbf{Machine-checked} (trust: proof).

\subsubsection*{\texttt{product-projection-continuous}}
\emph{for every ms, if ms is a sequence of metric spaces, then for every w, if w is a summable sequence of positive weights, then for every n in nn, product-proj(ms, n) is continuous from product-metric-w(ms, w) to ms(n)}
{\small\[
\forall \mathit{ms}.\; (\operatorname{is\text{-}ms\text{-}sequence}(\mathit{ms}) \Rightarrow \forall w.\; (\operatorname{summable\text{-}weight}(w) \Rightarrow \forall n \in \mathbb{N}.\; \operatorname{is\text{-}continuous}(\operatorname{product\text{-}metric\text{-}w}(\mathit{ms}, w), \operatorname{ms}(n), \operatorname{product\text{-}proj}(\mathit{ms}, n))))
\]}
\noindent\textbf{Machine-checked} (trust: proof).

\subsubsection*{\texttt{product-convergence-coordinatewise}}
\emph{for every ms, if ms is a sequence of metric spaces, then for every w, if w is a summable sequence of positive weights, then for every seq in fun(nn, product-carrier(ms)), lv in product-carrier(ms), seq converges to lv in product-metric-w(ms, w) if and only if for every n\_ in nn, vnb-lambda(k, nn, (seq(k))(n\_)) converges to lv(n\_) in ms(n\_)}
{\small\[
\forall \mathit{ms}.\; (\operatorname{is\text{-}ms\text{-}sequence}(\mathit{ms}) \Rightarrow \forall w.\; (\operatorname{summable\text{-}weight}(w) \Rightarrow \forall \mathit{seq} \in (\mathbb{N} \to \operatorname{product\text{-}carrier}(\mathit{ms})), \mathit{lv} \in \operatorname{product\text{-}carrier}(\mathit{ms}).\; (\operatorname{converges\text{-}to}(\operatorname{product\text{-}metric\text{-}w}(\mathit{ms}, w), \mathit{seq}, \mathit{lv}) \Leftrightarrow \forall \operatorname{n\_} \in \mathbb{N}.\; \operatorname{converges\text{-}to}(\operatorname{ms}(\operatorname{n\_}), (\lambda k \in \mathbb{N}.\; \operatorname{seq}(k)(\operatorname{n\_})), \operatorname{lv}(\operatorname{n\_})))))
\]}
\noindent\textbf{Machine-checked} (trust: proof).

\subsubsection*{\texttt{compact-countable-product}}
\emph{for every ms, if ms is a sequence of metric spaces, then if for every n in nn, ms(n) is compact, then product-metric(ms) is compact}
{\small\[
\forall \mathit{ms}.\; (\operatorname{is\text{-}ms\text{-}sequence}(\mathit{ms}) \Rightarrow (\forall n \in \mathbb{N}.\; \operatorname{is\text{-}compact}(\operatorname{ms}(n)) \Rightarrow \operatorname{is\text{-}compact}(\operatorname{product\text{-}metric}(\mathit{ms}))))
\]}
\noindent\textbf{Machine-checked} (trust: proof). Sequential compactness of a countable metric product, obtained by a coordinatewise diagonal subsequence.

\subsection{Normed-field metric}

\subsubsection*{\texttt{nf-metric-space-is-metric-space}}
\emph{for every normed field nf, nf-metric-space(nf) is a metric space}
{\small\[
\forall \mathit{nf}.\; (\operatorname{is\text{-}normed\text{-}field}(\mathit{nf}) \Rightarrow \operatorname{is\text{-}metric\text{-}space}(\operatorname{nf\text{-}metric\text{-}space}(\mathit{nf})))
\]}
\noindent\textbf{Machine-checked} (trust: proof).

\end{document}
